% Generated mechanically from the frozen formal-defos.tex target.
% Do not edit this generated file; edit the bound locale source instead.


% OK, start here.
%

\setcounter{chapter}{89}
\chapter{形式形变理论}



\phantomsection
\label{formal-defos-section-phantom}




\section{引言}
\label{formal-defos-section-introduction}

\noindent
本章建立一种可在 Stacks 项目后续章节中应用的形式形变理论，
其论述紧随 Rim \cite[Exposee VI]{SGA7-I} 和 Schlessinger \cite{Sch}。
我们强烈建议初次接触这一主题的读者先阅读 Schlessinger 的论文；
该文的一般性已足以应付大多数应用，而且采用这种形式形变理论的论文
事实上大多都会使用 Schlessinger 的结果。

\medskip\noindent
设 $\Lambda$ 为完备诺特局部环，剩余域为 $k$；以
$\mathcal{C}_\Lambda$ 表示剩余域为 $k$ 的阿廷局部
$\Lambda$-代数所成的范畴。给定函子
$F : \mathcal{C}_\Lambda \to \textit{Sets}$，并假设 $F(k)$
是单元素集合。Schlessinger 的论文引入了条件 (H1)--(H4)，使得：
\begin{enumerate}
\item $F$ 有一个“包络”当且仅当 (H1)--(H3) 成立。
\item $F$ 是 pro-可表的当且仅当 (H1)--(H4) 成立。
\end{enumerate}
本章的目的是完全按照 Rim 的论文，从以下两方面推广这些结果：
\begin{enumerate}
\item[(A)] 以 $\mathcal{C}_\Lambda$ 上的群胚余纤维范畴
$\mathcal{F}$ 取代函子 $F$；见第 \ref{formal-defos-section-CLambda} 节。
\item[(B)] 我们允许 $\Lambda$ 为诺特环，并令 $\Lambda \to k$
为到一个域的有限环映射。范畴 $\mathcal{C}_\Lambda$ 的对象是阿廷局部
$\Lambda$-代数 $A$，并配备一个给定的等同 $A/\mathfrak m_A = k$。
\end{enumerate}
条件 $F(k)$ 为单元素集合的对应物，是 $\mathcal{F}(k)$ 为平凡群胚。
若 $\mathcal{F}$ 满足此条件，我们称它为一个{\it 预形变范畴}；
但一般而言我们并不作此假设。Rim 的论文
\cite[Exposee VI]{SGA7-I} 是本文各项结果的原始来源。
我们还要提及很有用的论文 \cite{Talpo-Vistoli}；该文也用群胚讨论形变理论，
但一般性不及本文。

\medskip\noindent
范畴 $\mathcal{C}_\Lambda$ 的“完备化”
$\widehat{\mathcal{C}}_\Lambda$ 起着重要作用。
$\widehat{\mathcal{C}}_\Lambda$ 的对象是诺特完备局部
$\Lambda$-代数 $R$，其剩余域配备了与 $k$ 的等同；见
第 \ref{formal-defos-section-category-completion-CLambda} 节。
一方面，$\mathcal{C}_\Lambda \subset \widehat{\mathcal{C}}_\Lambda$
是严格满子范畴；另一方面，$\widehat{\mathcal{C}}_\Lambda$
是 $\mathcal{C}_\Lambda$ 的 pro-对象范畴的满子范畴。若函子
$\mathcal{C}_\Lambda \to \textit{Sets}$ 同构于可表函子
$\underline{R} = \Mor_{\widehat{\mathcal{C}}_\Lambda}(R, -)$
在 $\mathcal{C}_\Lambda$ 上的限制，其中
$R \in \Ob(\widehat{\mathcal{C}}_\Lambda)$，则称该函子为
{\it pro-可表的}。

\medskip\noindent
{\it 群胚余纤维范畴}与群胚纤维范畴对偶；我们将在
第 \ref{formal-defos-section-preliminary} 节引入它们。
$\mathcal{C}_\Lambda$ 上群胚余纤维范畴之间的一个{\it 光滑}态射，
是指满足对象的无穷小提升判据的态射；见
第 \ref{formal-defos-section-smooth-morphisms} 节。
这与形式光滑环映射的定义类似；见《代数》定义
\ref{algebra-definition-formally-smooth}，并且恰好与
《可表示性判据》第 \ref{criteria-section-formally-smooth} 节中的概念对偶。
这是一个重要概念，因为我们最终要证明某些群胚余纤维范畴
具有光滑的 pro-可表表示，这与《代数叠》第
\ref{algebraic-section-stack-to-presentation} 和
\ref{algebraic-section-smooth-groupoid-gives-algebraic-stack} 节中
对代数叠的刻画十分相似。
$\mathcal{C}_\Lambda$ 上群胚余纤维范畴 $\mathcal{F}$ 的一个{\it 泛形式对象}，是其完备化中的对象
$\xi \in \widehat{\mathcal{F}}(R)$，使相伴态射
$\underline{\xi} : \underline{R}|_{\mathcal{C}_\Lambda} \to \mathcal{F}$
为光滑态射。

\medskip\noindent
在第 \ref{formal-defos-section-schlessinger-conditions} 节中，
我们在 $\mathcal{F}$ 上定义条件 (S1) 和 (S2)，以推广 Schlessinger 的
(H1) 和 (H2)。Schlessinger 的 (H3) 所对应的条件——即
$\mathcal{F}$ 有有限维切空间——不另命名。
发展这套理论的一个关键步骤，是为满足 (S1)、(S2) 和 (H3) 的
预形变范畴证明泛形式对象的存在性；见引理
\ref{formal-defos-lemma-versal-object-existence}。
Schlessinger 对函子 $F : \mathcal{C}_\Lambda \to \textit{Sets}$
所定义的{\it 包络}，在我们的术语中，是一个泛形式对象
$\xi \in \widehat{F}(R)$，且其诱导的切空间映射
$d\underline{\xi} : T\underline{R}|_{\mathcal{C}_\Lambda} \to TF$
为同构。
文献中常把包络称为“极小泛”对象。我们不采用这一叫法，理由如下。
一个函子可能有泛形式对象，却没有包络。此外，我们将在
第 \ref{formal-defos-section-minimal-versal} 节证明：若预形变范畴有泛形式对象，
则它总有一个{\it 极小}泛形式对象（定义见
定义 \ref{formal-defos-definition-minimal-versal}），并且该对象在同构意义下唯一；
见引理 \ref{formal-defos-lemma-minimal-versal}。
然而，极小泛形式对象未必在切空间上诱导同构！（见例
\ref{formal-defos-example-do-not-get-S2} 和
\ref{formal-defos-example-smooth-continued}。）

\medskip\noindent
在牢记上述差异的前提下，
定理 \ref{formal-defos-theorem-miniversal-object-existence}
正是上面 (1) 的直接推广：当 $\mathcal{F}$ 是函子时，它恢复
Schlessinger 的结果；在条件 (S1) 和 (S2) 成立时，它又通过切空间映射
$d\underline{\xi} : T\underline{R}|_{\mathcal{C}_\Lambda} \to TF$
刻画极小泛形式对象。

\medskip\noindent
在第 \ref{formal-defos-section-RS-condition} 节中，
我们在 $\mathcal{F}$ 上定义 Rim 的条件 (RS)，以推广 Schlessinger 的
(H4)。所谓{\it 形变范畴}，是指满足 (RS) 的预形变范畴。
pro-可表函子的对应物，是 $\mathcal{C}_\Lambda$ 上那些
{\it 可由 $\mathcal{C}_\Lambda$ 上函子中的光滑 pro-可表群胚表示}的群胚余纤维范畴；见定义
\ref{formal-defos-definition-groupoid-in-functors}、
\ref{formal-defos-definition-prorepresentable-groupoid-in-functors} 和
\ref{formal-defos-definition-smooth-groupoid-in-functors}。
这种表示概念把 $\mathcal{F}$ 各纤维的群胚结构纳入考虑。
在定理 \ref{formal-defos-theorem-presentation-deformation-groupoid} 中，
我们证明：$\mathcal{F}$ 可由函子中的光滑 pro-可表群胚表示，
当且仅当 $\mathcal{F}$ 的切空间和无穷小自同构空间均为有限维。
这是上面 (2) 的推广；当 $\mathcal{F}$ 是函子时，它归结为
Schlessinger 的结果。最后还有第 \ref{formal-defos-section-minimality} 节，
其中讨论如何利用极小泛形式对象构造函子中的光滑 pro-可表群胚的
极小表示；这种表示在同构意义下唯一。

\medskip\noindent
我们还得到对 Schlessinger 条件的如下概念性解释。
若预形变范畴 $\mathcal{F}$ 满足 (RS)，则相伴的同构类函子
$\overline{\mathcal{F}}: \mathcal{C}_\Lambda \to \textit{Sets}$
满足 (H1) 和 (H2)（见引理
\ref{formal-defos-lemma-RS-implies-S1-S2} 和
\ref{formal-defos-lemma-S1-S2-associated-functor}）。
反过来，若函子 $F : \mathcal{C}_\Lambda \to \textit{Sets}$
自然地产生于一个群胚余纤维范畴 $\mathcal{F}$ 的同构类函子，
那么实践中，用来证明 $F$ 满足 (H1) 和 (H2) 的论证，似乎也总能证明
$\mathcal{F}$ 满足 (RS)。相关例子见《形变问题》第
\ref{examples-defos-section-introduction} 节。
此外，若 $\mathcal{F}$ 满足 (RS)，则 $\overline{\mathcal{F}}$ 的条件
(H4) 可以用 $\mathcal{F}$ 中对象的自同构能否延拓来作简单解释
（引理 \ref{formal-defos-lemma-RS-associated-functor}）。
这些观察表明，应把 (RS) 看作形变理论中基本的黏合条件。




\section{记号与约定}
\label{formal-defos-section-notations-conventions}

\noindent
环均指带 $1$ 的交换环。局部环 $A$ 的极大理想记为
$\mathfrak{m}_A$。正整数集合记为
$\mathbf{N} = \{1, 2, 3, \ldots\}$。若 $U$ 是范畴
$\mathcal{C}$ 的一个对象，则以 $\underline{U}$ 表示函子
$\Mor_\mathcal{C}(U, -): \mathcal{C} \to \textit{Sets}$；见备注
\ref{formal-defos-remarks-cofibered-groupoids} 的
(\ref{formal-defos-item-definition-yoneda})。注意：这可能与其他章节的记号冲突；
在那里我们有时用 $\underline{U}$ 表示
$h_U(-) = \Mor_\mathcal{C}(-, U)$。

\medskip\noindent
本章始终设 $\Lambda$ 为诺特环，并设
$\Lambda \to k$ 是从 $\Lambda$ 到一个域的有限环映射。
该映射的核记为 $\mathfrak m_\Lambda$，像记为 $k' \subset k$。
可以证明 $\mathfrak m_\Lambda$ 是极大理想，
$k' = \Lambda/\mathfrak m_\Lambda$ 是域，且扩张 $k/k'$ 有限。
见公式 (\ref{formal-defos-equation-k-prime}) 附近的讨论。


\section{基范畴}
\label{formal-defos-section-CLambda}

\noindent
动机。形式形变理论的一个重要应用，是研究由代数空间表示的判据。
设给定局部诺特基底 $S$ 和函子
$F : (\Sch/S)_{fppf}^{opp} \to \textit{Sets}$。
设 $k$ 是 $S$ 上的有限型域，也就是说，给定有限型态射
$\Spec(k) \to S$。
Artin 的判据之一要求：对任意元素 $x \in F(\Spec(k))$，
与三元组 $(S, k, x)$ 相伴的预形变函子应为 pro-可表的。
由《态射》引理 \ref{morphisms-lemma-point-finite-type}，
$k$ 在 $S$ 上为有限型这一条件意味着存在仿射开集
$\Spec(\Lambda) \subset S$，使 $k$ 是有限 $\Lambda$-代数。
这说明了为何本章始终采用如下基范畴。

\begin{definition}
\label{formal-defos-definition-CLambda}
设 $\Lambda$ 为诺特环，$\Lambda \to k$ 为有限环映射，其中 $k$ 是域。
定义{\it $\mathcal{C}_\Lambda$} 为如下范畴：
\begin{enumerate}
\item 对象为二元组 $(A, \varphi)$，其中 $A$ 是阿廷局部
$\Lambda$-代数，而 $\varphi : A/\mathfrak m_A \to k$ 是
$\Lambda$-代数同构；
\item 态射 $f : (B, \psi) \to (A, \varphi)$ 为局部
$\Lambda$-代数同态，且满足
$\varphi \circ (f \bmod \mathfrak m) = \psi$。
\end{enumerate}
若 $\Lambda$ 是诺特完备局部环，而 $k$ 是其剩余域，
则称我们处于{\it 经典情形}。
\end{definition}

\noindent
注意，若 $\Lambda \to k$ 满射，且 $A$ 是阿廷局部
$\Lambda$-代数，则等同 $\varphi$ 一旦存在便是唯一的。
而且在这种情形下，任意 $\Lambda$-代数映射 $A \to B$
都与这些等同相容。因此，此时 $\mathcal{C}_\Lambda$
就是剩余域“为”$k$ 的阿廷局部 $\Lambda$-代数所成的范畴。
在一般情形下，我们也滥用记号，直接以 $A$ 表示
$\mathcal{C}_\Lambda$ 的对象。此外，我们常写
$A/\mathfrak m = k$，即假装 $\mathcal{C}_\Lambda$ 中所有环的
剩余域都是 $k$（由于 $\mathcal{C}_\Lambda$ 中所有环映射都与给定等同相容，
这不会造成任何问题）。
本章余下部分固定基环 $\Lambda$ 和域 $k$。
$\mathcal{C}_\Lambda$ 将作为下文所考虑余纤维范畴的基范畴。

\begin{definition}
\label{formal-defos-definition-small-extension}
设 $f: B \to A$ 是 $\mathcal{C}_\Lambda$ 中的环映射。
若 $f$ 为满射，并且 $\Ker(f)$ 是被 $\mathfrak{m}_B$ 零化的
非零主理想，则称 $f$ 为一个{\it 小扩张}。
\end{definition}

\noindent
由下面的引理，我们常可把 $\mathcal{C}_\Lambda$ 中涉及满环映射的论证
归约到小扩张的情形。

\begin{lemma}
\label{formal-defos-lemma-factor-small-extension}
设 $f: B \to A$ 是 $\mathcal{C}_\Lambda$ 中的满环映射。
则 $f$ 可分解为一列小扩张的复合。
\end{lemma}

\begin{proof}
设 $I$ 为 $f$ 的核。由于 $B$ 是阿廷环，极大理想
$\mathfrak{m}_B$ 幂零；设 $\mathfrak{m}_B^n = 0$。于是 $f$ 有分解
$$
B = B/I\mathfrak{m}_B^{n-1} \to B/I\mathfrak{m}_B^{n-2} \to
\ldots \to B/I \cong A
$$
其中每个映射均为满射，且其核被极大理想零化。
因此只需证明 $f$ 本身具有这一性质的情形，也就是说，
$I$ 被 $\mathfrak{m}_B$ 零化的情形。此时 $I$ 是有限维
$k$-向量空间；见《代数》引理
\ref{algebra-lemma-artinian-finite-length}。
取 $I$ 作为 $k$-向量空间的一组基 $x_1, \ldots, x_n$，便得到分解
$$
B \to B/(x_1) \to \ldots \to  B/(x_1, \ldots, x_n) \cong  A
$$
其中每个映射都是小扩张。
\end{proof}

\noindent
下一个引理说明：剩余域为 $k$ 的局部 $\Lambda$-代数上的模，
其长度可由它作为 $\Lambda$-模的长度计算。为解释陈述中的记号，
设 $k' \subset k$ 为我们固定的有限环映射 $\Lambda \to k$ 的像。
注意 $k' \subset k$ 是有限环扩张。因此 $k'$ 是域，而 $k/k'$
是有限域扩张；见《代数》引理
\ref{algebra-lemma-integral-under-field}。
又因 $\Lambda \to k'$ 为满射，其核是极大理想
$\mathfrak m_\Lambda$。因此
\begin{equation}
\label{formal-defos-equation-k-prime}
[k : k'] = [k : \Lambda/\mathfrak m_\Lambda] < \infty
\end{equation}
且在经典情形下 $k = k'$。本章始终固定记号
$k' = \Lambda/\mathfrak m_\Lambda$。

\begin{lemma}
\label{formal-defos-lemma-length}
设 $A$ 是剩余域为 $k$ 的局部 $\Lambda$-代数，$M$ 是 $A$-模。则
$[k : k'] \text{length}_A(M) = \text{length}_\Lambda(M)$。
在经典情形下，
$\text{length}_A(M) = \text{length}_\Lambda(M)$。
\end{lemma}

\begin{proof}
若 $M$ 是单 $A$-模，则 $M \cong k$（作为 $A$-模）；见《代数》引理
\ref{algebra-lemma-characterize-length-1}。
此时 $\text{length}_A(M) = 1$，且
$\text{length}_\Lambda(M) = [k' : k]$；% P10-E0420
见《代数》引理 \ref{algebra-lemma-dimension-is-length}。
若 $\text{length}_A(M)$ 有限，则取 $M$ 的一个 $A$-子模滤过，
使其逐次商均为单模，再由长度的可加性得到结论；见《代数》引理
\ref{algebra-lemma-length-additive}。
若 $\text{length}_A(M)$ 无限，则结论由显然不等式
$\text{length}_A(M) \leq \text{length}_\Lambda(M)$ 得出。
\end{proof}

\begin{lemma}
\label{formal-defos-lemma-surjective}
设 $A \to B$ 是 $\mathcal{C}_\Lambda$ 中的环映射。下列条件等价：
\begin{enumerate}
\item $f$ 为满射；% P10-E0421
\item $\mathfrak m_A/\mathfrak m_A^2 \to \mathfrak m_B/\mathfrak m_B^2$
为满射；
\item $\mathfrak m_A/(\mathfrak m_\Lambda A + \mathfrak m_A^2)
\to \mathfrak m_B/(\mathfrak m_\Lambda B + \mathfrak m_B^2)$ 为满射。
\end{enumerate}
\end{lemma}

\begin{proof}
对 $\mathcal{C}_\Lambda$ 中任意环映射 $f : A \to B$，都有
$f(\mathfrak m_A) \subset \mathfrak m_B$；例如，这是因为
$\mathfrak m_A$ 和 $\mathfrak m_B$ 分别是 $A$ 和 $B$ 中幂零元的集合。
设 $f$ 为满射。取 $y \in \mathfrak m_B$，再取 $x \in A$ 使
$f(x) = y$。由于 $f$ 诱导同构
$A/\mathfrak m_A \to B/\mathfrak m_B$，可知 $x \in \mathfrak m_A$。
故诱导映射
$\mathfrak m_A/\mathfrak m_A^2 \to \mathfrak m_B/\mathfrak m_B^2$
为满射。这证明了 (1) 蕴含 (2)。

\medskip\noindent
显然 (2) 蕴含 (3)。映射 $A \to B$ 给出典范交换图
$$
\xymatrix{
\mathfrak m_\Lambda/\mathfrak m_\Lambda^2 \otimes_{k'} k \ar[r] \ar[d] &
\mathfrak m_A/\mathfrak m_A^2 \ar[r] \ar[d] &
\mathfrak m_A/(\mathfrak m_\Lambda A + \mathfrak m_A^2) \ar[r] \ar[d] & 0 \\
\mathfrak m_\Lambda/\mathfrak m_\Lambda^2 \otimes_{k'} k \ar[r] &
\mathfrak m_B/\mathfrak m_B^2 \ar[r] &
\mathfrak m_B/(\mathfrak m_\Lambda B + \mathfrak m_B^2) \ar[r] & 0
}
$$
其两行正合。因此若 (3) 成立，则 (2) 也成立。

\medskip\noindent
设 (2) 成立。由 Nakayama 引理（《代数》引理
\ref{algebra-lemma-NAK}），要证明 $A \to B$ 为满射，
只需证明 $A/\mathfrak m_A \to B/\mathfrak m_AB$ 为满射。
（注意 $\mathfrak m_A$ 是幂零理想。）由于
$k = A/\mathfrak m_A = B/\mathfrak m_B$，只需证明
$\mathfrak m_AB \to \mathfrak m_B$ 为满射。再次应用 Nakayama 引理，
只需证明
$\mathfrak m_AB/\mathfrak m_A\mathfrak m_B \to
\mathfrak m_B/\mathfrak m_B^2$
为满射，而这正是我们的假设。
\end{proof}

\noindent
若 $A \to B$ 是 $\mathcal{C}_\Lambda$ 中的环映射，则映射
$\mathfrak m_A/(\mathfrak m_\Lambda A + \mathfrak m_A^2)
\to \mathfrak m_B/(\mathfrak m_\Lambda B + \mathfrak m_B^2)$
就是相对余切空间上的映射。正式定义如下。

\begin{definition}
\label{formal-defos-definition-tangent-space-ring}
设 $R \to S$ 是局部环之间的局部同态。所谓 $R$ 在 $S$ 上的
{\it 相对余切空间}% P10-E0422
\footnote{注意：我们稍后将看到，在本章的一般设定下，
$\Lambda$ 上对象 $A \in \mathcal{C}_\Lambda$ 的切空间，
不应简单地定义为相对余切空间的 $k$-线性对偶。
事实上，相对余切空间的正确定义是
$\Omega_{S/R} \otimes_S S/\mathfrak m_S$。}
是如下 $S/\mathfrak m_S$-向量空间：
$\mathfrak m_S/(\mathfrak m_R S + \mathfrak m_S^2)$。
\end{definition}

\noindent
若 $f_1: A_1 \to A$ 和 $f_2: A_2 \to A$ 是两个环映射，
则纤维积 $A_1 \times_A A_2$ 是 $A_1 \times A_2$ 的子环，
由两个投影到 $A$ 的像相等的元素组成。本章将反复考虑
$f_1$ 和 $f_2$ 属于 $\mathcal{C}_\Lambda$ 时涉及这种纤维积的条件。
但该纤维积未必是 $\mathcal{C}_\Lambda$ 的对象。

\begin{example}
\label{formal-defos-example-fibre-product}
设 $p$ 为素数，$n \in \mathbf{N}$。令
$\Lambda = \mathbf{F}_p(t_1, t_2, \ldots, t_n)$，并令
$k = \mathbf{F}_p(x_1, \ldots, x_n)$；映射 $\Lambda \to k$ 由
$t_i \mapsto x_i^p$ 给出。令 $A = k[\epsilon] = k[x]/(x^2)$。
则 $A$ 是 $\mathcal{C}_\Lambda$ 的对象。设
$D : k \to k$ 是 $k$ 在 $\Lambda$ 上的一个导子，例如
$D = \partial/\partial x_i$。则映射
$$
f_D : k \longrightarrow k[\epsilon], \quad
a \mapsto a + D(a)\epsilon
$$
是 $\mathcal{C}_\Lambda$ 中的态射。令 $A_1 = A_2 = k$，并令
$f_1 = f_{\partial/\partial x_1}$、$f_2(a) = a$。于是
$A_1 \times_A A_2 = \{a \in k \mid \partial/\partial x_1(a) = 0\}$，
它不满射到 $k$。因此该纤维积不是 $\mathcal{C}_\Lambda$ 的对象。
\end{example}

\noindent
可以证明，这一问题只会在剩余域扩张
$k/k'$（见 (\ref{formal-defos-equation-k-prime})）不可分，并且
$f_1$、$f_2$ 都不是满射时发生。

\begin{lemma}
\label{formal-defos-lemma-fiber-product-CLambda}
设 $f_1 : A_1 \to A$ 和 $f_2 : A_2 \to A$ 是
$\mathcal{C}_\Lambda$ 中的环映射。则：
\begin{enumerate}
\item 若 $f_1$ 或 $f_2$ 为满射，则
$A_1 \times_A A_2$ 属于 $\mathcal{C}_\Lambda$；
\item 若 $f_2$ 为小扩张，则
$A_1 \times_A A_2 \to A_1$ 也是小扩张；
\item 若域扩张 $k/k'$ 可分，则
$A_1 \times_A A_2$ 属于 $\mathcal{C}_\Lambda$。
\end{enumerate}
\end{lemma}

\begin{proof}
映射 $\Lambda \to A_1$ 和 $\Lambda \to A_2$ 诱导映射
$\Lambda \to A_1 \times_A A_2$，从而使环
$A_1 \times_A A_2$ 成为 $\Lambda$-代数。它是局部环，唯一极大理想为
$$
\mathfrak m_{A_1} \times_{\mathfrak m_A} \mathfrak m_{A_2} =
\Ker(A_1 \times_A A_2 \longrightarrow k)
$$
一个环是阿廷环，当且仅当它作为自身的模具有有限长度；见《代数》引理
\ref{algebra-lemma-artinian-finite-length}。
由于 $A_1$ 和 $A_2$ 是阿廷环，引理 \ref{formal-defos-lemma-length} 表明
$\text{length}_\Lambda(A_1)$ 和 $\text{length}_\Lambda(A_2)$，
从而也表明 $\text{length}_\Lambda(A_1 \times A_2)$，都是有限的。
又因 $A_1 \times_A A_2 \subset A_1 \times A_2$ 是
$\Lambda$-子模，便有
$\text{length}_{A_1 \times_A A_2}(A_1 \times_A A_2) \leq
\text{length}_\Lambda(A_1 \times_A A_2)$ 有限。因此
$A_1 \times_A A_2$ 是阿廷环。故阻碍
$A_1 \times_A A_2$ 成为 $\mathcal{C}_\Lambda$ 对象的唯一可能，
是它的剩余域经上述映射
$A_1 \times_A A_2 \to A \to A/\mathfrak m_A = k$
映到 $k$ 的一个真子域。

\medskip\noindent
(1) 的证明。若 $f_2$ 为满射，则投影
$A_1 \times_A A_2 \to A_1$ 为满射。因此复合
$A_1 \times_A A_2 \to A_1 \to A_1/\mathfrak m_{A_1} = k$
为满射，从而 $A_1 \times_A A_2$ 是
$\mathcal{C}_\Lambda$ 的对象。

\medskip\noindent
(2) 的证明。若 $f_2$ 为小扩张，则 $A_2 \to A$ 与
$A_1 \times_A A_2 \to A_1$ 都是满射，且二者具有相同的核。
因此 $A_1 \times_A A_2 \to A_1$ 的核是 $1$-维
$k$-向量空间，故 $A_1 \times_A A_2 \to A_1$ 为小扩张。

\medskip\noindent
(3) 的证明。选取 $\overline{x} \in k$，使
$k = k'(\overline{x})$（见《域》引理
\ref{fields-lemma-primitive-element}）。
设 $P'(T) \in k'[T]$ 是 $\overline{x}$ 在 $k'$ 上的极小多项式。
由于 $k/k'$ 可分，
$\text{d}P/\text{d}T(\overline{x}) \not = 0$。
在满射 $\Lambda[T] \to k'[T]$ 下，选取映到 $P'$ 的首一多项式
$P \in \Lambda[T]$。由于 $A, A_1, A_2$ 都是 Hensel 环；见《代数》引理
\ref{algebra-lemma-local-dimension-zero-henselian}，
可找到 $x, x_1, x_2 \in A, A_1, A_2$，使
$P(x) = 0, P(x_1) = 0, P(x_2) = 0$，且 $x,x_1,x_2$
在 $k$ 中的像均为 $\overline{x}$。
由唯一性，$x_1,x_2$ 都映到 $x \in A$；见《代数》引理
\ref{algebra-lemma-uniqueness}。所以
$(x_1, x_2) \in A_1 \times_A A_2$。
于是 $A_1 \times_A A_2$ 的剩余域包含 $k$ 在 $k'$ 上的一个生成元，
结论得证。
\end{proof}


\noindent
下面定义 $\mathcal{C}_\Lambda$ 中的本质满射。
引理 \ref{formal-defos-lemma-essential-surjection} 给出了
$\mathcal{C}_\Lambda$ 中一个满射为本质满射的充要条件。

\begin{definition}
\label{formal-defos-definition-essential-surjection}
设 $f: B \to A$ 是 $\mathcal{C}_\Lambda$ 中的环映射。
若其具有下列性质，则称 $f$ 为{\it 本质满射}：
\begin{enumerate}
\item $f$ 为满射；
\item 若 $g: C \to B$ 是 $\mathcal{C}_\Lambda$ 中的环映射，
且 $f \circ g$ 为满射，则 $g$ 为满射。
\end{enumerate}
\end{definition}

\noindent
利用引理 \ref{formal-defos-lemma-surjective}，可以如下刻画
$\mathcal{C}_\Lambda$ 中的本质满射。

\begin{lemma}
\label{formal-defos-lemma-essential-surjection-mod-squares}
设 $f: B \to A$ 是 $\mathcal{C}_\Lambda$ 中的环映射。
下列条件等价：
\begin{enumerate}
\item $f$ 是本质满射；
\item 映射 $B/\mathfrak m_B^2 \to A/\mathfrak m_A^2$ 是本质满射；
\item 映射
$B/(\mathfrak m_\Lambda B + \mathfrak m_B^2) \to
A/(\mathfrak m_\Lambda A + \mathfrak m_A^2)$ 是本质满射。
\end{enumerate}
\end{lemma}

\begin{proof}
假设 (3) 成立。设 $C \to B$ 是 $\mathcal{C}_\Lambda$ 中的环映射，
且 $C \to A$ 为满射。则
$C \to A/(\mathfrak m_\Lambda A + \mathfrak m_A^2)$ 也为满射。
由假设，$C \to B/(\mathfrak m_\Lambda B + \mathfrak m_B^2)$
为满射。两次应用引理 \ref{formal-defos-lemma-surjective}，可知
$C \to B$ 为满射。

\medskip\noindent
假设 (1) 成立。设
$C \to B/(\mathfrak m_\Lambda B + \mathfrak m_B^2)$
是 $\mathcal{C}_\Lambda$ 中的态射，且
$C \to A/(\mathfrak m_\Lambda A + \mathfrak m_A^2)$ 为满射。令
$C' = C \times_{B/(\mathfrak m_\Lambda B + \mathfrak m_B^2)} B$。
由引理 \ref{formal-defos-lemma-fiber-product-CLambda}，它是
$\mathcal{C}_\Lambda$ 的对象。
注意 $C' \to A/(\mathfrak m_\Lambda A + \mathfrak m_A^2)$
仍为满射，故由引理 \ref{formal-defos-lemma-surjective}，$C' \to A$ 为满射。
于是由假设，$C' \to B$ 为满射。因而
$C' \to B/(\mathfrak m_\Lambda B + \mathfrak m_B^2)$ 为满射；
由 $C'$ 的构造，这又蕴含
$C \to B/(\mathfrak m_\Lambda B + \mathfrak m_B^2)$ 为满射。

\medskip\noindent
第一段证明了 (3) $\Rightarrow$ (1)，第二段证明了
(1) $\Rightarrow$ (3)。(2) 与 (3) 的等价，是 (1) 与 (3)
等价的一个特殊情形，故证明完成。
\end{proof}

\noindent
为了进一步分析 $\mathcal{C}_\Lambda$ 中的本质满射，
引入一些记号。设 $A$ 是 $\mathcal{C}_\Lambda$ 的对象；
更一般地，也可取配备了 $\Lambda$-代数满射 $A \to k$
的任意 $\Lambda$-代数。有典范正合序列
\begin{equation}
\label{formal-defos-equation-sequence}
\mathfrak m_A/\mathfrak m_A^2 \xrightarrow{\text{d}_A}
\Omega_{A/\Lambda} \otimes_A k \to
\Omega_{k/\Lambda} \to 0
\end{equation}
见《代数》引理 \ref{algebra-lemma-differential-seq}。
注意 $\Omega_{k/\Lambda} = \Omega_{k/k'}$，其中 $k'$ 如
(\ref{formal-defos-equation-k-prime}) 所定义。以 $H_1(L_{k/\Lambda})$
表示 $k$ 在 $\Lambda$ 上的朴素余切复形的第一同调模；见《代数》定义
\ref{algebra-definition-naive-cotangent-complex}。
于是可以把 (\ref{formal-defos-equation-sequence}) 延长为正合序列
\begin{equation}
\label{formal-defos-equation-sequence-extended}
H_1(L_{k/\Lambda}) \to
\mathfrak m_A/\mathfrak m_A^2 \xrightarrow{\text{d}_A}
\Omega_{A/\Lambda} \otimes_A k \to
\Omega_{k/\Lambda} \to 0,
\end{equation}
见《代数》引理 \ref{algebra-lemma-exact-sequence-NL}。
若 $B \to A$ 是 $\mathcal{C}_\Lambda$ 中的环映射；
更一般地，若它是两个均配备到 $k$ 的 $\Lambda$-代数满射的
$\Lambda$-代数之间的映射，则得到交换图
\begin{equation}
\label{formal-defos-equation-sequence-functorial}
\vcenter{
\xymatrix{
H_1(L_{k/\Lambda}) \ar[r] \ar@{=}[d] &
\mathfrak m_B/\mathfrak m_B^2 \ar[r]_{\text{d}_B} \ar[d] &
\Omega_{B/\Lambda} \otimes_B k \ar[r] \ar[d] &
\Omega_{k/\Lambda} \ar[r] \ar@{=}[d] & 0 \\
H_1(L_{k/\Lambda}) \ar[r] &
\mathfrak m_A/\mathfrak m_A^2 \ar[r]^{\text{d}_A} &
\Omega_{A/\Lambda} \otimes_A k \ar[r] &
\Omega_{k/\Lambda} \ar[r] & 0
}
}
\end{equation}
其两行均正合。

\begin{lemma}
\label{formal-defos-lemma-H1-separable-case}
有典范映射
$$
\mathfrak m_\Lambda/\mathfrak m_\Lambda^2 \longrightarrow H_1(L_{k/\Lambda}).
$$
若 $k' \subset k$ 可分（例如 $k$ 的特征为零），则该映射诱导同构
$\mathfrak m_\Lambda/\mathfrak m_\Lambda^2 \otimes_{k'} k = H_1(L_{k/\Lambda})$。
若 $k = k'$（例如经典情形），则
$\mathfrak m_\Lambda/\mathfrak m_\Lambda^2 = H_1(L_{k/\Lambda})$。
复合
$$
\mathfrak m_\Lambda/\mathfrak m_\Lambda^2 \longrightarrow
H_1(L_{k/\Lambda}) \longrightarrow \mathfrak m_A/\mathfrak m_A^2
$$
来自典范映射 $\mathfrak m_\Lambda \to \mathfrak m_A$。
\end{lemma}

\begin{proof}
由于 $\Lambda \to k'$ 为满射且核为 $\mathfrak m_\Lambda$，
故 $H_1(L_{k'/\Lambda}) = \mathfrak m_\Lambda/\mathfrak m_\Lambda^2$。
上述映射来自朴素余切复形的函子性。若 $k' \subset k$ 可分，
则 $k' \to k$ 是平展环映射；见《代数》引理
\ref{algebra-lemma-etale-over-field}。
因此它的朴素余切复形具有平凡的同调群；见《代数》定义
\ref{algebra-definition-etale}。
把《代数》引理 \ref{algebra-lemma-exact-sequence-NL}
应用于环映射 $\Lambda \to k' \to k$，可得
$\mathfrak m_\Lambda/\mathfrak m_\Lambda^2 \otimes_{k'} k = H_1(L_{k/\Lambda})$。
最后一个断言的证明从略。
\end{proof}


\begin{lemma}
\label{formal-defos-lemma-essential-surjection}
设 $f: B \to A$ 是 $\mathcal{C}_\Lambda$ 中的环映射，
记号沿用 (\ref{formal-defos-equation-sequence-functorial})。
\begin{enumerate}
\item 引理 \ref{formal-defos-lemma-surjective} 中刻画 $f$ 为满射的等价条件，
还等价于下列条件：
\begin{enumerate}
\item $\Im(\text{d}_B) \to \Im(\text{d}_A)$ 为满射；
\item 映射 $\Omega_{B/\Lambda} \otimes_B k \to
\Omega_{A/\Lambda} \otimes_A k$ 为满射。
\end{enumerate}
\item 下列条件等价：
\begin{enumerate}
\item $f$ 是本质满射（见引理
\ref{formal-defos-lemma-essential-surjection-mod-squares}）；
\item 映射 $\Im(\text{d}_B) \to \Im(\text{d}_A)$ 为同构；
\item 映射 $\Omega_{B/\Lambda} \otimes_B k \to
\Omega_{A/\Lambda} \otimes_A k$ 为同构。
\end{enumerate}
\item 若 $k/k'$ 可分，则 $f$ 为本质满射，当且仅当映射
$\mathfrak m_B/(\mathfrak m_\Lambda B + \mathfrak m_B^2) \to
\mathfrak m_A/(\mathfrak m_\Lambda A + \mathfrak m_A^2)$
为同构。
\item 若 $f$ 为小扩张，则 $f$ 非本质，当且仅当 $f$ 在
$\mathcal{C}_\Lambda$ 中有截面 $s: A \to B$，且
$f \circ s = \text{id}_A$。
\end{enumerate}
\end{lemma}

\begin{proof}
(1) 的证明。由 (\ref{formal-defos-equation-sequence-functorial})，
(1)(a) 与 (1)(b) 等价。此外，若 $A \to B$ 为满射，% P10-E0423
则 (1)(a) 与 (1)(b) 成立。设 (1)(a) 成立。
由于 $\text{d}_A$ 的核是 $H_1(L_{k/\Lambda})$ 的像，
而后者也映到 $\mathfrak m_B/\mathfrak m_B^2$，可知
$\mathfrak m_B/\mathfrak m_B^2 \to \mathfrak m_A/\mathfrak m_A^2$
为满射。因此由引理 \ref{formal-defos-lemma-surjective}，$B \to A$ 为满射。
这就完成了 (1) 的证明。

\medskip\noindent
(2) 的证明。(2)(b) 与 (2)(c) 的等价性直接来自
(\ref{formal-defos-equation-sequence-functorial})。

\medskip\noindent
设 (2)(b) 成立。令 $g : C \to B$ 为
$\mathcal{C}_\Lambda$ 中的环映射，且 $f \circ g$ 为满射。
由引理 \ref{formal-defos-lemma-surjective}，
$\mathfrak m_C/\mathfrak m_C^2 \to \mathfrak m_A/\mathfrak m_A^2$
为满射。因此 $\Im(\text{d}_C) \to \Im(\text{d}_A)$ 为满射；
再由假设可知 $\Im(\text{d}_C) \to \Im(\text{d}_B)$ 为满射。
由 (1)，$C \to B$ 为满射。

\medskip\noindent
设 (2)(a) 成立。则 $f$ 为满射，且
$\Omega_{B/\Lambda} \otimes_B k \to \Omega_{A/\Lambda} \otimes_A k$
为满射。令 $K$ 为其核。由
(\ref{formal-defos-equation-sequence-functorial})，注意到
$K = \text{d}_B(\Ker(\mathfrak m_B/\mathfrak m_B^2 \to
\mathfrak m_A/\mathfrak m_A^2))$。
选取 $k$-向量空间的分裂
$$
\Omega_{B/\Lambda} \otimes_B k =
\Omega_{A/\Lambda} \otimes_A k \oplus K
$$
映射 $\text{d} : B \to \Omega_{B/\Lambda}$ 经投影到 $K$
诱导映射 $D : B \to K$。令
$C = \{b \in B \mid D(b) = 0\}$。Leibniz 法则表明，
这是 $B$ 的一个 $\Lambda$-子代数。取 $\overline{x} \in k$，
再取映到 $\overline{x}$ 的 $x \in B$。若 $D(x) \not = 0$，
则可找到 $y \in \mathfrak m_B$，使 $D(y) = D(x)$。
于是 $x-y \in C$ 映到 $\overline{x}$。因此 $C \to k$ 为满射，
且 $C$ 是 $\mathcal{C}_\Lambda$ 的对象。类似地，取
$\omega \in \Im(\text{d}_A)$。由 (1)，可找到
$x \in \mathfrak m_B$，使 $\text{d}_B(x)$ 映到 $\omega$。
若 $D(x) \not = 0$，则可找到 $y \in \mathfrak m_B$，
它在 $\mathfrak m_A/\mathfrak m_A^2$ 中映为零，且
$D(y) = D(x)$。于是 $z=x-y$ 属于 $\mathfrak m_C$，而其像
$\text{d}_C(z) \in \Omega_{C/k} \otimes_C k$% P10-E0424
映到 $\omega$。故
$\Im(\text{d}_C) \to \Im(\text{d}_A)$ 为满射。
由 (1)，$C \to A$ 为满射；再由假设，$C \to B$ 为满射。
所以 $D=0$，也就是 $K=0$，即 (2)(c) 成立。
这完成了 (2) 的证明。

\medskip\noindent
(3) 的证明。若 $k'/k$ 可分，% P10-E0425
则由引理 \ref{formal-defos-lemma-H1-separable-case}，
$H_1(L_{k/\Lambda}) =
\mathfrak m_\Lambda/\mathfrak m_\Lambda^2 \otimes_{k'} k$。
于是
$\Im(\text{d}_A) =
\mathfrak m_A/(\mathfrak m_\Lambda A + \mathfrak m_A^2)$，
对 $B$ 也同样成立。因此 (3) 由 (2) 得出。

\medskip\noindent
(4) 的证明。$f$ 的截面 $s$ 不是满射（按定义，小扩张有非平凡核），
故 $f$ 不是本质满射。反过来，设 $f$ 是小扩张但不是本质满射。
选取 $\mathcal{C}_\Lambda$ 中的非满环映射 $C \to B$，
使 $C \to A$ 为满射。令 $C' \subset B$ 为 $C \to B$ 的像。
则 $C' \not = B$，但 $C'$ 满射到 $A$。由于
$f : B \to A$ 是小扩张，
$\text{length}_C(B) = \text{length}_C(A) + 1$。
又因 $C'$ 是 $B$ 的真子环，
$\text{length}_C(C') \leq \text{length}_C(A)$。
但 $C' \to A$ 为满射，所以事实上
$\text{length}_C(C') = \text{length}_C(A)$，且
$C' \to A$ 为同构。这便给出所需截面。
\end{proof}

\begin{example}
\label{formal-defos-example-essential-surjection}
设 $\Lambda = k[[x]]$ 是 $k$ 上 $1$ 个变量的幂级数环。
令 $A=k$、$B=\Lambda/(x^2)$。则由引理
\ref{formal-defos-lemma-essential-surjection}，$B \to A$ 是本质满射；
这是因为它是小扩张，且映射 $B \to A$ 在范畴
$\mathcal{C}_\Lambda$ 中没有右逆。然而映射
$$
k \cong \mathfrak m_B/\mathfrak m_B^2
\longrightarrow
\mathfrak m_A/\mathfrak m_A^2 = 0
$$
不是同构。因此，在引理 \ref{formal-defos-lemma-essential-surjection} (3) 中，
必须考虑相对余切空间之间的映射
$\mathfrak m_B/(\mathfrak m_\Lambda B + \mathfrak m_B^2) \to
\mathfrak m_A/(\mathfrak m_\Lambda A + \mathfrak m_A^2)$。
\end{example}


\section{完备化的基范畴}
\label{formal-defos-section-category-completion-CLambda}

\noindent
下面对范畴 $\mathcal{C}_\Lambda$ 的“完备化”，将作为
$\mathcal{C}_\Lambda$ 上群胚余纤维范畴之完备化的基范畴
（第 \ref{formal-defos-section-formal-objects} 节）。

\begin{definition}
\label{formal-defos-definition-completion-CLambda}
设 $\Lambda$ 是诺特环，且 $\Lambda \to k$ 是到域 $k$ 的有限环映射。
定义{\it $\widehat{\mathcal{C}}_\Lambda$} 为如下范畴：
\begin{enumerate}
\item 对象为偶 $(R, \varphi)$，其中 $R$ 是诺特完备局部
$\Lambda$-代数，且 $\varphi : R/\mathfrak m_R \to k$ 是
$\Lambda$-代数同构；
\item 态射 $f : (S, \psi) \to (R, \varphi)$ 为局部
$\Lambda$-代数同态，且满足
$\varphi \circ (f \bmod \mathfrak m) = \psi$。
\end{enumerate}
\end{definition}

\noindent
与定义 \ref{formal-defos-definition-CLambda} 之后的讨论一样，
我们通常把 $\widehat{\mathcal{C}}_\Lambda$ 的对象简记为 $R$，
并约定识别 $R/\mathfrak m_R = k$。
本节讨论范畴 $\widehat{\mathcal{C}}_\Lambda$ 的对象与态射的一些基本性质，
这些性质与上一节关于范畴 $\mathcal{C}_\Lambda$ 的讨论相对应。

\medskip\noindent
首先注意，任意对象 $A \in \mathcal{C}_\Lambda$ 也是
$\widehat{\mathcal{C}}_\Lambda$ 的对象，因为阿廷局部环总是诺特的，
并且关于其极大理想完备（该理想毕竟是幂零理想）。此外，由定义显然可见
$\mathcal{C}_\Lambda \subset \widehat{\mathcal{C}}_\Lambda$
是由所有阿廷环构成的严格满子范畴。反过来，事实上
$\widehat{\mathcal{C}}_\Lambda$ 的每个对象都是
$\mathcal{C}_\Lambda$ 中对象的极限。

\medskip\noindent
设 $R$ 是 $\widehat{\mathcal{C}}_\Lambda$ 的对象。
对 $n \in \mathbf{N}$，考虑环 $R_n = R/\mathfrak m_R^n$。
它们是具有唯一幂零素理想的诺特局部环，因而是阿廷环；见
《代数》命题 \ref{algebra-proposition-dimension-zero-ring}。
环映射
$$
\ldots \to R_{n + 1} \to R_n \to \ldots \to R_2 \to R_1 = k
$$
全都满射。按定义，$R$ 的完备性意即 $R = \lim R_n$。
若 $f : R \to S$ 是 $\widehat{\mathcal{C}}_\Lambda$ 中的环映射，
则得到一族环映射 $f_n : R_n \to S_n$，其极限就是给定映射。

\begin{lemma}
\label{formal-defos-lemma-surjective-cotangent-space}
设 $f: R \to S$ 是 $\widehat{\mathcal{C}}_\Lambda$ 中的环映射。
下列条件等价：
\begin{enumerate}
\item $f$ 为满射；
\item 映射
$\mathfrak m_R/\mathfrak m_R^2 \to \mathfrak m_S/\mathfrak m_S^2$
为满射；
\item 映射
$\mathfrak m_R/(\mathfrak m_\Lambda R + \mathfrak m_R^2) \to
\mathfrak m_S/(\mathfrak m_\Lambda S + \mathfrak m_S^2)$
为满射。
\end{enumerate}
\end{lemma}

\begin{proof}
注意，对 $n \geq 2$，相对余切空间之间有等式
$$
\mathfrak m_R/(\mathfrak m_\Lambda R + \mathfrak m_R^2)
=
\mathfrak m_{R_n}/(\mathfrak m_\Lambda R_n + \mathfrak m_{R_n}^2)
$$
对 $S$ 也同样成立。因此由引理 \ref{formal-defos-lemma-surjective}，
$R_n \to S_n$ 对所有 $n$ 都是满射。令 $K_n$ 为
$R_n \to S_n$ 的核。则序列
$$
0 \to K_n \to R_n \to S_n \to 0
$$
组成逆系统的正合序列。由于每个 $K_n$ 都是阿廷模，系统 $(K_n)$
满足 Mittag-Leffler 条件。因此由《代数》引理
\ref{algebra-lemma-ML-exact-sequence}，取极限保持正合性。
所以 $\lim R_n \to \lim S_n$ 为满射，即 $f$ 为满射。
\end{proof}

\begin{lemma}
\label{formal-defos-lemma-CLambdahat-pushouts}
范畴 $\widehat{\mathcal{C}}_\Lambda$ 有推出。
\end{lemma}

\begin{proof}
设 $R \to S_1$ 和 $R \to S_2$ 是
$\widehat{\mathcal{C}}_\Lambda$ 中的态射。考虑环
$C = S_1 \otimes_R S_2$。这个环有有限生成的极大理想
$\mathfrak m = \mathfrak m_{S_1} \otimes S_2 +
S_1 \otimes \mathfrak m_{S_2}$，其剩余域为 $k$。
令 $C^\wedge$ 为 $C$ 关于 $\mathfrak m$ 的完备化。
则 $C^\wedge$ 是关于极大理想
$\mathfrak m^\wedge = \mathfrak mC^\wedge$ 完备的诺特环，
其剩余域与 $k$ 识别；见《代数》引理
\ref{algebra-lemma-completion-Noetherian}。
因此 $C^\wedge$ 是 $\widehat{\mathcal{C}}_\Lambda$ 的对象。
于是 $S_1 \to C^\wedge$ 和 $S_2 \to C^\wedge$ 使
$C^\wedge$ 成为 $\widehat{\mathcal{C}}_\Lambda$ 中
$R$ 上的推出（细节从略）。
\end{proof}

\noindent
我们不会用到下面的引理。

\begin{lemma}
\label{formal-defos-lemma-CLambdahat-coproducts}
范畴 $\widehat{\mathcal{C}}_\Lambda$ 有对象对的余积。
\end{lemma}

\begin{proof}
设 $R$ 和 $S$ 是 $\widehat{\mathcal{C}}_\Lambda$ 的对象。
考虑环 $C = R \otimes_\Lambda S$。有典范满射
$C \to R \otimes_\Lambda S \to k \otimes_\Lambda k \to k$，
其中最后一个映射为乘法映射。$C \to k$ 的核是极大理想
$\mathfrak m$。注意，$\mathfrak m$ 由 $\mathfrak m_R C$、
$\mathfrak m_S C$，以及 $C$ 中有限多个映到
$k \otimes_\Lambda k \to k$ 之核的生成元的元素所生成。
因此 $\mathfrak m$ 是有限生成理想。令 $C^\wedge$ 为
$C$ 关于 $\mathfrak m$ 的完备化。则 $C^\wedge$ 是关于极大理想
$\mathfrak m^\wedge = \mathfrak mC^\wedge$ 完备的诺特环，
剩余域为 $k$；见《代数》引理
\ref{algebra-lemma-completion-Noetherian}。
因此 $C^\wedge$ 是 $\widehat{\mathcal{C}}_\Lambda$ 的对象。
于是 $R \to C^\wedge$ 和 $S \to C^\wedge$ 使
$C^\wedge$ 成为 $\widehat{\mathcal{C}}_\Lambda$ 中的余积
（细节从略）。
\end{proof}

\noindent
范畴中的空余积就是该范畴的始对象。在经典情形下，
$\widehat{\mathcal{C}}_\Lambda$ 有始对象，即 $\Lambda$ 本身。
更一般地，若 $k' = k$，则 $\Lambda$ 关于
$\mathfrak m_\Lambda$ 的完备化 $\Lambda^\wedge$ 是始对象。
还可以更一般地说，若 $k' \subset k$ 可分，则
$\widehat{\mathcal{C}}_\Lambda$ 也有始对象。具体地，选取首一多项式
$P \in \Lambda[T]$，使 $k \cong k'[T]/(P')$，其中
$p' \in k'[T]$ 是 $P$ 的像。% P10-E0426
则 $R = \Lambda^\wedge[T]/(P)$ 是始对象；见引理
\ref{formal-defos-lemma-fiber-product-CLambda} 的证明。

\medskip\noindent
若 $R$ 是上述始对象，则有
$\mathcal{C}_\Lambda = \mathcal{C}_R$ 和
$\widehat{\mathcal{C}}_\Lambda = \widehat{\mathcal{C}}_R$，
这实际上把本章的全部讨论归结到经典情形。
但是，若 $k' \subset k$ 不可分，则始对象不存在。

\begin{lemma}
\label{formal-defos-lemma-derivations-finite}
设 $S$ 是 $\widehat{\mathcal{C}}_\Lambda$ 的对象。则
$\dim_k \text{Der}_\Lambda(S, k) < \infty$。
\end{lemma}

\begin{proof}
选取 $x_1, \ldots, x_n \in \mathfrak m_S$，使其像构成相对余切空间
$\mathfrak m_S/(\mathfrak m_\Lambda S + \mathfrak m_S^2)$
的一组 $k$-基。选取 $y_1, \ldots, y_m \in S$，使其在 $k$ 中的像
在 $k'$ 上生成 $k$。我们断言
$\dim_k \text{Der}_\Lambda(S, k) \leq n + m$。
为证明这一点，只需证明：若 $D(x_i) = 0$ 且 $D(y_j) = 0$，
则 $D = 0$。设 $a \in S$。可找到多项式
$P = \sum \lambda_J y^J$，其中 $\lambda_J \in \Lambda$，
使其在 $k$ 中的像与 $a$ 在 $k$ 中的像相同。
由对所有 $j$ 都有 $D(y_j) = 0$ 这一假设可知
$D(a - P) = D(a) - D(P) = D(a)$。因而可以假设
$a \in \mathfrak m_S$。写成 $a = \sum a_i x_i$，其中
$a_i \in S$。由 Leibniz 法则，
$$
D(a) = \sum x_iD(a_i) + \sum a_iD(x_i) = \sum x_iD(a_i)
$$
因为我们假设了 $D(x_i) = 0$。又因 $x_i$ 在 $k$ 上的乘法为零，
所以 $\sum x_iD(a_i) = 0$。
\end{proof}

\begin{lemma}
\label{formal-defos-lemma-derivations-surjective}
设 $f : R \to S$ 是 $\widehat{\mathcal{C}}_\Lambda$ 的态射。
若 $\text{Der}_\Lambda(S, k) \to \text{Der}_\Lambda(R, k)$ 为单射，
则 $f$ 为满射。
\end{lemma}

\begin{proof}
若 $f$ 不是满射，则由引理
\ref{formal-defos-lemma-surjective-cotangent-space}，
$\mathfrak m_S/(\mathfrak m_R S + \mathfrak m_S^2)$ 非零。
因而 $Q = S/(f(R) + \mathfrak m_R S + \mathfrak m_S^2)$ 也非零。
通过 $f$，$Q$ 是 $k = R/\mathfrak m_R$-向量空间。
再通过 $S \to k$ 把 $Q$ 视为 $S$-模。商映射
$D : S \to Q$ 是一个 $R$-导子：若 $a_1, a_2 \in S$，
则可写成
$a_1 = f(b_1) + a_1'$ 和 $a_2 = f(b_2) + a_2'$，
其中 $b_1, b_2 \in R$ 且
$a_1', a_2' \in \mathfrak m_S$。此时对 $i = 1, 2$，
$b_i$ 和 $a_i$ 在 $k$ 中有相同的像，并且在 $Q$ 中
\begin{align*}
a_1a_2 & = (f(b_1) + a_1')(f(b_2) + a_2') \\
& = f(b_1)a_2' + f(b_2)a_1' \\
& = f(b_1)(f(b_2) + a_2') + f(b_2)(f(b_1) + a_1') \\
& = f(b_1)a_2 + f(b_2)a_1
\end{align*}
这就证明了 Leibniz 法则。因此 $D : S \to Q$ 是一个
$\Lambda$-导子，并且与 $R \to S$ 复合为零。
由于 $Q \not = 0$，也存在与 $R \to S$ 复合为零的导子
$D : S \to k$；也就是说，
$\text{Der}_\Lambda(S, k) \to \text{Der}_\Lambda(R, k)$
不是单射。
\end{proof}

\begin{lemma}
\label{formal-defos-lemma-m-adic-topology}
设 $R$ 是 $\widehat{\mathcal{C}}_\Lambda$ 的对象，$(J_n)$ 是满足
$\mathfrak m_R^n \subset J_n$ 的理想递减列。令
$J = \bigcap J_n$。则序列 $(J_n/J)$ 在 $R/J$ 上定义
$\mathfrak m_{R/J}$-进拓扑。
\end{lemma}

\begin{proof}
显然 $\mathfrak m_{R/J}^n \subset J_n/J$。因此只需证明：
对每个 $n$，存在 $N$，使
$J_N/J \subset \mathfrak m_{R/J}^n$。这等价于
$J_N \subset \mathfrak m_R^n + J$。对每个 $n$，环
$R/\mathfrak m_R^n$ 都是阿廷环，故存在 $N_n$，使
$$
J_{N_n} + \mathfrak m_R^n = J_{N_n + 1} + \mathfrak m_R^n = \ldots
$$
令 $E_n = (J_{N_n} + \mathfrak m_R^n)/\mathfrak m_R^n$，
并令 $E = \lim E_n \subset \lim R/\mathfrak m_R^n = R$。
注意 $E \subset J$：因为对任意 $f \in E$ 和任意 $m$，
当所有 $n \gg 0$ 时均有
$f \in J_m + \mathfrak m_R^n$，故由 Krull 交定理
$f \in J_m$；见《代数》引理
\ref{algebra-lemma-intersect-powers-ideal-module-zero}。
由于过渡映射 $E_n \to E_{n - 1}$ 全都满射，
$J$ 满射到 $E_n$。所以取 $N = N_n$ 即可。% P10-E0427
\end{proof}

\begin{lemma}
\label{formal-defos-lemma-limit-artinian}
设 $\ldots \to A_3 \to A_2 \to A_1$ 是
$\mathcal{C}_\Lambda$ 中的一列满环映射。若
$\dim_k (\mathfrak m_{A_n}/\mathfrak m_{A_n}^2)$ 有界，则
$S = \lim A_n$ 是 $\widehat{\mathcal{C}}_\Lambda$ 的对象，
并且理想 $I_n = \Ker(S \to A_n)$ 在 $S$ 上定义
$\mathfrak m_S$-进拓扑。
\end{lemma}

\begin{proof}
下文直接使用映射 $S \to A_n$ 对所有 $n$ 均为满射这一事实。
注意，映射
$\mathfrak m_{A_{n + 1}}/\mathfrak m_{A_{n + 1}}^2 \to
\mathfrak m_{A_n}/\mathfrak m_{A_n}^2$ 为满射；见引理
\ref{formal-defos-lemma-surjective-cotangent-space}。
因此，当 $n$ 充分大时，维数
$\dim_k (\mathfrak m_{A_n}/\mathfrak m_{A_n}^2)$ 稳定为某个整数，
记为 $r$。于是可找到
$x_1, \ldots, x_r \in \mathfrak m_S$，使其在 $A_n$ 中的像生成
$\mathfrak m_{A_n}$。此外，选取 $y_1, \ldots, y_t \in S$，
使其在 $k$ 中的像在 $\Lambda$ 上生成 $k$。从而得到环映射
$P = \Lambda[z_1, \ldots, z_{r + t}] \to S$，
$z_i \mapsto x_i$ 且 $z_{r + j} \mapsto y_j$，
使复合 $P \to S \to A_n$ 对所有 $n$ 都为满射。
令 $\mathfrak m \subset P$ 为 $P \to k$ 的核。
令 $R = P^\wedge$ 为 $P$ 的 $\mathfrak m$-进完备化；
这是 $\widehat{\mathcal{C}}_\Lambda$ 的对象。
仍有相容的满射系 $R \to A_n$，故得到映射 $R \to S$。
令 $J_n = \Ker(R \to A_n)$，并令 $J = \bigcap J_n$。
由引理 \ref{formal-defos-lemma-m-adic-topology}，
$R/J = \lim R/J_n = \lim A_n = S$，而且理想
$J_n/J = I_n$ 定义 $\mathfrak m$-进拓扑。
（注意，对每个 $n$，都有
$\mathfrak m_R^{N_n} \subset J_n$，其中 $N_n$ 依指标而定，
未必有 $N_n = n$；所以在应用该引理之前，可能需要重编号理想
$J_n$。）
\end{proof}

\begin{lemma}
\label{formal-defos-lemma-power-series}
设 $R', R \in \Ob(\widehat{\mathcal{C}}_\Lambda)$。
假设对 $R$ 的某个理想 $I$ 有 $R = R' \oplus I$。
取 $x_1, \ldots, x_r \in I$，使其像构成
$I/\mathfrak m_R I$ 的一组基。令
$S = R'[[X_1, \ldots, X_r]]$，并考虑把 $X_i$ 映到 $x_i$ 的
$R'$-代数映射 $S \to R$。假设对每个 $n \gg 0$，映射
$S/\mathfrak m_S^n \to R/\mathfrak m_R^n$ 在
$\mathcal{C}_\Lambda$ 中都有左逆。则 $S \to R$ 为同构。
\end{lemma}

\begin{proof}
由 $R = R' \oplus I$，有
$$
\mathfrak m_R/\mathfrak m_R^2 =
\mathfrak m_{R'}/\mathfrak m_{R'}^2 \oplus I/\mathfrak m_RI
$$
类似地，
$$
\mathfrak m_S/\mathfrak m_S^2 =
\mathfrak m_{R'}/\mathfrak m_{R'}^2 \oplus \bigoplus kX_i
$$
因此，当 $n > 1$ 时，映射
$S/\mathfrak m_S^n \to R/\mathfrak m_R^n$
在余切空间上诱导同构。于是由引理
\ref{formal-defos-lemma-surjective-cotangent-space}，左逆
$h_n : R/\mathfrak m_R^n \to S/\mathfrak m_S^n$ 为满射。
又因 $h_n$ 作为左逆是单射，它是同构。
所以典范满射 $S/\mathfrak m_S^n \to R/\mathfrak m_R^n$
全都是同构，结论得证。
\end{proof}

\section{群胚余纤维范畴}
\label{formal-defos-section-preliminary}

\noindent
在发展这套理论时，我们使用群胚{\it 余纤维}范畴。
我们假定群胚{\it 纤维}范畴的定义和基本性质为已知；见
《范畴》第 \ref{categories-section-fibred-groupoids} 节。

\begin{definition}
\label{formal-defos-definition-category-cofibred-groupoids}
设 $\mathcal{C}$ 是范畴。$\mathcal{C}$ 上的一个
{\it 群胚余纤维范畴}，是配备函子
$p: \mathcal{F} \to \mathcal{C}$ 的范畴 $\mathcal{F}$，
使得经由
$p^{opp}: \mathcal{F}^{opp} \to \mathcal{C}^{opp}$，
$\mathcal{F}^{opp}$ 成为 $\mathcal{C}^{opp}$ 上的群胚纤维范畴。
\end{definition}

\noindent
明确地说，若下列两个条件成立，则
$p: \mathcal{F} \to \mathcal{C}$ 是群胚余纤维范畴：
\begin{enumerate}
\item 对 $\mathcal{C}$ 中每个态射 $f: U \to V$ 以及
$U$ 上的每个对象 $x$，存在 $\mathcal{F}$ 中位于 $f$ 上方的态射
$x \to y$。
\item 对 $\mathcal{F}$ 中每一对态射
$a: x \to y$ 和 $b: x \to z$，以及满足
$p(b) = f \circ p(a)$ 的任意态射
$f: p(y) \to p(z)$，存在唯一的 $\mathcal F$ 中态射
$c: y \to z$，它位于 $f$ 上方且满足 $b = c \circ a$。
\end{enumerate}

\begin{remarks}
\label{formal-defos-remarks-cofibered-groupoids}
关于群胚纤维范畴的一切内容都直接转化到余纤维情形。
下面这些说明用于固定记号。设 $\mathcal{C}$ 是范畴。
\begin{enumerate}
\item 我们常从记号中省略函子
$p: \mathcal{F} \to \mathcal{C}$。
\item $\mathcal{C}$ 中对象 $U$ 上方的纤维范畴记作
$\mathcal{F}(U)$。其对象是 $\mathcal{F}$ 中位于 $U$ 上方的对象，
其态射是 $\mathcal{F}$ 中位于 $\text{id}_U$ 上方的态射。
若 $x, y$ 是 $\mathcal{F}(U)$ 的对象，有时把
$\Mor_{\mathcal{F}(U)}(x, y)$ 写作 $\Mor_U(x, y)$。
\item 纤维范畴 $\mathcal{F}(U)$ 是群胚；见《范畴》引理
\ref{categories-lemma-fibred-groupoids}。
因此 $\mathcal{F}(U)$ 中的所有态射都是同构。
有时把 $\Mor_{\mathcal{F}(U)}(x, x)$ 写作
$\text{Aut}_U(x)$。
\item
\label{formal-defos-item-pushforward}
设 $\mathcal{F}$ 是 $\mathcal{C}$ 上的群胚余纤维范畴，
$f: U \to V$ 是 $\mathcal{C}$ 中的态射，且
$x \in \Ob(\mathcal{F}(U))$。
$x$ 沿 $f$ 的一个{\it 推前}，是 $\mathcal{F}$ 中位于
$f$ 上方的态射 $x \to y$。推前在唯一同构意义下唯一
（见《范畴》定义 \ref{categories-definition-cartesian-over-C}
之后的讨论）。有时把 $x$ 沿 $f$ 的“这个”推前写作
$x \to f_*x$。
\item $\mathcal{F}$ 的一个{\it 推前选取}，是对上述每一对
$(x, f)$ 选取 $x$ 沿 $f$ 的一个推前。
由选择公理，可以为 $\mathcal{F}$ 作出这样的推前选取。
\item 设 $\mathcal{F}$ 是 $\mathcal{C}$ 上的群胚余纤维范畴。
给定 $\mathcal{F}$ 的一个推前选取，就有相伴的伪函子
$\mathcal{C} \to \textit{Groupoids}$。
我们绝不会用到这一构造，故略去细节。
\item
\label{formal-defos-item-cofibered-morphism}
$\mathcal{C}$ 上群胚余纤维范畴之间的态射，是与到
$\mathcal{C}$ 的投影交换的函子。若 $\mathcal{F}$ 和
$\mathcal{F}'$ 是 $\mathcal{C}$ 上的群胚余纤维范畴，
则从 $\mathcal{F}$ 到 $\mathcal{F}'$ 的态射记为
$\Mor_\mathcal{C}(\mathcal{F}, \mathcal{F}')$。
\item
\label{formal-defos-item-definition-cofibered-groupoids-2-category}
群胚余纤维范畴组成一个 $(2, 1)$-范畴
$\text{Cof}(\mathcal{C})$。它的 1-态射是
(\ref{formal-defos-item-cofibered-morphism}) 中所述的态射。
若 $p : \mathcal{F} \to C$ 和
$p': \mathcal{F}' \to \mathcal{C}$ 是群胚余纤维范畴，% P10-E0428
而 $\varphi, \psi : \mathcal{F} \to \mathcal{F}'$ 是
$1$-态射，则 2-态射 $t : \varphi \to \psi$ 是满足下述条件的
函子态射：对所有 $x \in \Ob(\mathcal{F})$，
$p'(t_x) = \text{id}_{p(x)}$。
\item
\label{formal-defos-item-construction-associated-cofibered-groupoid}
设 $F : \mathcal{C} \to \textit{Groupoids}$ 是函子。
如下可构造与 $F$ 相伴的群胚余纤维范畴
$\mathcal{F} \to \mathcal{C}$。
$\mathcal{F}$ 的对象是偶 $(U, x)$，其中
$U \in \Ob(\mathcal{C})$ 且 $x \in \Ob(F(U))$。
一个态射 $(U, x) \to (V, y)$ 是偶 $(f, a)$，其中
$f \in \Mor_\mathcal{C}(U, V)$ 且
$a \in \Mor_{F(V)}(F(f)(x), y)$。
函子 $\mathcal{F} \to \mathcal{C}$ 把 $(U, x)$ 映到 $U$。
见《范畴》第 \ref{categories-section-presheaves-groupoids} 节。
\item
\label{formal-defos-item-associated-functor-isomorphism-classes}
设 $\mathcal{F}$ 是 $\mathcal{C}$ 上的群胚余纤维范畴。
对 $U \in \Ob(\mathcal{C})$，令
$\overline{\mathcal{F}}(U)$ 等于范畴 $\mathcal{F}(U)$
的同构类集合。% P10-E0429
若 $f : U \to V$ 是 $\mathcal{C}$ 的态射，
则把 $x$ 的同构类映到 $x$ 的一个推前 $f_*x$ 的同构类，
便得到集合映射
$\overline{\mathcal{F}}(U) \to \overline{\mathcal{F}}(V)$；
见 (\ref{formal-defos-item-pushforward})。% P10-E0430
于是 $\overline{\mathcal{F}} : \mathcal{C} \to \textit{Sets}$
是函子。类似地，若
$\varphi : \mathcal{F} \to \mathcal{G}$ 是余纤维范畴之间的态射，
则以
$\overline{\varphi}: \overline{\mathcal{F}} \to  \overline{\mathcal{G}}$
表示相伴的函子态射。
\item
\label{formal-defos-item-convention-cofibered-sets}
设 $F: \mathcal{C} \to \textit{Sets}$ 是函子。
可以把集合视为离散范畴，即只有恒等态射的群胚。
于是构造 (\ref{formal-defos-item-construction-associated-cofibered-groupoid})
给出与 $F$ 相伴的集合余纤维范畴。
这定义了从函子范畴
$\mathcal{C} \to \textit{Sets}$ 到 $\mathcal{C}$ 上
群胚余纤维范畴之范畴的一个充分忠实嵌入。
我们把函子范畴与它在这一嵌入下的像识别。
因此，若 $F : \mathcal{C} \to \textit{Sets}$ 是函子，
则仍以 $F$ 表示相伴的集合余纤维范畴；若
$\varphi : F \to G$ 是函子态射，则仍以 $\varphi$
表示对应的集合余纤维范畴态射，反之亦然。
见《范畴》第 \ref{categories-section-fibred-in-sets} 节。
\item
\label{formal-defos-item-definition-yoneda}
设 $U$ 是 $\mathcal{C}$ 的对象。以 $\underline{U}$ 表示函子
$\Mor_\mathcal{C}(U, -): \mathcal{C} \to
\textit{Sets}$。这定义了从 $\mathcal
C^{opp}$ 到函子范畴
$\mathcal{C} \to \textit{Sets}$ 的充分忠实嵌入。
因此，若 $f : U \to V$ 是态射，则仍以 $f$ 表示所诱导的态射
$\underline{V} \to \underline{U}$，反之亦然。
\item
\label{formal-defos-item-fibre-product}
群胚余纤维范畴的纤维积：若
$\mathcal{F} \to \mathcal{H}$ 和
$\mathcal{G} \to \mathcal{H}$ 是
$\mathcal{C}_\Lambda$ 上群胚余纤维范畴之间的态射，
则可把它们作为 $\mathcal{C}_\Lambda$ 上范畴的 2-纤维积构造，
用作其 2-纤维积；见《范畴》引理
\ref{categories-lemma-2-product-categories-over-C} 中的描述。
\item
\label{formal-defos-item-product}
群胚余纤维范畴的积：若 $\mathcal{F}$ 和 $\mathcal{G}$ 是
$\mathcal{C}_\Lambda$ 上的群胚余纤维范畴，则按定义，
其积为 $2$-纤维积
$\mathcal{F} \times_{\mathcal{C}_\Lambda} \mathcal{G}$；
见《范畴》引理
\ref{categories-lemma-2-product-categories-over-C} 中的描述。
\item
\label{formal-defos-item-definition-restricting-base-category}
限制基范畴：设 $p : \mathcal{F} \to \mathcal{C}$ 是
群胚余纤维范畴，且 $\mathcal{C}'$ 是 $\mathcal{C}$ 的满子范畴。
限制 $\mathcal{F}|_{\mathcal{C}'}$ 是 $\mathcal{F}$ 的满子范畴，
其对象是位于 $\mathcal{C}'$ 对象上方的那些对象。
经由函子
$p|_{\mathcal{C}'}: \mathcal{F}|_{\mathcal{C}'} \to \mathcal{C}'$，
它是群胚余纤维范畴。
\end{enumerate}
\end{remarks}

\section{pro-可表函子与预形变范畴}
\label{formal-defos-section-cofibered-groupoids}

\noindent
我们的基本目标，是理解 $\mathcal{C}_\Lambda$ 和
$\widehat{\mathcal{C}}_\Lambda$ 上的群胚余纤维范畴。
由于 $\mathcal{C}_\Lambda$ 是
$\widehat{\mathcal{C}}_\Lambda$ 的满子范畴，可以把
$\widehat{\mathcal{C}}_\Lambda$ 上的群胚余纤维范畴限制到
$\mathcal{C}_\Lambda$；见说明
\ref{formal-defos-remarks-cofibered-groupoids}
(\ref{formal-defos-item-definition-restricting-base-category})。
尤其可对函子这样做，因而也可对可表函子这样做。
由此在 $\mathcal{C}_\Lambda$ 上所得的函子称为 pro-可表函子。

\begin{definition}
\label{formal-defos-definition-prorepresentable}
设 $F : \mathcal{C}_\Lambda \to \textit{Sets}$ 是函子。
若对某个 $R \in \Ob(\widehat{\mathcal{C}}_\Lambda)$，
存在函子同构
$F \cong \underline{R}|_{\mathcal{C}_\Lambda}$，
则称 $F$ 是{\it pro-可表的}。
\end{definition}

\noindent
注意，若 $F : \mathcal{C}_\Lambda \to \textit{Sets}$
由 $R \in \Ob(\widehat{\mathcal{C}}_\Lambda)$ pro-表示，则
$$
F(k) = \Mor_{\widehat{\mathcal{C}}_\Lambda}(R, k) = \{*\}
$$
是单元集。形变理论中出现的
$\mathcal{C}_\Lambda$ 上群胚余纤维范畴也常满足类似条件。% P10-E0431

\begin{definition}
\label{formal-defos-definition-predeformation-category}
一个{\it 预形变范畴} $\mathcal{F}$，是
$\mathcal{C}_\Lambda$ 上的群胚余纤维范畴，并且
$\mathcal{F}(k)$ 等价于恰有一个对象和一个态射的范畴；
也就是说，$\mathcal{F}(k)$ 至少含有一个对象，且任意两个对象之间
都有唯一态射。预形变范畴之间的一个{\it 态射}，是
$\mathcal{C}_\Lambda$ 上群胚余纤维范畴之间的态射。
\end{definition}

\noindent
预形变范畴具有如下特征。设
$x_0 \in \Ob(\mathcal{F}(k))$。则 $\mathcal{F}$ 的每个对象
都配备了到 $x_0$ 的唯一态射。事实上，若 $x$ 是
$\mathcal{F}$ 中位于 $A$ 上方的对象，则可以选取推前
$x \to q_*x$，其中 $q : A \to k$ 为商映射。
存在唯一同构 $q_*x \to x_0$，而复合
$x \to q_*x \to x_0$ 就是所需态射。

\begin{remark}
\label{formal-defos-remark-predeformation-functor}
若函子 $F: \mathcal{C}_\Lambda \to \textit{Sets}$ 的相伴集合
余纤维范畴是预形变范畴，即 $F(k)$ 是单元集，
则称其为{\it 预形变函子}。
因此，若 $\mathcal{F}$ 是预形变范畴，则
$\overline{\mathcal{F}}$ 是预形变函子。
\end{remark}

\begin{remark}
\label{formal-defos-remark-localize-cofibered-groupoid}
设 $p : \mathcal{F} \to \mathcal{C}_\Lambda$ 是群胚余纤维范畴，
且 $x \in \Ob(\mathcal{F}(k))$。以 $\mathcal{F}_x$
表示位于 $x$ 上的对象所成范畴。
$\mathcal{F}_x$ 的对象是箭头 $y \to x$。
$\mathcal{F}_x$ 中的一个态射
$(y \to x) \to (z \to x)$ 是交换图
$$
\xymatrix{
y \ar[rr] \ar[dr] & & z \ar[dl] \\
& x &
}
$$
有遗忘函子 $\mathcal{F}_x \to \mathcal{F}$。
把函子 $p_x : \mathcal{F}_x \to \mathcal{C}_\Lambda$ 定义为复合
$\mathcal{F}_x \to \mathcal{F} \xrightarrow{p} \mathcal{C}_\Lambda$。
于是 $p_x : \mathcal{F}_x \to \mathcal{C}_\Lambda$
是预形变范畴（证明从略）。由此，对
$\mathcal{C}_\Lambda$ 上的任意群胚余纤维范畴以及任意
$x \in \Ob(\mathcal{F}(k))$，都可得到一个预形变范畴。
\end{remark}

\section{形式对象与完备化范畴}
\label{formal-defos-section-formal-objects}

\noindent
本节讨论如何在 $\mathcal{C}_\Lambda$ 上的群胚余纤维范畴与
$\widehat{\mathcal{C}}_\Lambda$ 上的群胚余纤维范畴之间
往返转换。

\begin{definition}
\label{formal-defos-definition-formal-objects}
设 $\mathcal{F}$ 是 $\mathcal{C}_\Lambda$ 上的群胚余纤维范畴。
$\mathcal{F}$ 的{\it 形式对象范畴 $\widehat{\mathcal{F}}$}
是具有如下对象和态射的范畴。
\begin{enumerate}
\item $\mathcal{F}$ 的一个
{\it 形式对象 $\xi = (R, \xi_n, f_n)$}，由
$\widehat{\mathcal{C}}_\Lambda$ 的一个对象 $R$、
一族以 $n \in \mathbf{N}$ 为指标的对象 $\xi_n$，它们属于
$\mathcal{F}(R/\mathfrak m_R^n)$，以及位于投影
$R/\mathfrak m_R^{n + 1} \to R/\mathfrak m_R^n$
上方的态射 $f_n : \xi_{n + 1} \to \xi_n$ 所组成。
\item 设 $\xi = (R, \xi_n, f_n)$ 和
$\eta = (S, \eta_n, g_n)$ 是 $\mathcal{F}$ 的形式对象。
形式对象的一个{\it 态射 $a : \xi \to \eta$}，由
$\widehat{\mathcal{C}}_\Lambda$ 中的映射
$a_0 : R \to S$，以及一族 $\mathcal{F}$ 中位于
$R/\mathfrak m_R^n \to S/\mathfrak m_S^n$ 上方的态射
$a_n : \xi_n \to \eta_n$ 所组成，并且对每个 $n$，图
$$
\xymatrix{
\xi_{n + 1} \ar[r]_{f_n} \ar[d]_{a_{n + 1}} & \xi_n \ar[d]^{a_n} \\
\eta_{n + 1} \ar[r]^{g_n} & \eta_n
}
$$
交换。
\end{enumerate}
\end{definition}

\noindent
形式对象范畴配备函子
$\widehat{p}: \widehat{\mathcal{F}} \to
\widehat{\mathcal{C}}_\Lambda$；它把对象
$(R, \xi_n, f_n)$ 映到 $R$，并把态射
$(R, \xi_n, f_n) \to (S, \eta_n, g_n)$ 映到 $R \to S$。

\begin{lemma}
\label{formal-defos-lemma-completion-cofibred}
设 $p : \mathcal{F} \to \mathcal{C}_\Lambda$ 是群胚余纤维范畴。
则
$\widehat{p} : \widehat{\mathcal{F}} \to
\widehat{\mathcal{C}}_\Lambda$
是群胚余纤维范畴。
\end{lemma}

\begin{proof}
设 $R \to S$ 是 $\widehat{\mathcal{C}}_\Lambda$ 中的环映射，
且 $(R, \xi_n, f_n)$ 是 $\widehat{\mathcal{F}}$ 的对象。
对每个 $n$，选取 $\xi_n$ 沿
$R/\mathfrak m_R^n \to S/\mathfrak m_S^n$ 的一个推前
$\xi_n \to \eta_n$。对每个 $n$，存在唯一的
$\mathcal{F}$ 中态射 $g_n : \eta_{n + 1} \to \eta_n$，
它位于 $S/\mathfrak m_S^{n + 1} \to S/\mathfrak m_S^n$
上方，且使图
$$
\xymatrix{
\xi_{n + 1} \ar[d] \ar[r]_{f_n} & \xi_n \ar[d] \\
\eta_{n + 1} \ar[r]^{g_n} & \eta_n
}
$$
交换（由群胚余纤维范畴的第一条公理）。% P10-E0432
因而得到位于 $R \to S$ 上方的态射
$(R, \xi_n, f_n) \to (S, \eta_n, g_n)$；
也就是说，群胚余纤维范畴的第一条公理对
$\widehat{\mathcal{F}}$ 成立。为验证第二条公理，设有
$\widehat{\mathcal{F}}$ 中的态射
$a : (R, \xi_n, f_n) \to (S, \eta_n, g_n)$ 和
$b : (R, \xi_n, f_n) \to (T, \theta_n, h_n)$，
以及 $\widehat{\mathcal{C}}_\Lambda$ 中满足
$c_0 \circ a_0 = b_0$ 的态射 $c_0 : S \to T$。
由 $\mathcal{F}$ 的群胚余纤维范畴第二条公理，
得到位于 $S/\mathfrak m_S^n \to T/\mathfrak m_T^n$
上方的唯一映射 $c_n : \eta_n \to \theta_n$，使
$c_n \circ a_n = b_n$。令 $c = (c_n)_{n \geq 0}$，
便得到所需的 $\widehat{\mathcal{F}}$ 中态射
$c : (S, \eta_n, g_n) \to (T, \theta_n, h_n)$
（验证 $h_n \circ c_{n + 1} = c_n \circ g_n$ 的步骤从略）。
\end{proof}

\begin{definition}
\label{formal-defos-definition-completion}
设 $p : \mathcal{F} \to \mathcal{C}_\Lambda$ 是群胚余纤维范畴。
群胚余纤维范畴
$\widehat{p} : \widehat{\mathcal  F} \to
\widehat{\mathcal{C}}_\Lambda$
称为 $\mathcal{F}$ 的{\it 完备化}。
\end{definition}

\noindent
若 $\mathcal{F}$ 是 $\mathcal C_\Lambda$ 上的群胚余纤维范畴，
我们已经用 $R$ 的极大理想幂滤过定义了
$R \in \Ob(\widehat{\mathcal{C}}_\Lambda)$ 时的
$\widehat{\mathcal{F}}(R)$。现在设
$\mathcal{I} = (I_n)$ 是由理想给出的 $R$ 的滤过，
它诱导 $\mathfrak{m}_R$-进拓扑。
定义 $\widehat{\mathcal{F}}_\mathcal{I}(R)$ 为具有如下对象和态射的范畴：
\begin{enumerate}
\item 对象是一族
$(\xi_n, f_n)_{n \in \mathbf{N}}$，其中
$\xi_n$ 是 $\mathcal{F}(R/I_n)$ 的对象，而
$f_n : \xi_{n + 1} \to \xi_n$ 是位于投影
$R/I_{n + 1} \to R/I_n$ 上方的态射。
\item 态射
$a : (\xi_n, f_n) \to (\eta_n, g_n)$ 由一族
$\mathcal{F}(R/I_n)$ 中的态射
$a_n : \xi_n \to \eta_n$ 组成，并且对每个 $n$，图
$$
\xymatrix{
\xi_{n + 1} \ar[r]^{f_n} \ar[d]_{a_{n + 1}} & \xi_n \ar[d]^{a_n} \\
\eta_{n + 1} \ar[r]^{g_n} & \eta_n
}
$$
交换。
\end{enumerate}

\begin{lemma}
\label{formal-defos-lemma-formal-objects-different-filtration}
在上述情形下，$\widehat{\mathcal{F}}_\mathcal{I}(R)$ 等价于范畴
$\widehat{\mathcal{F}}(R)$。
\end{lemma}

\begin{proof}
可如下定义一个等价
$\widehat{\mathcal{F}}_\mathcal{I}(R) \to
\widehat{\mathcal{F}}(R)$。对每个 $n$，令 $m(n)$ 为满足
$I_m \subset \mathfrak m_R^n$ 的最小 $m$。
给定 $\widehat{\mathcal{F}}_\mathcal{I}(R)$ 的对象
$(\xi_n, f_n)$，令 $\eta_n$ 为 $\xi_{m(n)}$ 沿
$R/I_{m(n)} \to R/\mathfrak m_R^n$ 的推前。
令 $g_n : \eta_{n + 1} \to \eta_n$ 为
$\mathcal{F}$ 中位于
$R/\mathfrak m_R^{n + 1} \to R/\mathfrak m_R^n$ 上方的唯一态射，
使图
$$
\xymatrix{
\xi_{m(n + 1)} \ar[rrr]_{f_{m(n)} \circ \ldots \circ f_{m(n + 1) - 1}} \ar[d]
& & & \xi_{m(n)} \ar[d] \\
\eta_{n + 1} \ar[rrr]^{g_n} & & & \eta_n
}
$$
交换（余纤维范畴的公理保证存在性和唯一性）。
函子
$\widehat{\mathcal{F}}_\mathcal{I}(R) \to
\widehat{\mathcal{F}}(R)$ 把 $(\xi_n, f_n)$ 映到
$(R, \eta_n, g_n)$。验证这确为范畴等价的步骤从略。
\end{proof}

\begin{remark}
\label{formal-defos-remark-different-sequence-ideals}
设 $p : \mathcal{F} \to \mathcal{C}_\Lambda$ 是群胚余纤维范畴。
假设对每个 $R \in \Ob(\widehat{\mathcal{C}}_\Lambda)$，
都给定了由理想构成的 $R$ 的滤过 $\mathcal{I}_R$。
若 $\mathcal{I}_R$ 对所有 $R$ 都在 $R$ 上诱导
$\mathfrak m_R$-进拓扑，则可模仿
$\widehat{\mathcal{F}}$ 的定义来定义范畴
$\widehat{\mathcal{F}}_\mathcal{I}$。
该范畴配备态射
$\widehat{p}_\mathcal{I} : \widehat{\mathcal{F}}_\mathcal{I} \to
\widehat{\mathcal{C}}_\Lambda$，从而成为群胚余纤维范畴，
并且 $\widehat{\mathcal{F}}_\mathcal{I}(R)$ 与上面定义的
$\widehat{\mathcal{F}}_{\mathcal{I}_R}(R)$ 同构。
在对象 $R \in \Ob(\widehat{\mathcal{C}}_\Lambda)$ 上使用引理
\ref{formal-defos-lemma-formal-objects-different-filtration} 的等价，
可知群胚余纤维范畴
$\widehat{\mathcal{F}}_\mathcal{I}$ 和
$\widehat{\mathcal{F}}$ 等价。
\end{remark}

\begin{remark}
\label{formal-defos-remark-completion-functor}
设 $F: \mathcal{C}_\Lambda \to \textit{Sets}$ 是函子。
把函子与集合余纤维范畴识别，则 $F$ 的完备化是函子
$\widehat{F} : \widehat{\mathcal{C}}_\Lambda \to \textit{Sets}$，
由 $\widehat{F}(S) = \lim F(S/\mathfrak{m}_S^{n})$ 给出。
这与 Schlessinger 论文 \cite{Sch} 中的定义一致。
\end{remark}

\begin{remark}
\label{formal-defos-remark-restrict-completion}
设 $\mathcal{F}$ 是 $\mathcal{C}_\Lambda$ 上的群胚余纤维范畴。
我们断言有典范等价
$$
can :
\widehat{\mathcal{F}}|_{\mathcal{C}_\Lambda}
\longrightarrow
\mathcal{F}.
$$
事实上，设 $A \in \Ob(\mathcal{C}_\Lambda)$，且
$(A, \xi_n, f_n)$ 是
$\widehat{\mathcal{F}}|_{\mathcal{C}_\Lambda}(A)$ 的对象。
由于 $A$ 是阿廷环，存在满足
$\mathfrak m_A^m = 0$ 的最小 $m \in \mathbf{N}$。
此时 $can$ 把 $(A, \xi_n, f_n)$ 映到 $\xi_m$。
由《范畴》引理
\ref{categories-lemma-equivalence-fibred-categories}，
这个函子是群胚余纤维范畴之间的等价；这是因为由引理
\ref{formal-defos-lemma-formal-objects-different-filtration}，它在每个纤维范畴上
都是等价，而且局部阿廷环 $A$ 上的
$\mathfrak m_A$-进拓扑来自零理想。
我们将经常通过函子 $can$ 的一个拟逆，把 $\mathcal{F}$
与 $\widehat{\mathcal{F}}$ 的一个满子范畴识别。
\end{remark}

\begin{remark}
\label{formal-defos-remark-completion-morphism}
设 $\varphi : \mathcal{F} \to \mathcal{G}$ 是
$\mathcal{C}_\Lambda$ 上群胚余纤维范畴之间的态射。
则诱导出 $\widehat{\mathcal{C}}_\Lambda$ 上群胚余纤维范畴之间的态射
$\widehat{\varphi}: \widehat{\mathcal{F}} \to
\widehat{\mathcal{G}}$。
它把 $\widehat{\mathcal{F}}$ 的对象
$\xi = (R, \xi_n, f_n)$ 映到
$(R, \varphi(\xi_n), \varphi(f_n))$，并把
$\widehat{\mathcal{F}}$ 的对象 $\xi$ 与 $\eta$ 之间的态射
$(a_0 : R \to S, a_n : \xi_n \to \eta_n)$ 映到
$(a_0 : R \to S, \varphi(a_n) : \varphi(\xi_n) \to \varphi(\eta_n))$。
最后，若 $t : \varphi \to \varphi'$ 是群胚余纤维范畴的
$1$-态射 $\varphi, \varphi': \mathcal{F} \to \mathcal{G}$
之间的 $2$-态射，则得到 $2$-态射
$\widehat{t} : \widehat{\varphi} \to \widehat{\varphi}'$。
具体地，对上述 $\xi = (R, \xi_n, f_n)$，令
$\widehat{t}_\xi = (t_{\varphi(\xi_n)})$。% P10-E0433
因此，完备化定义了 $2$-范畴之间的函子
$$
\widehat{~} :
\text{Cof}(\mathcal{C}_\Lambda)
\longrightarrow
\text{Cof}(\widehat{\mathcal{C}}_\Lambda)
$$
它从 $\mathcal{C}_\Lambda$ 上群胚余纤维范畴的 $2$-范畴
到 $\widehat{\mathcal{C}}_\Lambda$ 上群胚余纤维范畴的
$2$-范畴。
\end{remark}

\begin{remark}
\label{formal-defos-remark-completion-restriction-adjoint}
我们断言，说明 \ref{formal-defos-remark-completion-morphism} 的完备化函子与
说明 \ref{formal-defos-remarks-cofibered-groupoids}
(\ref{formal-defos-item-definition-restricting-base-category}) 的限制函子
$|_{\mathcal{C}_\Lambda} : \text{Cof}(\widehat{\mathcal{C}}_\Lambda)
\to \text{Cof}(\mathcal{C}_\Lambda)$
在如下精确意义下“2-伴随”。设
$\mathcal{F} \in \Ob(\text{Cof}(\mathcal{C}_\Lambda))$ 且
$\mathcal{G} \in \Ob(\text{Cof}(\widehat{\mathcal{C}}_\Lambda))$。
则有范畴等价
$$
\Phi :
\Mor_{\mathcal{C}_\Lambda}(
\mathcal{G}|_{\mathcal{C}_\Lambda}, \mathcal{F})
\longrightarrow
\Mor_{\widehat{\mathcal{C}}_\Lambda}(\mathcal{G}, \widehat{\mathcal{F}})
$$
为描述这一等价，我们如下定义典范态射
$\mathcal{G} \to
\widehat{\mathcal{G}|_{\mathcal{C}_\Lambda}}$ 和
$\widehat{\mathcal{F}}|_{\mathcal{C}_\Lambda} \to \mathcal{F}$：
\begin{enumerate}
\item 设 $R \in \Ob(\widehat{\mathcal{C}}_\Lambda))$，并设 $\xi$
为纤维范畴 $\mathcal{G}(R)$ 的对象。% P10-E0434
对每个 $n \in \mathbf{N}$，选取 $\xi$ 到
$R/\mathfrak m_R^n$ 的推前 $\xi \to \xi_n$，并令
$f_n : \xi_{n + 1} \to \xi_n$ 为所诱导的态射。
于是 $\mathcal{G} \to
\widehat{\mathcal{G}|_{\mathcal{C}_\Lambda}}$
把 $\xi$ 映到 $(R, \xi_n, f_n)$。
\item 这是说明 \ref{formal-defos-remark-restrict-completion} 的等价
$can : \widehat{\mathcal{F}}|_{\mathcal{C}_\Lambda} \to \mathcal{F}$。
\end{enumerate}
说明这些之后，等价
$\Phi : \Mor_{\mathcal{C}_\Lambda}(
\mathcal{G}|_{\mathcal{C}_\Lambda}, \mathcal{F})  \to
\Mor_{\widehat{\mathcal{C}}_\Lambda}(\mathcal{G},
\widehat{\mathcal{F}})$
把态射
$\varphi : \mathcal{G}|_{\mathcal{C}_\Lambda} \to \mathcal{F}$
映到
$$
\mathcal{G} \to \widehat{\mathcal{G}|_{\mathcal{C}_\Lambda}}
\xrightarrow{\widehat{\varphi}} \widehat{\mathcal{F}}
$$
$\Phi$ 有拟逆
$\Psi :
\Mor_{\widehat{\mathcal{C}}_\Lambda}(
\mathcal{G}, \widehat{\mathcal{F}}) \to
\Mor_{\mathcal{C}_\Lambda}(
\mathcal{G}|_{\mathcal{C}_\Lambda}, \mathcal{F})$，
它把 $\psi : \mathcal{G} \to \widehat{\mathcal{F}}$ 映到
$$
\mathcal{G}|_{\mathcal{C}_\Lambda}
\xrightarrow{\psi|_{\mathcal{C}_\Lambda}}
\widehat{\mathcal{F}}|_{\mathcal{C}_\Lambda} \to \mathcal{F}.
$$
验证 $\Phi$ 与 $\Psi$ 互为拟逆的步骤从略。
我们也不讨论 $\Phi$ 的函子性（因为这会进入我们无论如何都想避免的
3-范畴领域）。
\end{remark}

\begin{remark}
\label{formal-defos-remark-completion-restriction-cofset-adjoint}
对范畴 $\mathcal{C}$，以 $\text{CofSet}(\mathcal{C})$
表示 $\mathcal{C}$ 上集合余纤维范畴所成的范畴。
这是一个 $1$-范畴，与函子范畴
$\mathcal{C} \to \textit{Sets}$ 同构。% P10-E0435
见说明 \ref{formal-defos-remarks-cofibered-groupoids}
(\ref{formal-defos-item-convention-cofibered-sets})。
完备化函子和限制函子分别限制为函子
$\widehat{~} : \text{CofSet}(\mathcal{C}_\Lambda) \to
\text{CofSet}(\widehat{\mathcal{C}}_\Lambda)$ 和
$|_{\mathcal{C}_\Lambda} :
\text{CofSet}(\widehat{\mathcal{C}}_\Lambda) \to
\text{CofSet}(\mathcal{C}_\Lambda)$，
仍以相同符号表示。作为集合余纤维范畴之范畴上的函子，
完备化与限制在通常的 1-范畴意义下伴随：说明
\ref{formal-defos-remark-completion-restriction-adjoint} 中的同一构造，
对
$F \in \Ob(\text{CofSet}(\mathcal{C}_\Lambda))$ 和
$G \in \Ob(\text{CofSet}(\widehat{\mathcal{C}}_\Lambda))$
定义函子性双射
$$
\Mor_{\mathcal{C}_\Lambda}(G|_{\mathcal{C}_\Lambda}, F)
\longrightarrow
\Mor_{\widehat{\mathcal{C}}_\Lambda}(G, \widehat{F})
$$
同样，映射 $\widehat{F}|_{\mathcal{C}_\Lambda} \to F$
是同构。
\end{remark}

\begin{remark}
\label{formal-defos-remark-restrict-complete-continuous-functor}
设 $G : \widehat{\mathcal{C}}_\Lambda \to \textit{Sets}$
是与极限交换的函子。则说明
\ref{formal-defos-remark-completion-restriction-adjoint} 中所述的映射
$G \to \widehat{G|_{\mathcal{C}_\Lambda}}$ 是同构。
事实上，若 $S$ 是 $\widehat{\mathcal{C}}_\Lambda$ 的对象，
则有典范双射
$$
\widehat{G|_{\mathcal{C}_\Lambda}}(S) =
\lim_n G(S/\mathfrak{m}_S^n) =
G(\lim_n S/\mathfrak{m}_S^n) = G(S).
$$
特别地，若 $R$ 是 $\widehat{\mathcal{C}}_\Lambda$ 的对象，则
$\underline{R} =
\widehat{\underline{R}|_{\mathcal{C}_\Lambda}}$，
因为根据极限的定义，可表函子 $\underline{R}$ 与极限交换。
\end{remark}

\begin{remark}
\label{formal-defos-remark-formal-objects-yoneda}
设 $R$ 是 $\widehat{\mathcal{C}}_\Lambda$ 的对象。
它按说明 \ref{formal-defos-remarks-cofibered-groupoids}
(\ref{formal-defos-item-definition-yoneda}) 定义函子
$\underline{R}: \widehat{\mathcal{C}}_\Lambda \to \textit{Sets}$。
一如往常，我们把这个函子与相伴的集合余纤维范畴识别。
若 $\mathcal{F}$ 是 $\mathcal{C}_\Lambda$ 上的余纤维范畴，
则有范畴等价
\begin{equation}
\label{formal-defos-equation-formal-objects-maps}
\Mor_{\mathcal{C}_\Lambda}(
\underline{R}|_{\mathcal{C}_\Lambda}, \mathcal{F})
\longrightarrow
\widehat{\mathcal{F}}(R).
\end{equation}
它由复合
$$
\Mor_{\mathcal{C}_\Lambda}(
\underline{R}|_{\mathcal{C}_\Lambda}, \mathcal{F})
\xrightarrow{\Phi}
\Mor_{\widehat{\mathcal{C}}_\Lambda}(
\underline{R}, \widehat{\mathcal{F}})
\xrightarrow{\sim}
\widehat{\mathcal{F}}(R)
$$
给出，其中 $\Phi$ 如说明
\ref{formal-defos-remark-completion-restriction-adjoint} 所述，
而第二个等价来自 2-Yoneda 引理（即《范畴》引理
\ref{categories-lemma-yoneda-2category} 的余纤维版本）。
明确地说，这个等价把态射
$\varphi : \underline{R}|_{\mathcal{C}_\Lambda} \to \mathcal{F}$
映到 $\widehat{\mathcal{F}}(R)$ 中的形式对象
$(R, \varphi(R \to R/\mathfrak{m}_R^n), \varphi(f_n))$，
其中 $f_n : R/\mathfrak m_R^{n + 1} \to
R/\mathfrak m_R^n$ 是投影。

\medskip\noindent
假设已经为 $\mathcal{F}$ 选取了推前。
给定任意 $\xi \in \Ob(\widehat{\mathcal{F}}(R))$，
我们构造一个明确的态射
$\underline{\xi} :
\underline{R}|_{\mathcal{C}_\Lambda} \to \mathcal{F}$，
它在 (\ref{formal-defos-equation-formal-objects-maps}) 下映到 $\xi$。
具体地，设 $\xi = (R, \xi_n, f_n)$。
$\underline{R}|_{\mathcal{C}_\Lambda}$ 的对象 $\alpha$
等同于 $\widehat{\mathcal{C}}_\Lambda$ 中以阿廷环
$A$ 为靶的态射 $\alpha : R \to A$。
令 $m \in \mathbf{N}$ 为满足 $\mathfrak m_A^m = 0$ 的最小整数。
于是 $\alpha$ 唯一地经由
$\alpha_m : R/\mathfrak m_R^m \to A$ 分解，并可令
$\underline{\xi}(\alpha) = \alpha_{m, *}\xi_m$。
我们略去 $\underline{\xi}$ 在态射上的描述，也略去证明
$\underline{\xi}$ 经由
(\ref{formal-defos-equation-formal-objects-maps}) 映到 $\xi$ 的步骤。

\medskip\noindent
假设已经为 $\widehat{\mathcal{F}}$ 选取了推前。
在这种情况下，《范畴》引理
\ref{categories-lemma-yoneda-2category} 的证明给出 2-Yoneda
等价的一个明确拟逆
$$
\iota :
\widehat{\mathcal{F}}(R) \longrightarrow
\Mor_{\widehat{\mathcal{C}}_\Lambda}(
\underline{R}, \widehat{\mathcal{F}})
$$
它把 $\xi$ 映到态射
$\iota(\xi) : \underline{R} \to \widehat{\mathcal{F}}$；
该态射把
$f \in \underline{R}(S) =
\Mor_{\mathcal{C}_\Lambda}(R, S)$ 映到 $f_*\xi$。% P10-E0436
于是 (\ref{formal-defos-equation-formal-objects-maps}) 的一个拟逆为
$$
\widehat{\mathcal{F}}(R)
\xrightarrow{\iota}
\Mor_{\widehat{\mathcal{C}}_\Lambda}(
\underline{R}, \widehat{\mathcal{F}})
\xrightarrow{\Psi}
\Mor_{\mathcal{C}_\Lambda}(
\underline{R}|_{\mathcal{C}_\Lambda}, \mathcal{F})
$$
其中 $\Psi$ 如说明
\ref{formal-defos-remark-completion-restriction-adjoint} 所述。
给定 $\xi \in \Ob(\widehat{\mathcal{F}}(R))$，有
$\Psi(\iota(\xi)) \cong \underline{\xi}$，
其中 $\underline{\xi}$ 如上一段所述；这是因为在范畴等价
(\ref{formal-defos-equation-formal-objects-maps}) 下，两者都映到 $\xi$。
使用
$\underline{R} =
\widehat{\underline{R}|_{\mathcal{C}_\Lambda}}$
（见说明 \ref{formal-defos-remark-restrict-complete-continuous-functor}），
并展开 $\Phi$ 和 $\Psi$ 的定义，可得
$\iota(\xi)$ 与 $\underline{\xi}$ 的完备化同构。
\end{remark}

\begin{remark}
\label{formal-defos-remark-formal-objects-yoneda-map}
设 $\mathcal{F}$ 是 $\mathcal{C}_\Lambda$ 上的群胚余纤维范畴。
设 $\xi = (R, \xi_n, f_n)$ 和
$\eta = (S, \eta_n, g_n)$ 是 $\mathcal{F}$ 的形式对象。
令 $a = (a_n) : \xi \to \eta$ 为形式对象之间的态射，
即 $\widehat{\mathcal{F}}$ 的态射。
令 $f = \widehat{p}(a) = a_0 : R \to S$ 为 $a$ 在
$\widehat{\mathcal{C}}_\Lambda$ 中的投影。于是得到 $2$-交换图
$$
\xymatrix{
\underline{R}|_{\mathcal{C}_\Lambda} \ar[rd]_{\underline{\xi}} & &
\underline{S}|_{\mathcal{C}_\Lambda} \ar[ll]^f \ar[ld]^{\underline{\eta}} \\
& \mathcal{F}
}
$$
其中 $\underline{\xi}$ 和 $\underline{\eta}$ 是说明
\ref{formal-defos-remark-formal-objects-yoneda} 中构造的态射。
为说明这一点，设 $\alpha : S \to A$ 是
$\underline{S}|_{\mathcal{C}_\Lambda}$ 的对象（见上引文献）。
令 $m \in \mathbf{N}$ 为满足 $\mathfrak m_A^m = 0$ 的最小整数。
得到交换图
$$
\xymatrix{
R \ar[d]^f \ar[r] & R/\mathfrak m_R^m \ar[d]_{f_m} \ar[rd]^{\beta_m} \\
S \ar[r] & S/\mathfrak m_S^m \ar[r]^{\alpha_m} & A
}
$$
其底行箭头的复合为 $\alpha$。此时
$\underline{\eta}(\alpha) = \alpha_{m, *}\eta_m$ 且
$\underline{\xi}(\alpha \circ f) = \beta_{m, *}\xi_m$。
态射 $a_m : \xi_m \to \eta_m$ 位于 $f_m$ 上方，
故得到位于 $\text{id}_A$ 上方的典范态射
$$
\underline{\xi}(\alpha \circ f) = \beta_{m, *}\xi_m
\longrightarrow
\underline{\eta}(\alpha) = \alpha_{m, *}\eta_m
$$
使图
$$
\xymatrix{
\xi_m \ar[r] \ar[d]^{a_m} &  \beta_{m, *}\xi_m \ar[d] \\
\eta_m \ar[r] &  \alpha_{m, *}\eta_m
}
$$
由群胚余纤维范畴的公理而交换。这定义了函子变换
$\underline{\xi} \circ f \to \underline{\eta}$，
它见证了本说明第一个图的 2-交换性。
\end{remark}

\begin{remark}
\label{formal-defos-remark-spell-out-formal-object}
按照说明 \ref{formal-defos-remark-formal-objects-yoneda}，
给出 $\mathcal{F}$ 的形式对象 $\xi$，等价于给出一个
pro-可表函子 $U : \mathcal{C}_\Lambda \to \textit{Sets}$
以及一个态射 $U \to \mathcal{F}$。
\end{remark}

\section{光滑态射}
\label{formal-defos-section-smooth-morphisms}

\noindent
本节讨论 $\mathcal{C}_\Lambda$ 上群胚余纤维范畴之间的光滑态射。

\begin{definition}
\label{formal-defos-definition-smooth-morphism}
设 $\varphi : \mathcal{F} \to \mathcal{G}$ 是
$\mathcal{C}_\Lambda$ 上群胚余纤维范畴之间的态射。
若 $\varphi$ 满足下列条件，则称它{\it 光滑}：
设 $B \to A$ 是 $\mathcal{C}_\Lambda$ 中的满环映射，
$y \in \Ob(\mathcal{G}(B)), x \in
\Ob(\mathcal{F}(A))$，且
$y \to \varphi(x)$ 是位于 $B \to A$ 上方的态射。
则存在 $x' \in \Ob(\mathcal{F}(B))$、位于
$B \to A$ 上方的态射 $x' \to x$，以及位于
$\text{id}: B \to B$ 上方的态射 $\varphi(x') \to y$，
使图
$$
\xymatrix{
\varphi(x') \ar[r] \ar[dr] & y \ar[d] \\
& \varphi(x)
}
$$
交换。
\end{definition}

\begin{lemma}
\label{formal-defos-lemma-smoothness-small-extensions}
设 $\varphi : \mathcal{F} \to \mathcal{G}$ 是
$\mathcal{C}_\Lambda$ 上群胚余纤维范畴之间的态射。
若只对小扩张 $B \to A$ 假设定义
\ref{formal-defos-definition-smooth-morphism} 中的条件成立，
则 $\varphi$ 光滑。
\end{lemma}

\begin{proof}
设 $B \to A$ 是 $\mathcal{C}_\Lambda$ 中的满环映射。
设 $y \in \Ob(\mathcal{G}(B))$、
$x \in \Ob(\mathcal{F}(A))$，且
$y \to \varphi(x)$ 是位于 $B \to A$ 上方的态射。
由引理 \ref{formal-defos-lemma-factor-small-extension}，可把
$B \to A$ 分解为小扩张
$B = B_n \to B_{n-1} \to \ldots \to B_0 = A$。
对 $n$ 作归纳。若 $n = 1$，结论由假设成立。
若 $n > 1$，记
$f : B = B_n \to B_{n - 1}$ 及
$g : B_{n - 1} \to B_0 = A$。
选取 $y$ 沿 $f$ 的推前 $y \to f_* y$，
使态射 $y \to \varphi(x)$ 分解为
$y \to f_* y \to \varphi(x)$。
由归纳假设，可以找到位于
$g : B_{n - 1} \to A$ 上方的 $x_{n - 1} \to x$，
以及位于
$\text{id} : B_{n - 1} \to B_{n - 1}$ 上方的
$a : \varphi(x_{n - 1}) \to f_*y$，使图
$$
\xymatrix{
\varphi(x_{n - 1}) \ar[r]_-a \ar[dr] & f_*y \ar[d] \\
& \varphi(x)
}
$$
交换。把假设应用于 $y \to f_*y$ 与
$a^{-1} : f_*y \to \varphi(x_{n - 1})$ 的复合
$y \to \varphi(x_{n - 1})$，可得位于
$B_n \to B_{n - 1}$ 上方的
$x_n \to x_{n - 1}$，以及位于
$\text{id} : B_n \to B_n$ 上方的
$\varphi(x_n) \to y$，使图
$$
\xymatrix{
\varphi(x_n) \ar[r] \ar[d] & y \ar[d] \\
\varphi(x_{n - 1}) \ar[r]^-a \ar[dr] & f_*y \ar[d] \\
& \varphi(x)
}
$$
交换。于是复合 $x_n \to x_{n - 1} \to x$ 与
$\varphi(x_n) \to y$ 就是光滑性定义所需的态射。
\end{proof}

\begin{remark}
\label{formal-defos-remark-smoothness-2-categorical}
设 $\varphi : \mathcal{F} \to \mathcal{G}$ 是
$\mathcal{C}_\Lambda$ 上群胚余纤维范畴之间的态射，
且 $B \to A$ 是 $\mathcal{C}_\Lambda$ 中的环映射。
为纤维范畴 $\mathcal{F}(B)$ 和 $\mathcal{G}(B)$ 中的对象
选取沿 $B \to A$ 的推前，就确定了函子
$\mathcal{F}(B) \to \mathcal{F}(A)$ 和
$\mathcal{G}(B) \to \mathcal{G}(A)$，它们组成 $2$-交换图
$$
\xymatrix{
\mathcal{F}(B) \ar[r]^{\varphi} \ar[d] & \mathcal{G}(B) \ar[d] \\
\mathcal{F}(A) \ar[r]^{\varphi}        & \mathcal{G}(A) .
}
$$
因此诱导出函子
$\mathcal{F}(B) \to \mathcal{F}(A)
\times_{\mathcal{G}(A)} \mathcal{G}(B)$。
展开定义可知，$\varphi : \mathcal{F} \to \mathcal{G}$ 光滑，
当且仅当每当 $B \to A$ 为满射时，这个诱导函子本质满
（或者等价地，由引理
\ref{formal-defos-lemma-smoothness-small-extensions}，每当
$B \to A$ 为小扩张时它本质满）。
\end{remark}

\begin{remark}
\label{formal-defos-remark-compare-smooth-schlessinger}
说明 \ref{formal-defos-remark-smoothness-2-categorical} 中对光滑态射的刻画，
类似于 Schlessinger 的函子光滑态射概念；比较
\cite[Definition 2.2.]{Sch}。事实上，当
$\mathcal{F}$ 和 $\mathcal{G}$ 是集合余纤维范畴时，
我们的概念与 Schlessinger 的概念等价。
具体地，在这种情况下，设
$F, G : \mathcal{C}_\Lambda \to \textit{Sets}$ 为对应函子；
见说明 \ref{formal-defos-remarks-cofibered-groupoids}
(\ref{formal-defos-item-convention-cofibered-sets})。
则 $F \to G$ 光滑，当且仅当对
$\mathcal{C}_\Lambda$ 中的每个满环映射
$B \to A$，映射
$F(B) \to F(A) \times_{G(A)} G(B)$ 都满射。
\end{remark}

\begin{remark}
\label{formal-defos-remark-smooth-to-iso-classes}
设 $\mathcal{F}$ 是 $\mathcal{C}_\Lambda$ 上的群胚余纤维范畴。
则态射 $\mathcal{F} \to \overline{\mathcal{F}}$ 光滑。
事实上，设 $f : B \to A$ 是
$\mathcal{C}_\Lambda$ 中的环映射。
设 $x \in \Ob(\mathcal{F}(A))$，并设
$\overline{y} \in \overline{\mathcal{F}}(B)$
是某个 $y \in \Ob(\mathcal{F}(B))$ 的同构类，且满足
$\overline{f_*y} = \overline{x}$。
那么只需取 $x' = y$、隐含的位于 $B \to A$ 上方的态射
$x' = y \to x$，以及等式
$\overline{x'} = \overline{y}$，就解出了定义
\ref{formal-defos-definition-smooth-morphism} 中的问题。
\end{remark}

\noindent
若 $R \to S$ 是环映射
$\widehat{\mathcal{C}}_\Lambda$，则在函子
$\underline{S}, \underline{R}: \widehat{\mathcal{C}}_\Lambda
\to \textit{Sets}$ 之间诱导出态射
$\underline{S} \to \underline{R}$。% P10-E0437
在这种情况下，限制
$\underline{S}|_{\mathcal{C}_\Lambda} \to
\underline{R}|_{\mathcal{C}_\Lambda}$ 的光滑性是熟悉的概念：

\begin{lemma}
\label{formal-defos-lemma-smooth-morphism-power-series}
设 $R \to S$ 是 $\widehat{\mathcal{C}}_\Lambda$ 中的环映射。
则诱导态射
$\underline{S}|_{\mathcal{C}_\Lambda} \to
\underline{R}|_{\mathcal{C}_\Lambda}$
光滑，当且仅当 $S$ 是 $R$ 上的幂级数环。
\end{lemma}

\begin{proof}
先设 $S$ 是 $R$ 上的幂级数环，记
$S = R[[x_1, \ldots, x_n]]$。
$\underline{S}|_{\mathcal{C}_\Lambda} \to
\underline{R}|_{\mathcal{C}_\Lambda}$ 的光滑性意指如下性质
（见说明 \ref{formal-defos-remark-compare-smooth-schlessinger}）：
给定 $\mathcal{C}_\Lambda$ 中的满环映射 $B \to A$、
环映射 $R \to B$，以及使实线图
$$
\xymatrix{
S \ar[r] \ar@{..>}[rd] & A \\
R \ar[u] \ar[r] & B \ar[u]
}
$$
交换的环映射 $S \to A$，存在使该图交换的虚线箭头。
（注意这与《代数》定义
\ref{algebra-definition-formally-smooth} 的相似性。）
为构造虚线箭头，选取元素 $b_i \in B$，使其在 $A$ 中的像
等于 $x_i$ 在 $A$ 中的像。注意
$b_i \in \mathfrak m_B$，因为 $x_i$ 映到
$\mathfrak m_A$ 的元素。因此存在唯一的 $R$-代数映射
$R[[x_1, \ldots, x_n]] \to B$，它把 $x_i$ 映到 $b_i$，
并可作为所需虚线箭头。

\medskip\noindent
反过来，假设
$\underline{S}|_{\mathcal{C}_\Lambda} \to
\underline{R}|_{\mathcal{C}_\Lambda}$ 光滑。
选取 $x_1, \ldots, x_n \in S$，使其像构成 $S$ 在 $R$ 上的
相对余切空间
$\mathfrak m_S/(\mathfrak m_R S + \mathfrak m_S^2)$ 的一组基。
令 $T = R[[X_1, \ldots, X_n]]$。注意
$$
S/(\mathfrak m_R S + \mathfrak m_S^2) \cong
R/\mathfrak m_R[x_1, \ldots, x_n]/(x_ix_j)
$$
以及
$$
T/(\mathfrak m_R T + \mathfrak m_T^2) \cong
R/\mathfrak m_R[X_1, \ldots, X_n]/(X_iX_j).
$$
令
$S/(\mathfrak m_R S + \mathfrak m_S^2) \to
T/(\mathfrak m_R T + \mathfrak m_T^2)$
为把 $x_i$ 的类映到 $X_i$ 的类所定义的局部 $R$-代数同构。
令 $f_1 : S \to T/(\mathfrak m_R T + \mathfrak m_T^2)$
为复合
$S \to S/(\mathfrak m_R S + \mathfrak m_S^2)
\to T/(\mathfrak m_R T + \mathfrak m_T^2)$。
假设
$\underline{S}|_{\mathcal{C}_\Lambda} \to
\underline{R}|_{\mathcal{C}_\Lambda}$
光滑，意即可以把 $f_1$ 提升为映射
$f_2 : S \to T/\mathfrak{m}_T^2$，再提升为
$f_3 : S \to T/\mathfrak{m}_T^3$，并如此继续，对所有
$n \geq 1$ 都成立。于是得到局部 $R$-代数的诱导映射
$f : S \to T = \lim T/\mathfrak m_T^n$。
由 $f_1$ 的选取，映射 $f$ 在相对余切空间上诱导同构
$\mathfrak m_S/(\mathfrak m_R S + \mathfrak m_S^2) \to
\mathfrak m_T/(\mathfrak m_R T + \mathfrak m_T^2)$。
因此由引理 \ref{formal-defos-lemma-surjective-cotangent-space}，
$f$ 为满射（此处把 $f$ 视为
$\widehat{\mathcal{C}}_R$ 中的映射）。
选取 $X_i \in T$ 在 $f$ 下的原像 $y_i \in S$。
由于 $T$ 是 $R$ 上的幂级数环，存在局部 $R$-代数同态
$s : T \to S$，把 $X_i$ 映到 $y_i$。
由构造 $f \circ s = \text{id}$，故 $s$ 是单射。
但由于 $f$ 在相对余切空间上诱导同构，$s$ 也如此；
再次由引理 \ref{formal-defos-lemma-surjective-cotangent-space}，
它还是满射。因此 $s$ 和 $f$ 都是同构。
\end{proof}

\noindent
光滑态射具有如下函子性质。

\begin{lemma}
\label{formal-defos-lemma-smooth-properties}
设 $\varphi : \mathcal{F} \to \mathcal{G}$ 和
$\psi : \mathcal{G} \to \mathcal{H}$ 是
$\mathcal{C}_\Lambda$ 上群胚余纤维范畴之间的态射。
\begin{enumerate}
\item 若 $\varphi$ 和 $\psi$ 光滑，则
$\psi \circ \varphi$ 光滑。
\item 若 $\varphi$ 本质满且
$\psi \circ \varphi$ 光滑，则 $\psi$ 光滑。
\item 若 $\mathcal{G}' \to \mathcal{G}$ 是群胚余纤维范畴之间的态射，
且 $\varphi$ 光滑，则
$\mathcal{F} \times_\mathcal{G} \mathcal{G}' \to
\mathcal{G}'$ 光滑。
\end{enumerate}
\end{lemma}

\begin{proof}
(1) 和 (2) 直接由定义得出。(3) 的证明从略。提示：使用说明
\ref{formal-defos-remark-smoothness-2-categorical} 给出的光滑性表述，
并使用 $\mathcal{F} \times_\mathcal{G} \mathcal{G}'$
是 $2$-纤维积这一事实；见说明
\ref{formal-defos-remarks-cofibered-groupoids} (\ref{formal-defos-item-fibre-product})。
\end{proof}

\begin{lemma}
\label{formal-defos-lemma-smooth-morphism-essentially-surjective}
设 $\varphi : \mathcal{F} \to \mathcal{G}$ 是
$\mathcal{C}_\Lambda$ 上群胚余纤维范畴之间的光滑态射。
假设 $\varphi : \mathcal{F}(k) \to \mathcal{G}(k)$ 本质满。
则 $\varphi : \mathcal{F} \to \mathcal{G}$ 和
$\widehat{\varphi} : \widehat{\mathcal{F}} \to
\widehat{\mathcal{G}}$ 都本质满。
\end{lemma}

\begin{proof}
设 $y$ 是 $\mathcal{G}$ 中位于
$A \in \Ob(\mathcal{C}_\Lambda)$ 上方的对象。
设 $y \to y_0$ 是 $y$ 沿 $A \to k$ 的一个推前。
由 $\varphi : \mathcal{F}(k) \to \mathcal{G}(k)$
本质满这一假设，存在 $\mathcal{F}$ 中位于 $k$ 上方的对象
$x_0$ 以及同构 $y_0 \to \varphi(x_0)$。
$\varphi$ 的光滑性蕴含：存在 $\mathcal{F}$ 中位于
$A$ 上方的对象 $x$，使其像 $\varphi(x)$ 与 $y$ 同构。
因此 $\varphi : \mathcal{F} \to \mathcal{G}$ 本质满。

\medskip\noindent
设 $\eta = (R, \eta_n, g_n)$ 是
$\widehat{\mathcal{G}}$ 的对象。
我们构造 $\widehat{\mathcal{F}}$ 的对象 $\xi$ 以及同构
$\eta \to \varphi(\xi)$。由
$\varphi : \mathcal{F}(k) \to \mathcal{G}(k)$ 本质满这一假设，
对某个 $\xi_1 \in \Ob(\mathcal{F}(k))$，存在
$\mathcal{G}(k)$ 中的态射
$\eta_1 \to \varphi(\xi_1)$。态射
$\eta_2 \xrightarrow{g_1} \eta_1 \to \varphi(\xi_1)$
位于满环映射 $R/\mathfrak m_R^2 \to k$ 上方，
故由 $\varphi$ 的光滑性，存在
$\xi_2 \in \Ob(\mathcal{F}(R/\mathfrak m_R^2))$、
位于 $R/\mathfrak m_R^2 \to k$ 上方的态射
$f_1: \xi_2 \to \xi_1$，以及态射
$\eta_2 \to \varphi(\xi_2)$，使图
$$
\xymatrix{
\varphi(\xi_2)  \ar[r]^{\varphi(f_1)} &  \varphi(\xi_{1})   \\
\eta_2   \ar[u] \ar[r]^{g_1}  & \eta_1  \ar[u] \\
}
$$
交换。继续这一过程，就构造出
$\widehat{\mathcal{F}}$ 的对象
$\xi = (R, \xi_n, f_n)$，以及
$\widehat{\mathcal{G}}(R)$ 中的态射
$\eta \to \varphi(\xi) =
(R, \varphi(\xi_n), \varphi(f_n))$。
\end{proof}

\noindent
稍后我们要从 pro-可表函子构造到预形变范畴
$\mathcal{F}$ 的光滑态射。由说明
\ref{formal-defos-remark-formal-objects-yoneda} 中的讨论，
这些态射对应于 $\mathcal{F}$ 的某些形式对象。
更确切地说，它们就是所谓 $\mathcal{F}$ 的泛形式对象。

\begin{definition}
\label{formal-defos-definition-versal}
设 $\mathcal{F}$ 是群胚余纤维范畴，$\xi$ 是
$\mathcal{F}$ 的形式对象，位于
$R \in \Ob(\widehat{\mathcal{C}}_\Lambda)$ 上方。
若说明 \ref{formal-defos-remark-formal-objects-yoneda} 中的对应态射
$\underline{\xi}: \underline{R}|_{\mathcal{C}_\Lambda} \to
\mathcal{F}$ 光滑，则称 $\xi$ 是{\it 泛的}。
\end{definition}

\begin{remark}
\label{formal-defos-remark-versal-object}
设 $\mathcal{F}$ 是 $\mathcal C_\Lambda$ 上的群胚余纤维范畴，
且 $\xi$ 是 $\mathcal{F}$ 的形式对象。
由光滑性的定义，$\xi$ 的泛性等价于如下条件：
若
$$
\xymatrix{
& y \ar[d] \\
\xi \ar[r] & x
}
$$
是 $\widehat{\mathcal{F}}$ 中的图，且 $y \to x$
位于阿廷环之间的满射 $B \to A$ 上方
（可以假设它是小扩张），则存在态射 $\xi \to y$，使
$$
\xymatrix{
& y \ar[d] \\
\xi \ar[r] \ar[ur] & x
}
$$
交换。特别地，$\xi$ 为泛的这一条件不依赖于说明
\ref{formal-defos-remark-formal-objects-yoneda} 中构造
$\underline{\xi} :
\underline{R}|_{\mathcal{C}_\Lambda} \to \mathcal{F}$
时所作的推前选取。
\end{remark}

\begin{lemma}
\label{formal-defos-lemma-versal-object-quasi-initial}
设 $\mathcal{F}$ 是预形变范畴，且 $\xi$ 是
$\mathcal{F}$ 的泛形式对象。对
$\widehat{\mathcal{F}}$ 的任意形式对象 $\eta$，% P10-E0438
都存在态射 $\xi \to \eta$。
\end{lemma}

\begin{proof}
按假设，态射
$\underline{\xi} :
\underline{R}|_{\mathcal{C}_\Lambda} \to \mathcal{F}$
光滑。于是
$\iota(\xi) : \underline{R} \to \widehat{\mathcal{F}}$
是 $\underline{\xi}$ 的完备化；见说明
\ref{formal-defos-remark-formal-objects-yoneda}。
由引理 \ref{formal-defos-lemma-smooth-morphism-essentially-surjective}，
存在 $\underline{R}$ 的对象 $f$，使
$\iota(\xi)(f) = \eta$。此时 $f$ 是
$\widehat{\mathcal{C}}_\Lambda$ 中的环映射
$f : R \to S$。而 $\iota(\xi)(f) = \eta$ 意即
$f_*\xi \cong \eta$，这正是说存在位于 $f$ 上方的态射
$\xi \to \eta$。
\end{proof}

\section{光滑范畴或无阻碍范畴}
\label{formal-defos-section-smooth}

\noindent
设 $p : \mathcal{F} \to \mathcal{C}_\Lambda$ 是群胚余纤维范畴。
利用恒等函子，可以把 $\mathcal{C}_\Lambda$ 视为
$\mathcal{C}_\Lambda$ 上的群胚余纤维范畴。
这样，
$p : \mathcal{F} \longrightarrow \mathcal{C}_\Lambda$
就成为 $\mathcal{C}_\Lambda$ 上群胚余纤维范畴之间的态射。

\begin{definition}
\label{formal-defos-definition-cofibered-groupoid-projection-smooth}
设 $p : \mathcal{F} \to \mathcal{C}_\Lambda$ 是群胚余纤维范畴。
若其结构态射 $p$ 在定义 \ref{formal-defos-definition-smooth-morphism}
的意义下光滑，则称 $\mathcal{F}$ 是{\it 光滑的}或
{\it 无阻碍的}。
\end{definition}

\noindent
这是 $\mathcal{C}_\Lambda$ 上群胚余纤维范畴光滑性的“绝对”概念，
尽管更准确的说法是 $\mathcal{F}$ 在 $\Lambda$ 上光滑。
使用“{\it $\mathcal{F}$ 无阻碍}”这一说法时必须谨慎：
可能出现这种情况：$\mathcal{F}$ 有一个阻碍理论，
且该理论的阻碍空间非零，尽管 $\mathcal{F}$ 光滑。

\begin{remark}
\label{formal-defos-remark-smooth-on-top}
假设 $\mathcal{F}$ 是预形变范畴，并且有来自预形变范畴
$\mathcal{U}$ 的光滑态射
$\varphi : \mathcal U \to \mathcal{F}$。
则由引理 \ref{formal-defos-lemma-smooth-morphism-essentially-surjective}，
$\varphi$ 本质满；所以由引理 \ref{formal-defos-lemma-smooth-properties}，
$p: \mathcal{F} \to \mathcal{C}_\Lambda$ 光滑，
当且仅当复合
$\mathcal U \xrightarrow{\varphi} \mathcal{F} \xrightarrow{p}
\mathcal{C}_\Lambda$ 光滑；也就是说，$\mathcal{F}$ 光滑，
当且仅当 $\mathcal{U}$ 光滑。
\end{remark}

\begin{lemma}
\label{formal-defos-lemma-smooth}
设 $R \in \Ob(\widehat{\mathcal{C}}_\Lambda)$。下列条件等价：
\begin{enumerate}
\item $\underline{R}|_{\mathcal{C}_\Lambda}$ 光滑；
\item $\Lambda \to R$ 关于 $\mathfrak m_R$-进拓扑形式光滑；
\item $\Lambda \to R$ 平坦，并且
$R \otimes_\Lambda k'$ 在 $k'$ 上几何正则；
\item $\Lambda \to R$ 平坦，并且
$k' \to R \otimes_\Lambda k'$ 关于
$\mathfrak m_R$-进拓扑形式光滑。
\end{enumerate}
在经典情形下，这些条件还等价于
\begin{enumerate}
\item[(5)] 对某个 $n$，$R$ 同构于
$\Lambda[[x_1, \ldots, x_n]]$。
\end{enumerate}
\end{lemma}

\begin{proof}
态射
$p : \underline{R}|_{\mathcal{C}_\Lambda} \to
\mathcal{C}_\Lambda$
光滑，意即给定 $\mathcal{C}_\Lambda$ 中的满射
$B \to A$ 和映射 $R \to A$，可以在图
$$
\xymatrix{
R \ar[r] \ar@{-->}[rd] & A \\
\Lambda \ar[r] \ar[u] & B \ar[u]
}
$$
中找到虚线箭头。若 $\Lambda \to R$ 关于
$\mathfrak m_R$-进拓扑形式光滑，这当然成立；见
《代数进阶》定义
\ref{more-algebra-definition-formally-smooth-adic} 和
\ref{more-algebra-definition-formally-smooth}。
反过来，若上述性质成立，则由《代数进阶》引理
\ref{more-algebra-lemma-fs-local}，
$\Lambda \to R$ 关于 $\mathfrak m_R$-进拓扑形式光滑。
所以 (1) 与 (2) 等价。

\medskip\noindent
(2)、(3) 与 (4) 的等价性是《代数进阶》命题
\ref{more-algebra-proposition-fs-flat-fibre-fs}。
与 (5) 的等价性，例如可由引理
\ref{formal-defos-lemma-smooth-morphism-power-series} 以及经典情形下
$\mathcal{C}_\Lambda$ 与
$\underline{\Lambda}|_{\mathcal{C}_\Lambda}$ 相同这一事实得出。
\end{proof}

\begin{lemma}
\label{formal-defos-lemma-smooth-power-series-classical}
设 $\mathcal{F}$ 是预形变范畴，$\xi$ 是
$\mathcal{F}$ 的泛形式对象，位于
$R \in \Ob(\widehat{\mathcal{C}}_\Lambda)$ 上方。
下列条件等价：
\begin{enumerate}
\item $\mathcal{F}$ 无阻碍；
\item $\Lambda \to R$ 关于 $\mathfrak m_R$-进拓扑形式光滑。
\end{enumerate}
在经典情形下，这些条件还等价于
\begin{enumerate}
\item[(3)] 对某个 $n$，有
$R \cong \Lambda[[x_1, \ldots, x_n]]$。\ % P10-E0439
\end{enumerate}
\end{lemma}

\begin{proof}
若 (1) 成立，即 $\mathcal{F}$ 无阻碍，则复合
$$
\underline{R}|_{\mathcal{C}_\Lambda}
\xrightarrow{\underline{\xi}}
\mathcal{F}
\to
\mathcal{C}_\Lambda
$$
光滑；见引理 \ref{formal-defos-lemma-smooth-properties}。
所以由引理 \ref{formal-defos-lemma-smooth}，(2) 成立。
反过来，若 (2) 成立，则上述复合光滑，
而且由引理 \ref{formal-defos-lemma-versal-object-quasi-initial}，
第一个箭头本质满。因此由引理
\ref{formal-defos-lemma-smooth-properties}，第二个箭头光滑；
按定义，这意即 $\mathcal{F}$ 无阻碍。
经典情形下与 (3) 的等价性由引理 \ref{formal-defos-lemma-smooth} 得出。
\end{proof}

\begin{lemma}
\label{formal-defos-lemma-exists-smooth}
存在 $R \in \Ob(\widehat{\mathcal{C}}_\Lambda)$，
使引理 \ref{formal-defos-lemma-smooth} 的等价条件成立，并且还有
$H_1(L_{k/\Lambda}) = \mathfrak m_R/\mathfrak m_R^2$ 及
$\Omega_{R/\Lambda} \otimes_R k = \Omega_{k/\Lambda}$。
\end{lemma}

\begin{proof}
在经典情形下取 $R = \Lambda$。更一般地，若剩余域扩张
$k/k'$ 可分，则存在 $\Lambda$ 的完备化
$\Lambda^\wedge$ 的唯一有限平展扩张
$\Lambda^\wedge \to R$
（《代数》引理 \ref{algebra-lemma-complete-henselian} 和
\ref{algebra-lemma-henselian-cat-finite-etale}），
它在剩余域上诱导扩张 $k/k'$。

\medskip\noindent
一般情形下作如下处理。选取光滑 $\Lambda$-代数 $P$
和 $\Lambda$-代数满射 $P \to k$。
（例如，可令 $P$ 为多项式代数。）
以 $\mathfrak m_P$ 表示 $P \to k$ 的核。
Jacobi--Zariski 序列——见
(\ref{formal-defos-equation-sequence-extended}) 和《代数》引理
\ref{algebra-lemma-exact-sequence-NL}——是正合序列
$$
0 \to H_1(\NL_{k/\Lambda}) \to
\mathfrak m_P/\mathfrak m_P^2 \to
\Omega_{P/\Lambda} \otimes_P k \to
\Omega_{k/\Lambda} \to 0
$$
左端有 $0$，是因为 $P/k$ 光滑；% P10-E0440
所以 $\NL_{P/\Lambda}$ 拟同构于置于次数 $0$ 的有限投射模，
进而 $H_1(\NL_{P/\Lambda} \otimes_P k) = 0$。
假设 $f \in \mathfrak m_P$ 映到
$\Omega_{P/\Lambda} \otimes_P k$ 中的非零元素。
令 $P' = P/(f)$，便有 $\Lambda$-代数满射 $P' \to k$。
注意 $P'$ 在 $\mathfrak m_{P'}$ 处光滑：
这由《态射进阶》引理
\ref{more-morphisms-lemma-slice-smooth-given-element} 得出。
因此，用 $P'$ 的一个主局部化替换 $P$ 后，
$\dim(\mathfrak m_P/\mathfrak m_P^2)$ 减小。
有限次重复这一过程后，可以假设映射
$\mathfrak m_P/\mathfrak m_P^2 \to
\Omega_{P/\Lambda} \otimes_P k$ 为零，于是该正合序列分解为同构
$H_1(L_{k/\Lambda}) = \mathfrak m_P/\mathfrak m_P^2$ 和
$\Omega_{P/\Lambda} \otimes_P k = \Omega_{k/\Lambda}$。

\medskip\noindent
令 $R$ 为 $P$ 的 $\mathfrak m_P$-进完备化。
则 $R$ 是 $\widehat{\mathcal{C}}_\Lambda$ 的对象。
事实上，它是完备局部诺特环（见《代数》引理
\ref{algebra-lemma-completion-Noetherian-Noetherian}），
且其剩余域与 $k$ 识别。我们断言这个 $R$ 满足要求。

\medskip\noindent
首先注意，映射 $P \to R$ 诱导同构
$\mathfrak m_P/\mathfrak m_P^2 =
\mathfrak m_R/\mathfrak m_R^2$ 和
$\Omega_{P/\Lambda} \otimes_P k =
\Omega_{R/\Lambda} \otimes_R k$。
这是因为 $\mathfrak m_P/\mathfrak m_P^2$ 与
$\Omega_{P/\Lambda} \otimes_P k$ 都只依赖于
$\Lambda$-代数 $P/\mathfrak m_P^2$；见《代数》引理
\ref{algebra-lemma-differential-mod-power-ideal}。
对 $R$ 也同样如此，而且
$P/\mathfrak m_P^2 = R/\mathfrak m_R^2$。
利用 Jacobi--Zariski 序列
(\ref{formal-defos-equation-sequence-functorial}) 的函子性，
并利用同样结论对 $P$ 成立，可得
$H_1(L_{k/\Lambda}) = \mathfrak m_R/\mathfrak m_R^2$ 和
$\Omega_{R/\Lambda} \otimes_R k = \Omega_{k/\Lambda}$。

\medskip\noindent
最后，由于 $\Lambda \to P$ 光滑，故由《代数》命题
\ref{algebra-proposition-smooth-formally-smooth}，
$\Lambda \to P$ 形式光滑。
于是由《代数进阶》引理
\ref{more-algebra-lemma-formally-smooth}，
$\Lambda \to P$ 关于 $\mathfrak m_P$-进拓扑形式光滑。
由《代数进阶》引理
\ref{more-algebra-lemma-formally-smooth-completion}，
这一性质传递到完备化 $R$，证明完成。
事实上，只要
$\underline{R}|_{\mathcal{C}_\Lambda}$ 光滑，
$R$ 就同构于 $\Lambda$ 上某个光滑代数的完备化，
但我们不会用到这一点。
\end{proof}

\begin{example}
\label{formal-defos-example-smooth-explicit}
下面给出引理 \ref{formal-defos-lemma-exists-smooth} 中 $R$ 的一个更明确的例子。
设 $p$ 为素数，且 $n \in \mathbf{N}$。
令 $\Lambda = \mathbf{F}_p(t_1, t_2, \ldots, t_n)$，
$k = \mathbf{F}_p(x_1, \ldots, x_n)$，并令映射
$\Lambda \to k$ 由 $t_i \mapsto x_i^p$ 给出。此时可取
$$
R = \Lambda[x_1, \ldots, x_n]^\wedge_{(x_1^p - t_1, \ldots, x_n^p - t_n)}
$$
在这个例子中无法做得“更好”；也就是说，不能用
$\widehat{\mathcal{C}}_\Lambda$ 中更小的光滑对象来逼近
$\mathcal{C}_\Lambda$（可以论证 $R$ 的维数至少为 $n$，
因为映射
$\Omega_{R/\Lambda} \otimes_R k \to \Omega_{k/\Lambda}$
为满射）。稍后将更详细地讨论这一现象。
\end{example}

\section{Schlessinger 条件}
\label{formal-defos-section-schlessinger-conditions}

\noindent
下文经常考虑范畴 $\mathcal{C}_\Lambda$ 中环的纤维积
$A_1 \times_A A_2$。例 \ref{formal-defos-example-fibre-product}
已经表明，这样的纤维积未必总是 $\mathcal{C}_\Lambda$ 的对象。
不过，在下文几乎所有情形中，两个映射
$A_i \to A$ 之一都是满射；由引理
\ref{formal-defos-lemma-fiber-product-CLambda}，
$A_1 \times_A A_2$ 将是 $\mathcal{C}_\Lambda$ 的对象。
我们将不再说明而直接使用这一结果。

\medskip\noindent
以 $k[\epsilon]$ 表示 $k$ 上的对偶数环。
更一般地，对 $k$-向量空间 $V$，以 $k[V]$ 表示如下
$k$-代数：其底层向量空间为 $k \oplus V$，乘法由
$(a, v) \cdot (a', v') = (aa', av' + a'v)$ 给出。
当 $V = k$ 时，$k[V]$ 就是 $k$ 上的对偶数环。
对任意有限维 $k$-向量空间 $V$，环 $k[V]$ 都属于
$\mathcal{C}_\Lambda$。

\begin{definition}
\label{formal-defos-definition-S1-S2}
设 $\mathcal{F}$ 是 $\mathcal C_\Lambda$ 上的群胚余纤维范畴。
如下定义 $\mathcal{F}$ 上的{\it 条件 (S1) 和 (S2)}：
\begin{enumerate}
\item[(S1)] $\mathcal{F}$ 中每个图
$$
\vcenter{
\xymatrix{
           & x_2 \ar[d] \\
x_1 \ar[r] & x
}
}
\quad\text{位于下图上方}\quad
\vcenter{
\xymatrix{
           & A_2 \ar[d] \\
A_1 \ar[r] & A
}
}
$$
若其底图属于 $\mathcal{C}_\Lambda$ 且 $A_2 \to A$ 满射，
则可补成交换图
$$
\vcenter{
\xymatrix{
y \ar[r] \ar[d] & x_2 \ar[d] \\
x_1 \ar[r]      & x
}
}
\quad\text{位于下图上方}\quad
\vcenter{
\xymatrix{
A_1 \times_A A_2 \ar[r] \ar[d] & A_2 \ar[d] \\
A_1 \ar[r]      & A.
}
}
$$
\item[(S2)]
对 $\mathcal{F}$ 中位于 $\mathcal{C}_\Lambda$ 中如下形式图上方的图，
(S1) 条件成立：
$$
\xymatrix{
          & k[\epsilon] \ar[d] \\
A  \ar[r] & k.
}
$$
此外，若 $\mathcal{F}$ 中有两个交换图
$$
\vcenter{
\xymatrix{
y \ar[r]_c \ar[d]_a & x_\epsilon \ar[d]^e \\
x \ar[r]^d          & x_0
}
}
\quad\text{和}\quad
\vcenter{
\xymatrix{
y' \ar[r]_{c'} \ar[d]_{a'} & x_\epsilon \ar[d]^e \\
x \ar[r]^d                 & x_0
}
}
\quad\text{位于下图上方}\quad
\vcenter{
\xymatrix{
A \times_k k[\epsilon] \ar[r] \ar[d] & k[\epsilon] \ar[d] \\
A  \ar[r] & k
}
}
$$
则存在
$\mathcal{F}(A \times_k k[\epsilon])$ 中的态射
$b : y \to y'$，使 $a = a' \circ b$。
\end{enumerate}
\end{definition}

\noindent
可以用纤维范畴部分地解释条件 (S1) 与 (S2) 的含义。
设 $f_1 : A_1 \to A$ 和 $f_2 : A_2 \to A$ 是
$\mathcal{C}_\Lambda$ 中的环映射，且 $f_2$ 满射。
以 $p_i : A_1 \times_A A_2 \to A_i$ 表示投影映射。
假设已经为 $\mathcal{F}$ 作了推前选取。
于是环的交换图转化为纤维范畴的 $2$-交换图
$$
\xymatrix{
\mathcal{F}(A_1 \times_A A_2) \ar[r]_-{p_{2, *}} \ar[d]_{p_{1, *}} &
\mathcal{F}(A_2) \ar[d]^{f_{2, *}} \\
\mathcal{F}(A_1) \ar[r]^{f_{1, *}} & \mathcal{F}(A)
}
$$
从而得到到范畴 $2$-纤维积的函子
\begin{equation}
\label{formal-defos-equation-compare}
\mathcal{F}(A_1 \times_A A_2) \to
\mathcal{F}(A_1) \times_{\mathcal{F}(A)} \mathcal{F}(A_2)
\end{equation}
条件 (S1) 要求这个函子本质满。
条件 (S2) 的第一部分要求：若 $f_2$ 等于映射
$k[\epsilon] \to k$，这个函子是本质满的。% P10-E0441
此外，在这种情况下，(S2) 的第二部分蕴含：
在靶中变为同构的两个对象在源中同构
（但它与这一陈述{\it 并不}等价）。
按照定义中的方式陈述这些条件的优点，是无需作任何选取。

\begin{lemma}
\label{formal-defos-lemma-S1-small-extensions}
设 $\mathcal{F}$ 是 $\mathcal C_\Lambda$ 上的群胚余纤维范畴。
若只在 $A_2 \to A$ 为小扩张时假设 (S1) 的条件成立，
则 $\mathcal{F}$ 满足 (S1)。
\end{lemma}

\begin{proof}
证明从略。提示：应用引理
\ref{formal-defos-lemma-factor-small-extension}，并作类似于引理
\ref{formal-defos-lemma-smoothness-small-extensions} 证明中的归纳。
\end{proof}

\begin{remark}
\label{formal-defos-remark-compare-S1-S2-schlessinger}
当 $\mathcal{F}$ 是集合余纤维范畴时，条件 (S1) 与 (S2)
正是 Schlessinger 论文 \cite{Sch} 中的条件 (H1) 与 (H2)。
具体地，对函子
$F: \mathcal{C}_\Lambda \to \textit{Sets}$，
条件 (S1) 与 (S2) 表述为：
\begin{enumerate}
\item [(S1)] 若 $A_1 \to A$ 和 $A_2 \to A$ 是
$\mathcal{C}_\Lambda$ 中的映射，且 $A_2 \to A$ 满射，
则诱导映射
$F(A_1 \times_A A_2) \to F(A_1) \times_{F(A)} F(A_2)$
为满射。
\item [(S2)] 若 $A \to k$ 是 $\mathcal{C}_\Lambda$ 中的映射，
则诱导映射
$F(A \times_k k[\epsilon]) \to F(A) \times_{F(k)} F(k[\epsilon])$
为双射。
\end{enumerate}
映射
$F(A \times_k k[\epsilon]) \to F(A) \times_{F(k)} F(k[\epsilon])$
的单射性来自条件 (S2) 的第二部分以及所有态射都是恒等态射这一事实。
\end{remark}

\begin{lemma}
\label{formal-defos-lemma-S2-extensions}
设 $\mathcal{F}$ 是 $\mathcal{C}_\Lambda$ 上的群胚余纤维范畴。
若 $\mathcal{F}$ 满足 (S2)，则对任意有限维
$k$-向量空间 $V$，把 $k[\epsilon]$ 换成 $k[V]$ 后，
(S2) 的条件仍成立。
\end{lemma}

\begin{proof}
若 $\mathcal{F}$ 是集合余纤维范畴，即对应于函子
$F : \mathcal{C}_\Lambda \to \textit{Sets}$，
则结论来自说明 \ref{formal-defos-remark-compare-S1-S2-schlessinger}
中对 $F$ 的 (S2) 的描述，以及同构
$k[V] \cong k[\epsilon] \times_k \ldots \times_k k[\epsilon]$；
右边有 $\dim_k V$ 个因子。
本章余下部分只会用到函子情形。

\medskip\noindent
一般情形下，对 $\dim(V)$ 作归纳。
若 $\dim(V) = 1$，则 $k[V] \cong k[\epsilon]$，
结论由假设成立。若 $\dim(V) > 1$，写成
$V = V' \oplus k\epsilon$。选取图
$$
\vcenter{
\xymatrix{
& x_V \ar[d] \\
x \ar[r] & x_0
}
}
\quad\text{位于下图上方}\quad
\vcenter{
\xymatrix{
& k[V] \ar[d] \\
A \ar[r] & k
}
}
$$
选取位于 $k[V] \to k[V']$ 上方的态射
$x_V \to x_{V'}$，以及位于
$k[V] \to k[\epsilon]$ 上方的态射
$x_V \to x_\epsilon$。
注意，态射 $x_V \to x_0$ 分别分解为
$x_V \to x_{V'} \to x_0$ 和
$x_V \to x_\epsilon \to x_0$。
由归纳假设，可以找到图
$$
\vcenter{
\xymatrix{
y' \ar[d] \ar[r] & x_{V'} \ar[d] \\
x \ar[r] & x_0
}
}
\quad\text{位于下图上方}\quad
\vcenter{
\xymatrix{
A \times_k k[V'] \ar[d] \ar[r] & k[V'] \ar[d] \\
A \ar[r] & k
}
}
$$
这给出交换图
$$
\vcenter{
\xymatrix{
& x_\epsilon \ar[d] \\
y' \ar[r] & x_0
}
}
\quad\text{位于下图上方}\quad
\vcenter{
\xymatrix{
& k[\epsilon] \ar[d] \\
A \times_k k[V'] \ar[r] & k
}
}
$$
因此，由 (S2) 得到交换图
$$
\vcenter{
\xymatrix{
y \ar[d] \ar[r] & x_\epsilon \ar[d] \\
y' \ar[r] & x_0
}
}
\quad\text{位于下图上方}\quad
\vcenter{
\xymatrix{
(A \times_k k[V']) \times_k k[\epsilon] \ar[d] \ar[r] &
k[\epsilon] \ar[d] \\
A \times_k k[V'] \ar[r] & k
}
}
$$
注意
$(A \times_k k[V']) \times_k k[\epsilon] =
A \times_k k[V' \oplus k\epsilon]
= A \times_k k[V]$。
我们断言 $y$ 可置于正确的交换图中。
为说明这一点，取位于
$A \times_k k[V] \to k[V]$ 上方的态射
$y \to y_V$。可以把态射
$y \to y' \to x_{V'}$ 和 $y \to x_\epsilon$
经由态射 $y \to y_V$ 分解
（由群胚余纤维范畴的公理）。
因此，$y_V$ 与 $x_V$ 都可置于交换图
$$
\vcenter{
\xymatrix{
y_V \ar[r] \ar[d] & x_\epsilon \ar[d] \\
x_{V'} \ar[r]          & x_0
}
}
\quad\text{和}\quad
\vcenter{
\xymatrix{
x_V \ar[r] \ar[d] & x_\epsilon \ar[d] \\
x_{V'} \ar[r]                 & x_0
}
}
$$
中。因此，由 (S2) 的第二部分，存在同构
$y_V \to x_V$，它与 $y_V \to x_{V'}$ 和
$x_V \to x_{V'}$ 相容，特别地也与到 $x_0$ 的映射相容。
于是复合 $y \to y_V \to x_V$ 可置于所需交换图
$$
\vcenter{
\xymatrix{
y \ar[r] \ar[d] & x_V \ar[d] \\
x \ar[r] & x_0
}
}
\quad\text{位于下图上方}\quad
\vcenter{
\xymatrix{
A \times_k k[V] \ar[d] \ar[r] & k[V] \ar[d] \\
A \ar[r] & k
}
}
$$
中。由此可见，把 $k[\epsilon]$ 换成 $k[V]$ 后，
$(S2)$ 的第一部分成立。

\medskip\noindent
为证明第二部分，设给定两个交换图
$$
\vcenter{
\xymatrix{
y \ar[r] \ar[d] & x_V \ar[d] \\
x \ar[r] & x_0
}
}
\quad\text{和}\quad
\vcenter{
\xymatrix{
y' \ar[r] \ar[d] & x_V \ar[d] \\
x \ar[r] & x_0
}
}
\quad\text{位于下图上方}\quad
\vcenter{
\xymatrix{
A \times_k k[V] \ar[d] \ar[r] & k[V] \ar[d] \\
A \ar[r] & k
}
}
$$
我们将使用本证明第一段所引入的态射
$x_V \to x_{V'} \to x_0$ 和
$x_V \to x_\epsilon \to x_0$。
选取位于
$A \times_k k[V] \to A \times_k k[V']$ 上方的态射
$y \to y_{V'}$ 和 $y' \to y'_{V'}$。
余纤维范畴的公理蕴含可以找到交换图
$$
\vcenter{
\xymatrix{
y_{V'} \ar[r] \ar[d] & x_{V'} \ar[d] \\
x \ar[r] & x_0
}
}
\quad\text{和}\quad
\vcenter{
\xymatrix{
y'_{V'} \ar[r] \ar[d] & x_{V'} \ar[d] \\
x \ar[r] & x_0
}
}
\quad\text{位于下图上方}\quad
\vcenter{
\xymatrix{
A \times_k k[V'] \ar[d] \ar[r] & k[V'] \ar[d] \\
A \ar[r] & k
}
}
$$
由归纳假设，得到同构
$b : y_{V'} \to y'_{V'}$，
它与态射 $y_{V'} \to x$ 和
$y'_{V'} \to x$ 相容，特别地与到 $x_0$ 的态射相容。
于是有交换图
$$
\vcenter{
\xymatrix{
y \ar[r] \ar[d] & x_\epsilon \ar[d] \\
y'_{V'} \ar[r] & x_0
}
}
\quad\text{和}\quad
\vcenter{
\xymatrix{
y' \ar[r] \ar[d] & x_\epsilon \ar[d] \\
y'_{V'} \ar[r] & x_0
}
}
\quad\text{位于下图上方}\quad
\vcenter{
\xymatrix{
A \times_k k[\epsilon] \ar[d] \ar[r] & k[\epsilon] \ar[d] \\
A \ar[r] & k
}
}
$$
其中态射 $y \to y'_{V'}$ 是复合
$y \to y_{V'} \xrightarrow{b} y'_{V'}$，
而态射 $y \to x_\epsilon$ 和
$y' \to x_\epsilon$ 分别是映射
$y \to x_V$ 和 $y' \to x_V$ 与态射
$x_V \to x_\epsilon$ 的复合。
于是 (S2) 的第二部分保证存在同构
$y \to y'$，它与到 $y'_{V'}$ 的映射相容，
特别地与到 $x$ 的映射相容
（因为 $b$ 与到 $x$ 的映射相容）。
\end{proof}

\begin{lemma}
\label{formal-defos-lemma-S1-S2-associated-functor}
设 $\mathcal{F}$ 是 $\mathcal{C}_\Lambda$ 上的群胚余纤维范畴。
\begin{enumerate}
\item 若 $\mathcal{F}$ 满足 (S1)，则
$\overline{\mathcal{F}}$ 也满足 (S1)。
\item 若 $\mathcal{F}$ 满足 (S2)，并且下列条件至少有一个成立，
则 $\overline{\mathcal{F}}$ 也满足 (S2)：
\begin{enumerate}
\item $\mathcal{F}$ 是预形变范畴；
\item 范畴 $\mathcal{F}(k)$ 是集合或集合型群胚；
\item 对 $\mathcal{F}$ 中位于
$k[\epsilon] \to k$ 上方的任意态射
$x_\epsilon \to x_0$，推前映射
$\text{Aut}_{k[\epsilon]}(x_\epsilon) \to
\text{Aut}_k(x_0)$ 为满射。
\end{enumerate}
\end{enumerate}
\end{lemma}

\begin{proof}
假设 $\mathcal{F}$ 满足 (S1)。
设 $f_i : A_i \to A$ 是 $\mathcal{C}_\Lambda$ 中的环映射，
且 $f_2$ 满射。设 $x_i \in \mathcal{F}(A_i)$，
并且推前 $f_{1, *}(x_1)$ 与 $f_{2, *}(x_2)$ 同构。
于是可以用 $x$ 表示 $\mathcal{F}$ 中位于 $A$ 上方、
且与这两个推前都同构的对象，% P10-E0442
并得到 (S1) 中的图。因此，找到 $\mathcal{F}$ 中位于
$A_1 \times_A A_2$ 上方的对象 $y$，
其到 $A_1$、相应地到 $A_2$ 的推前分别与
$x_1$、相应地与 $x_2$ 同构。
由此可见，(S1) 对 $\overline{\mathcal{F}}$ 成立。

\medskip\noindent
假设 $\mathcal{F}$ 满足 (S2)。
$\overline{\mathcal{F}}$ 的 (S2) 第一部分由与上述相同的论证得出。
$\overline{\mathcal{F}}$ 的 (S2) 第二部分意指：对
$\mathcal{C}_\Lambda$ 中的任意环 $A$，映射
$$
\overline{\mathcal{F}}(A \times_k k[\epsilon]) \to
\overline{\mathcal{F}}(A)
\times_{\overline{\mathcal{F}}(k)}
\overline{\mathcal{F}}(k[\epsilon])
$$
为单射。设
$y, y' \in \mathcal{F}(A \times_k k[\epsilon])$。
利用余纤维范畴的公理，可以选取交换图
$$
\vcenter{
\xymatrix{
y \ar[r]_c \ar[d]_a & x_\epsilon \ar[d]^e \\
x \ar[r]^d          & x_0
}
}
\quad\text{和}\quad
\vcenter{
\xymatrix{
y' \ar[r]_{c'} \ar[d]_{a'} & x'_\epsilon \ar[d]^{e'} \\
x' \ar[r]^{d'}                 & x'_0
}
}
\quad\text{位于下图上方}\quad
\vcenter{
\xymatrix{
A \times_k k[\epsilon] \ar[d] \ar[r] & k[\epsilon] \ar[d] \\
A \ar[r] & k
}
}
$$
假设存在 $\mathcal{F}(A)$ 中的同构
$\alpha : x \to x'$，以及
$\mathcal{F}(k[\epsilon])$ 中的同构
$\beta : x_\epsilon \to x'_\epsilon$。
这也意味着存在与 $\alpha$ 相容的同构
$\gamma : x_0 \to x'_0$。
为证明 $\overline{\mathcal{F}}$ 的 (S2)，必须证明
$\mathcal{F}(A \times_k k[\epsilon])$ 中存在同构
$y \to y'$。
由 $\mathcal{F}$ 的 (S2)，如果能选取同构
$\alpha$、$\beta$ 和 $\gamma$，使
$$
\xymatrix{
x \ar[d]^\alpha \ar[r] & x_0 \ar[d]^\gamma &
x_\epsilon \ar[d]^\beta \ar[l]^e \\
x' \ar[r] & x'_0 & x'_\epsilon \ar[l]_{e'}
}
$$
交换，这样的态射就存在
（因为这时可在前一个显示图中用 $x'$ 替换 $x$，
并用 $x'_\epsilon$ 替换 $x_\epsilon$）。
左侧方块由 $\gamma$ 的选取而交换。
可把 $e' \circ \beta$ 分解成
$\gamma' \circ e$，其中
$\gamma' : x_0 \to x'_0$ 是另一个映射。
现在的问题是能否安排 $\gamma = \gamma'$？
若 $\mathcal{F}(k)$ 是集合或集合型群胚，这显然成立。
此外，若
$\text{Aut}_{k[\epsilon]}(x_\epsilon) \to
\text{Aut}_k(x_0)$ 为满射，
则可在 $\beta$ 前复合一个 $x_\epsilon$ 的自同构，
其像为 $\gamma^{-1} \circ \gamma'$，从而使要求成立。
\end{proof}

\begin{lemma}
\label{formal-defos-lemma-S1-S2-localize}
设 $\mathcal{F}$ 是 $\mathcal{C}_\Lambda$ 上的群胚余纤维范畴，
且 $x_0 \in \Ob(\mathcal{F}(k))$。
设 $\mathcal{F}_{x_0}$ 是说明
\ref{formal-defos-remark-localize-cofibered-groupoid} 中构造的
$\mathcal{C}_\Lambda$ 上群胚余纤维范畴。
\begin{enumerate}
\item 若 $\mathcal{F}$ 满足 (S1)，则
$\mathcal{F}_{x_0}$ 也满足 (S1)。
\item 若 $\mathcal{F}$ 满足 (S2)，则
$\mathcal{F}_{x_0}$ 也满足 (S2)。
\end{enumerate}
\end{lemma}

\begin{proof}
$\mathcal{F}_{x_0}$ 中每个如定义
\ref{formal-defos-definition-S1-S2} 所述的图，都产生 $\mathcal{F}$ 中的图；
而条件 (S1) 或 (S2) 对 $\mathcal{F}$ 中该图给出的输出，
也可视为对 $\mathcal{F}_{x_0}$ 的输出。
\end{proof}

\begin{lemma}
\label{formal-defos-lemma-lifting-section}
设 $p: \mathcal{F} \to \mathcal{C}_\Lambda$ 是群胚余纤维范畴。
考虑 $\mathcal{F}$ 中的图
$$
\vcenter{
\xymatrix{
y \ar[r] \ar[d]_a & x_\epsilon \ar[d]_e \\
x \ar[r]^d        & x_0
}
}
\quad\text{位于下图上方}\quad
\vcenter{
\xymatrix{
A \times_k k[\epsilon] \ar[r] \ar[d] & k[\epsilon] \ar[d] \\
A \ar[r] & k.
}
}
$$
其中底图属于 $\mathcal{C}_\Lambda$。
假设 $\mathcal{F}$ 满足 (S2)。
则存在满足 $a \circ s = \text{id}_x$ 的态射
$s : x \to y$，当且仅当存在满足
$e \circ s_\epsilon = d$ 的态射
$s_\epsilon : x \to x_\epsilon$。
\end{lemma}

\begin{proof}
“仅当”方向显然。反过来，假设存在态射
$s_\epsilon : x \to x_\epsilon$，使
$e \circ s_\epsilon = d$。
注意，$p(s_\epsilon) : A \to k[\epsilon]$ 是与映射
$A \to k$ 相容的环映射。因此得到
$$
\sigma = (\text{id}_A, p(s_\epsilon)) :
A \to A \times_k k[\epsilon].
$$
选取推前 $x \to \sigma_*x$。
由构造，可把 $s_\epsilon$ 分解为
$x \to \sigma_*x \to x_\epsilon$。
此外，由于 $\sigma$ 是
$A \times_k k[\epsilon] \to A$ 的截面，
得到态射 $\sigma_*x \to x$，使
$x \to \sigma_*x \to x$ 为 $\text{id}_x$。
因为 $e \circ s_\epsilon = d$，图
$$
\xymatrix{
\sigma_*x \ar[r] \ar[d] & x_\epsilon \ar[d]_e \\
x \ar[r]^d        & x_0
}
$$
交换。因此，由 (S2) 得到态射
$\sigma_*x \to y$，使复合
$\sigma_*x \to y \to x$ 等于给定映射
$\sigma_*x \to x$。现在问题的解是取
$a : x \to y$ 为复合
$x \to \sigma_*x \to y$。% P10-E0443
\end{proof}

\begin{lemma}
\label{formal-defos-lemma-lifting-along-small-extension}
考虑预形变范畴 $\mathcal{F}$ 中的交换图
$$
\vcenter{
\xymatrix{
y \ar[r] \ar[d] & x_2 \ar[d]^{a_2} \\
x_1 \ar[r]^{a_1}        & x
}
}
\quad\text{位于下图上方}
\vcenter{
\xymatrix{
A_1 \times_A A_2 \ar[r] \ar[d] & A_2 \ar[d]^{f_2} \\
A_1 \ar[r]^{f_1} & A
}
}
$$
上方，其中底图属于 $\mathcal{C}_\Lambda$，且
$f_2 : A_2 \to A$ 是小扩张。
假设存在映射 $h : A_1 \to A_2$，使
$f_2 = f_1 \circ h$。% P10-E0444
令 $I = \Ker(f_2)$。考虑环映射
$$
g : A_1 \times_A A_2 \longrightarrow k[I] = k \oplus I, \quad
(u, v) \longmapsto \overline{u} \oplus (v - h(u))
$$
选取推前 $y \to g_*y$。假设 $\mathcal{F}$ 满足 (S2)。
若存在态射 $x_1 \to g_*y$，则存在态射
$b: x_1 \to x_2$，使 $a_1 =  a_2 \circ b$。
\end{lemma}

\begin{proof}
注意，
$\text{id}_{A_1} \times g :
A_1 \times_A A_2 \to A_1 \times_k k[I]$
是同构，且 $k[I] \cong k[\epsilon]$。因此有图
$$
\vcenter{
\xymatrix{
y \ar[r] \ar[d] & g_*y \ar[d] \\
x_1 \ar[r]        & x_0
}
}
\quad\text{位于下图上方}\quad
\vcenter{
\xymatrix{
A_1 \times_k k[\epsilon] \ar[r] \ar[d] & k[\epsilon] \ar[d] \\
A_1 \ar[r] & k.
}
}
$$
其中 $x_0$ 是 $\mathcal{F}$ 中位于 $k$ 上方的对象
（$\mathcal{F}$ 的每个对象都有到 $x_0$ 的唯一态射；
见定义 \ref{formal-defos-definition-predeformation-category} 之后的讨论）。
若有态射 $x_1 \to g_*y$，则引理
\ref{formal-defos-lemma-lifting-section} 给出映射
$y \to x_1$ 的截面 $s : x_1 \to y$。
将它与映射 $y \to x_2$ 复合，得到
$b : x_1 \to x_2$；它满足
$a_1 =  a_2 \circ b$，因为引理中的图交换且 $s$ 是截面。
\end{proof}





\section{函子的切空间}
\label{formal-defos-section-tangent-spaces-functors}

\noindent
设 $R$ 是环。用 $\text{Mod}_R$ 表示 $R$-模范畴，
用 $\text{Mod}^{fg}_R$ 表示有限生成 $R$-模范畴。

\begin{definition}
\label{formal-defos-definition-linear}
设 $L: \text{Mod}^{fg}_R \to \text{Mod}_R$，
相应地，设 $L: \text{Mod}_R \to \text{Mod}_R$ 是函子。
若对 $\text{Mod}^{fg}_R$、相应地 $\text{Mod}_R$ 中的每一对对象
$M,N$，映射
$$
L : \Hom_R(M, N) \longrightarrow \Hom_R(L(M), L(N))
$$
都是 $R$-模映射，则称 $L$ 是 {\it $R$-线性的}。
\end{definition}

\begin{remark}
\label{formal-defos-remark-linear-enriched-over-modules}
对任意两个在 $\text{Mod}_R$ 上富集的范畴之间的函子，
都可以定义 $R$-线性的概念。这里专门对函子
$L: \text{Mod}^{fg}_R \to \text{Mod}_R$ 和
$L: \text{Mod}_R \to \text{Mod}_R$
作出定义，是因为迄今所需的正是这两种情形。
\end{remark}

\begin{remark}
\label{formal-defos-remark-linear-functor}
若 $L: \text{Mod}^{fg}_R \to \text{Mod}_R$ 是 $R$-线性函子，
则 $L$ 保持有限积，并把零模送到零模；见
《同调》引理 \ref{homology-lemma-additive-additive}。
另一方面，若函子 $\text{Mod}^{fg}_R \to \textit{Sets}$
保持有限积，并把零模送到单元素集合，
则它可唯一地提升为 $R$-线性函子；见引理
\ref{formal-defos-lemma-linear-functor}。
\end{remark}

\begin{lemma}
\label{formal-defos-lemma-linear-functor}
设 $L: \text{Mod}^{fg}_R \to \textit{Sets}$，
相应地，设 $L: \text{Mod}_R \to \textit{Sets}$ 是函子。
假设 $L(0)$ 是单元素集合，并且 $L$ 保持有限积。
则存在唯一的 $R$-线性函子
$\widetilde{L} : \text{Mod}^{fg}_R \to \text{Mod}_R$，
相应地，存在唯一的
$\widetilde{L} : \text{Mod}_R \to \text{Mod}_R$，使图
$$
\vcenter{
\xymatrix{
& \text{Mod}_R \ar[dr]^{\text{遗忘}} &   \\
\text{Mod}^{fg}_R  \ar[ur]^{\widetilde{L}} \ar[rr]^{L} &  &
\textit{Sets}
}
}
\quad\text{相应地}\quad
\vcenter{
\xymatrix{
& \text{Mod}_R \ar[dr]^{\text{遗忘}} &   \\
\text{Mod}_R  \ar[ur]^{\widetilde{L}} \ar[rr]^{L} &  &
\textit{Sets}
}
}
$$
交换。
\end{lemma}

\begin{proof}
这里只证明 $L: \text{Mod}^{fg}_R \to \textit{Sets}$ 的情形。
设 $M$ 是有限生成 $R$-模。我们把 $\widetilde{L}(M)$
定义为集合 $L(M)$ 配备以下 $R$-模结构。

\medskip\noindent
乘法：若 $r \in R$，则 $L(M)$ 上的乘以 $r$ 定义为由乘法映射
$r \cdot : M \to M$ 诱导的映射 $L(M) \to L(M)$。

\medskip\noindent
加法：求和映射 $M \times M \to M: (m_1, m_2) \mapsto m_1 + m_2$
诱导映射 $L(M \times M) \to L(M)$。由假设，$L(M \times M)$
典范同构于 $L(M) \times L(M)$。于是 $L(M)$ 上的加法定义为映射
$L(M) \times L(M) \cong L(M \times M) \to L(M)$。

\medskip\noindent
零元：存在唯一映射 $0 \to M$。$L(M)$ 的零元是映射
$L(0) \to L(M)$ 的像。

\medskip\noindent
我们略去验证：这定义了 $R$-模 $\widetilde{L}(M)$，
并且这是唯一使其关于 $M$ 具有 $R$-线性函子性的结构。% P10-E0445
\end{proof}

\begin{lemma}
\label{formal-defos-lemma-morphism-linear-functors}
设 $L_1, L_2: \text{Mod}^{fg}_R \to \textit{Sets}$ 是把 $0$
送到单元素集合并保持有限积的函子。
设 $t : L_1 \to L_2$ 是函子态射。
则 $t$ 在引理 \ref{formal-defos-lemma-linear-functor} 所保证的函子之间诱导态射
$\widetilde{t} : \widetilde{L}_1 \to \widetilde{L}_2$；
对每个 $M \in \Ob(\text{Mod}^{fg}_R)$，它就是
$\widetilde{t}_M = t_M: \widetilde{L}_1(M) \to \widetilde{L}_2(M)$。
换言之，$t_M: \widetilde{L}_1(M) \to \widetilde{L}_2(M)$
是 $R$-模映射。
\end{lemma}

\begin{proof}
证明从略。
\end{proof}

\noindent
当 $R = K$ 是域时，$K$-线性函子
$L : \text{Mod}^{fg}_K \to \text{Mod}_K$ 由它的值
$L(K)$ 决定。

\begin{lemma}
\label{formal-defos-lemma-linear-functor-over-field}
设 $K$ 是域，并设 $L: \text{Mod}^{fg}_K \to
\text{Mod}_K$ 是 $K$-线性函子。则 $L$ 同构于函子
$L(K) \otimes_K - : \text{Mod}^{fg}_K \to
\text{Mod}_K$。
\end{lemma}

\begin{proof}
对 $V \in \Ob(\text{Mod}^{fg}_K)$，同构
$L(K) \otimes_K V \to L(V)$ 在单纯张量上由
$x \otimes v \mapsto L(f_v)(x)$ 给出，其中
$f_v: K \to V$ 是把 $1 \mapsto v$ 的 $K$-线性映射。
当 $V = K$ 时，这是由 $K$ 上的乘法给出的同构
$L(K) \otimes_K K \to L(K)$。% P10-E0446
对一般的 $V$，由 $V=K$ 的情形以及 $L$ 与有限积交换这一事实
（说明 \ref{formal-defos-remark-linear-functor}），它是同构。
\end{proof}

\noindent
对环 $R$ 和 $R$-模 $M$，令 $R[M]$ 表示这样的 $R$-代数：
其底层 $R$-模是 $R \oplus M$，乘法由
$(r, m) \cdot (r', m') = (rr', rm' + r'm)$ 给出。
当 $M = R$ 时，这是 $R$ 上的对偶数环，记作 $R[\epsilon]$。

\medskip\noindent
现在设 $S$ 是环，并假设 $R$ 是 $S$-代数。
赋值 $M \mapsto R[M]$ 决定函子
$\text{Mod}_R \to S\text{-Alg}/R$，其中 $S\text{-Alg}/R$
表示以 $R$ 为靶的 $S$-代数范畴。注意，
$S\text{-Alg}/R$ 有有限积：若 $A_1 \to R$ 和
$A_2 \to R$ 是两个对象，则 $A_1 \times_R A_2$ 是它们的积。

\begin{lemma}
\label{formal-defos-lemma-preserves-products}
设 $R$ 是 $S$-代数。则上述函子
$\text{Mod}_R \to S\text{-Alg}/R$ 保持有限积。
\end{lemma}

\begin{proof}
这不过是说：若 $M$、$N$ 是 $R$-模，则映射
$R[M \times N] \to R[M] \times_R R[N]$ 是
$S\text{-Alg}/R$ 中的同构。
\end{proof}

\begin{lemma}
\label{formal-defos-lemma-tangent-space-functor}
设 $R$ 是 $S$-代数，并设 $\mathcal{C}$ 是 $S\text{-Alg}/R$
的严格满子范畴，它包含所有
$M \in \Ob(\text{Mod}^{fg}_R)$ 所对应的 $R[M]$。
设 $F: \mathcal{C} \to \textit{Sets}$ 是函子。
假设 $F(R)$ 是单元素集合，并且对任意
$M, N \in \Ob(\text{Mod}^{fg}_R)$，诱导映射
$$
F(R[M] \times_R R[N]) \to F(R[M]) \times F(R[N])
$$
都是双射。则对任意 $M \in \Ob(\text{Mod}^{fg}_R)$，
$F(R[M])$ 都有自然的 $R$-模结构。
\end{lemma}

\begin{proof}
注意，$R \cong R[0]$ 且
$R[M] \times_R R[N] \cong R[M \times N]$，
所以由对 $\mathcal{C}$ 的假设，$R$ 与
$R[M] \times_R R[N]$ 都是 $\mathcal{C}$ 的对象。
因此，对 $F$ 的条件是有意义的。
引理 \ref{formal-defos-lemma-preserves-products} 中的函子
$\text{Mod}_R \to S\text{-Alg}/R$，由对 $\mathcal{C}$ 的假设，
限制为函子 $\text{Mod}^{fg}_R \to \mathcal{C}$。
令 $L$ 为复合
$\text{Mod}^{fg}_R \to \mathcal{C} \to \textit{Sets}$，即
$L(M) = F(R[M])$。
由引理 \ref{formal-defos-lemma-preserves-products} 及对 $F$ 的假设，
$L$ 保持有限积。因此，引理 \ref{formal-defos-lemma-linear-functor} 表明：
对任意 $M \in \Ob(\text{Mod}^{fg}_R)$，
$L(M) = F(R[M])$ 都有自然的 $R$-模结构。
\end{proof}

\begin{definition}
\label{formal-defos-definition-tangent-space-over-R}
设 $\mathcal{C}$ 是引理 \ref{formal-defos-lemma-tangent-space-functor}
中那样的范畴。设 $F : \mathcal{C} \to \textit{Sets}$ 是函子，
且 $F(R)$ 是单元素集合。$F$ 的{\it 切空间 $TF$} 定义为
$F(R[\epsilon])$。
\end{definition}

\noindent
当 $F : \mathcal{C} \to \textit{Sets}$ 满足引理
\ref{formal-defos-lemma-tangent-space-functor} 的假设时，
切空间 $TF$ 具有自然的 $R$-模结构。

\begin{example}
\label{formal-defos-example-tangent-space-functor}
因为 $\mathcal{C}_\Lambda$ 包含有限维向量空间 $V$ 对应的所有
$k[V]$，所以定义 \ref{formal-defos-definition-tangent-space-over-R} 可应用于
$S = \Lambda$、$R = k$、$\mathcal{C} = \mathcal{C}_\Lambda$，
以及预形变函子 $F : \mathcal{C}_\Lambda \to \textit{Sets}$。
切空间为 $TF = F(k[\epsilon])$。
\end{example}

\begin{example}
\label{formal-defos-example-tangent-space-prorepresentable-functor}
在例 \ref{formal-defos-example-tangent-space-functor} 中，设
$F : \mathcal{C}_\Lambda \to \textit{Sets}$ 是 pro-可表函子；
更具体地，对某个
$S \in \Ob(\widehat{\mathcal{C}}_\Lambda)$，设
$F = \underline{S}|_{\mathcal{C}_\Lambda}$。
则 $F$ 与任意极限交换，因而满足引理
\ref{formal-defos-lemma-tangent-space-functor} 的假设。计算得
$$
TF = F(k[\epsilon]) = \Mor_{\mathcal{C}_\Lambda}(S, k[\epsilon])
= \text{Der}_\Lambda(S, k)
$$
更一般地，对有限维 $k$-向量空间 $V$，有
$$
F(k[V]) = \Mor_{\mathcal{C}_\Lambda}(S, k[V]) = \text{Der}_\Lambda(S, V).
$$
% P10-E0447
具体地，与增广 $q : S \to k$ 和 $k[V] \to k$ 相容的
$\Lambda$-代数映射 $f : S \to k[V]$，对应于由
$s \mapsto f(s) - q(s)$ 定义的导子 $D$。
反过来，$\Lambda$-导子 $D : S \to V$ 对应于
$\mathcal{C}_\Lambda$ 中按公式 $f(s) = q(s) + D(s)$ 定义的
$f : S \to k[V]$。由于这些识别具有函子性，
$TF$ 与 $\text{Der}_\Lambda(S, k)$ 上的 $k$-向量空间结构彼此对应
（见引理 \ref{formal-defos-lemma-morphism-linear-functors}）。% P10-E0448
由引理 \ref{formal-defos-lemma-derivations-finite}，可知 $\dim_k TF$ 有限。
\end{example}

\begin{example}
\label{formal-defos-example-tangent-space-classical-prorepresentable-functor}
在经典情形下，例
\ref{formal-defos-example-tangent-space-prorepresentable-functor} 的计算可以简化。
确切地说，此时函子
$F = \underline{S}|_{\mathcal{C}_\Lambda}$ 的切空间就是
$S$ 在 $\Lambda$ 上的相对余切空间，即
$TF = T_{S/\Lambda}$。事实上，当域扩张 $k/k'$ 可分时，
这一结论仍然成立。参见《习题》中的习题
\ref{exercises-exercise-tangent-space-Zariski}。
\end{example}

\begin{lemma}
\label{formal-defos-lemma-morphism-tangent-spaces}
设 $F, G: \mathcal{C} \to \textit{Sets}$ 是满足引理
\ref{formal-defos-lemma-tangent-space-functor} 假设的函子。
设 $t : F \to G$ 是函子态射。
对任意 $M \in \Ob(\text{Mod}^{fg}_R)$，映射
$t_{R[M]}: F(R[M]) \to G(R[M])$ 是 $R$-模映射，
其中 $F(R[M])$ 与 $G(R[M])$ 配备引理
\ref{formal-defos-lemma-tangent-space-functor} 给出的 $R$-模结构。
特别地，$t_{R[\epsilon]} : TF \to TG$ 是 $R$-模映射。
\end{lemma}

\begin{proof}
这由引理 \ref{formal-defos-lemma-morphism-linear-functors} 得出。
\end{proof}

\begin{example}
\label{formal-defos-example-tangent-space-map-prorepresentable-functor}
假设 $f : R \to S$ 是 $\widehat{\mathcal{C}}_\Lambda$ 中的环映射。
令 $F = \underline{R}|_{\mathcal{C}_\Lambda}$，
$G = \underline{S}|_{\mathcal{C}_\Lambda}$。
环映射 $f$ 诱导函子变换 $G \to F$。
由引理 \ref{formal-defos-lemma-morphism-tangent-spaces}，得到 $k$-线性映射
$TG \to TF$。这就是映射
$$
TG = \text{Der}_\Lambda(S, k) \longrightarrow \text{Der}_\Lambda(R, k) = TF
$$
这一点来自例 \ref{formal-defos-example-tangent-space-prorepresentable-functor}
中的典范识别
$F(k[V]) = \text{Der}_\Lambda(R, V)$ 和
$G(k[V]) = \text{Der}_\Lambda(S, V)$，
以及切空间上映射的计算规则。
\end{example}

\begin{lemma}
\label{formal-defos-lemma-tangent-space-tensor}
设 $F: \mathcal{C} \to \textit{Sets}$ 是满足引理
\ref{formal-defos-lemma-tangent-space-functor} 假设的函子。
假设 $R = K$ 是域。则对任意有限维 $K$-向量空间 $V$，有
$F(K[V]) \cong TF \otimes_K V$。
\end{lemma}

\begin{proof}
这由引理 \ref{formal-defos-lemma-linear-functor-over-field} 得出。
\end{proof}





\section{预形变范畴的切空间}
\label{formal-defos-section-tangent-spaces}

\noindent
我们将用一般的定义 \ref{formal-defos-definition-tangent-space-over-R}
定义预形变函子的切空间。例 \ref{formal-defos-example-tangent-space-functor}
已经具体说明了这一点。考察同构类的关联函子，
即可把该定义应用于预形变范畴。

\begin{definition}
\label{formal-defos-definition-tangent-space}
设 $\mathcal{F}$ 是预形变范畴。
$\mathcal{F}$ 的{\it 切空间 $T \mathcal{F}$} 是纤维范畴
$\mathcal F(k[\epsilon])$ 中对象的同构类集合
$\overline{\mathcal{F}}(k[\epsilon])$。
\end{definition}

\noindent
因此，$T \mathcal{F}$ 就是关联函子
$\overline{\mathcal{F}}: \mathcal{C}_\Lambda \to \textit{Sets}$
的切空间。当 $\mathcal{F}$ 满足 (S2) 时，
它具有自然的向量空间结构；事实上，只要
$\overline{\mathcal{F}}$ 满足 (S2) 即可。

\begin{lemma}
\label{formal-defos-lemma-tangent-space-vector-space}
设 $\mathcal{F}$ 是预形变范畴，且
$\overline{\mathcal{F}}$ 满足 (S2)\footnote{例如，若
$\mathcal{F}$ 满足 (S2)，则此条件成立；见引理
\ref{formal-defos-lemma-S1-S2-associated-functor}。}。
则 $T \mathcal{F}$ 具有自然的 $k$-向量空间结构。
对任意有限维向量空间 $V$，都有
$\overline{\mathcal{F}}(k[V]) = T\mathcal{F} \otimes_k V$，
且这一等式关于 $V$ 具有函子性。
\end{lemma}

\begin{proof}
记
$F = \overline{\mathcal{F}} : \mathcal{C}_\Lambda \to \textit{Sets}$。
这是预形变函子，并且 $F$ 满足 (S2)。由引理
\ref{formal-defos-lemma-S2-extensions}
（以及说明 \ref{formal-defos-remark-compare-S1-S2-schlessinger}
对集合值函子的对应表述）% P10-E0449
可知，对每个有限维向量空间 $V$ 和每个
$A \in \Ob(\mathcal{C}_\Lambda)$，映射
$$
F(A \times_k k[V]) \longrightarrow F(A) \times F(k[V])
$$
都是双射。特别地，取 $A = k[W]$，可得
$F(k[W] \times_k k[V]) = F(k[W]) \times F(k[V])$。
换言之，引理 \ref{formal-defos-lemma-tangent-space-functor} 的假设成立，
从而 $TF = T \mathcal{F}$ 具有自然的 $k$-向量空间结构。
最后一个断言由引理 \ref{formal-defos-lemma-tangent-space-tensor} 得出。
\end{proof}

\noindent
预形变范畴的态射在切空间上诱导映射。

\begin{definition}
\label{formal-defos-definition-differential}
设 $\varphi : \mathcal{F} \to \mathcal{G}$ 是预形变范畴的态射。% P10-E0450
$\varphi$ 的{\it 微分
$d \varphi : T \mathcal{F} \to T \mathcal{G}$}
是把函子态射
$\overline{\varphi}: \overline{\mathcal{F}} \to \overline{\mathcal{G}}$
在 $A = k[\epsilon]$ 处取值所得的映射。
\end{definition}

\begin{lemma}
\label{formal-defos-lemma-k-linear-differential}
设 $\varphi : \mathcal{F} \to \mathcal{G}$ 是预形变范畴的态射。
假设 $\overline{\mathcal{F}}$ 和 $\overline{\mathcal{G}}$
都满足 (S2)。则
$d \varphi : T \mathcal{F} \to T \mathcal{G}$ 是 $k$-线性的。
\end{lemma}

\begin{proof}
在引理 \ref{formal-defos-lemma-tangent-space-vector-space} 的证明中，
已经看到 $\overline{\mathcal{F}}$ 和 $\overline{\mathcal{G}}$
满足引理 \ref{formal-defos-lemma-tangent-space-functor} 的假设。
因此，本引理由引理 \ref{formal-defos-lemma-morphism-tangent-spaces} 得出。
\end{proof}

\begin{remark}
\label{formal-defos-remark-tangent-space-cofibered-groupoid}
可以把切空间和微分的概念如下全局化到任意群胚余纤维范畴。
设 $\mathcal{F}$ 是 $\mathcal{C}_\Lambda$ 上的群胚余纤维范畴，
且 $x \in \Ob(\mathcal{F}(k))$。
如说明 \ref{formal-defos-remark-localize-cofibered-groupoid}，
得到预形变范畴 $\mathcal{F}_x$。定义
$$
T_x\mathcal{F} = T\mathcal{F}_x
$$
为 $\mathcal{F}$ 在 $x$ 处的{\it 切空间}。
若 $\varphi : \mathcal{F} \to \mathcal{G}$ 是
$\mathcal{C}_\Lambda$ 上群胚余纤维范畴的态射，且
$x \in \Ob(\mathcal{F}(k))$，则有诱导态射
$\varphi_x: \mathcal{F}_x \to \mathcal{G}_{\varphi(x)}$。
把 $\varphi$ 在 $x$ 处的{\it 微分
$d_x \varphi : T_x \mathcal{F} \to T_{\varphi(x)} \mathcal{G}$}
定义为映射
$d \varphi_x: T \mathcal{F}_x \to T \mathcal{G}_{\varphi(x)}$。
若 $\mathcal{F}$ 和 $\mathcal{G}$ 都满足 (S2)，
则所有这些切空间都具有自然的 $k$-向量空间结构，
而所有微分
$d_x \varphi : T_x \mathcal{F} \to T_{\varphi(x)} \mathcal{G}$
都是 $k$-线性的（使用引理
\ref{formal-defos-lemma-S1-S2-localize} 和
\ref{formal-defos-lemma-k-linear-differential}）。
\end{remark}

\noindent
在经典情形下，或当 $k/k'$ 是可分域扩张时，
以下观察没有实质内容，因为这时
$\text{Der}_\Lambda(k, k)$ 和 $\text{Der}_\Lambda(k, V)$ 都为零。
有典范识别
$$
\Mor_{\mathcal{C}_\Lambda}(k, k[\epsilon]) =
\text{Der}_\Lambda(k, k).
$$
确切地，对 $D \in \text{Der}_\Lambda(k, k)$，
令 $f_D : k \to k[\epsilon]$ 为映射
$a \mapsto a + D(a)\epsilon$。更一般地，
给定 $k$ 上的有限维向量空间 $V$，有
$$
\Mor_{\mathcal{C}_\Lambda}(k, k[V]) =
\text{Der}_\Lambda(k, V)
$$
并且仍用 $f_D$ 表示与导子 $D$ 相伴的映射。还有
$$
\Mor_{\mathcal{C}_\Lambda}(k[W], k[V]) =
\Hom_k(V, W) \oplus \text{Der}_\Lambda(k, V)
$$
% P10-E0451
其中 $(\varphi, D)$ 对应于映射
$f_{\varphi, D} : a + w \mapsto a + \varphi(w) + D(a)$。
有时把由导子 $D : k \to V$ 决定的 $k[V]$ 自同构写成
$f_{1, D} : a + v \to a + v + D(a)$。% P10-E0452
注意，$f_{1, D} \circ f_{1, D'} = f_{1, D + D'}$。

\medskip\noindent
设 $\mathcal{F}$ 是 $\mathcal{C}_\Lambda$ 上的预形变范畴，
并设 $x_0 \in \Ob(\mathcal{F}(k))$。
由上可得典范映射
$$
\gamma_V :
\text{Der}_\Lambda(k, V)
\longrightarrow
\overline{\mathcal{F}}(k[V])
$$
定义为 $D \mapsto f_{D, *}(x_0)$。此外，还有作用
$$
a_V : \text{Der}_\Lambda(k, V) \times \overline{\mathcal{F}}(k[V])
\longrightarrow
\overline{\mathcal{F}}(k[V])
$$
定义为 $(D, x) \mapsto f_{1, D, *}(x)$。
这两个映射相容；也就是说，由这些映射复合的直接计算可得
$f_{1, D, *}f_{D', *}x_0 = f_{D + D', *}x_0$。
注意，$\gamma_V$ 与 $a_V$ 不依赖于 $x_0$ 的选取，
因为 $x_0$ 在同构意义下唯一。

\begin{lemma}
\label{formal-defos-lemma-action-linear}
设 $\mathcal{F}$ 是 $\mathcal{C}_\Lambda$ 上的预形变范畴。
若 $\overline{\mathcal{F}}$ 满足 (S2)，则映射 $\gamma_V$
是 $k$-线性的，并且 $a_V(D, x) = x + \gamma_V(D)$。
\end{lemma}

\begin{proof}
在引理 \ref{formal-defos-lemma-tangent-space-vector-space} 的证明中，
已经看到函子 $V \mapsto \overline{\mathcal{F}}(k[V])$
把 $0$ 送到单元素集合，并把积送到积。
函子 $V \mapsto \text{Der}_\Lambda(k, V)$ 也如此。
因此，由引理 \ref{formal-defos-lemma-morphism-linear-functors}，
$\gamma_V$ 是线性的。
设 $D : k \to V$ 是 $\Lambda$-导子。
令 $D_1 : k \to V^{\oplus 2}$ 等于 $a \mapsto (D(a), 0)$。
则图
$$
\xymatrix{
k[V \times V] \ar[r]_{+} \ar[d]^{f_{1, D_1}} & k[V] \ar[d]^{f_{1, D}} \\
k[V \times V] \ar[r]^{+} & k[V]
}
$$
交换。展开定义，并利用
$\overline{F}(V \times V) = \overline{F}(V) \times \overline{F}(V)$，% P10-E0453
这意味着对所有 $x_1, x_2 \in \overline{F}(V)$，有
$a_D(x_1) + x_2 = a_D(x_1 + x_2)$。
因此，只需证明 $a_V(D, 0) = 0 + \gamma_V(D)$，
其中 $0 \in \overline{F}(V)$ 是零向量。
按定义，这是元素 $f_{0, *}(x_0)$。
由于 $f_D = f_{1, D} \circ f_0$，所需结论随即得出。
\end{proof}

\noindent
上述构造的一个特殊情形是映射
\begin{equation}
\label{formal-defos-equation-map}
\gamma : \text{Der}_\Lambda(k, k) \longrightarrow T\mathcal{F}
\end{equation}
以及作用
\begin{equation}
\label{formal-defos-equation-action}
a : \text{Der}_\Lambda(k, k) \times T\mathcal{F} \longrightarrow T\mathcal{F}
\end{equation}
它们对任意预形变范畴 $\mathcal{F}$ 都有定义。
注意，若 $\varphi : \mathcal{F} \to \mathcal{G}$ 是预形变范畴的态射，
则得到交换图
$$
\vcenter{
\xymatrix{
\text{Der}_\Lambda(k, k) \ar[r]_-\gamma \ar[rd]_\gamma &
T\mathcal{F} \ar[d]_{d\varphi} \\
& T\mathcal{G}
}
}
\quad\text{和}\quad
\vcenter{
\xymatrix{
\text{Der}_\Lambda(k, k) \times T\mathcal{F} \ar[r]_-a
\ar[d]_{1 \times d\varphi} &
T\mathcal{F} \ar[d]^{d\varphi} \\
\text{Der}_\Lambda(k, k) \times T\mathcal{G} \ar[r]^-a &
T\mathcal{G}
}
}
$$





\section{万有形式对象}
\label{formal-defos-section-versal-objects}

\noindent
万有形式对象的存在迫使 $\mathcal{F}$ 具有性质 (S1)。

\begin{lemma}
\label{formal-defos-lemma-versal-object-S1}
设 $\mathcal{F}$ 是预形变范畴。
假设 $\mathcal{F}$ 有万有形式对象。
则 $\mathcal{F}$ 满足 (S1)。
\end{lemma}

\begin{proof}
设 $\xi$ 是 $\mathcal{F}$ 的万有形式对象。设
$$
\xymatrix{
           & x_2 \ar[d] \\
x_1 \ar[r] & x
}
$$
是 $\mathcal{F}$ 中的图，并且 $x_2 \to x$
位于一个满环映射上方。由于自然态射
$\widehat{\mathcal{F}}|_{\mathcal{C}_\Lambda} \xrightarrow{\sim} \mathcal{F}$
是等价（见说明 \ref{formal-defos-remark-restrict-completion}），
也可以把这个图视为 $\widehat{\mathcal{F}}$ 中的图。
由引理 \ref{formal-defos-lemma-versal-object-quasi-initial}，存在态射
$\xi \to x_1$；所以由说明 \ref{formal-defos-remark-versal-object}，
还得到态射 $\xi \to x_2$，使图
$$
\xymatrix{
\xi \ar[r] \ar[d]          & x_2 \ar[d] \\
x_1 \ar[r] & x
}
$$
交换。若 $x_1 \to x$ 与 $x_2 \to x$ 分别位于环映射
$A_1 \to A$ 与 $A_2 \to A$ 上方，
则把 $\xi$ 推前到 $A_1 \times_A A_2$，
便得到 (S1) 所要求的对象 $y$。
\end{proof}

\noindent
当余纤维范畴满足 (S1) 和 (S2) 时，
可以如下刻画万有形式对象。

\begin{lemma}
\label{formal-defos-lemma-versal-criterion}
设 $\mathcal{F}$ 是满足 (S1) 和 (S2) 的预形变范畴。
设 $\xi$ 是 $\mathcal{F}$ 的形式对象，对应于
$\underline{\xi} : \underline{R}|_{\mathcal{C}_\Lambda} \to \mathcal{F}$；
见说明 \ref{formal-defos-remark-formal-objects-yoneda}。
则 $\xi$ 是万有的，当且仅当以下两个条件成立：
\begin{enumerate}
\item 切空间上的映射
$d\underline{\xi} : T\underline{R}|_{\mathcal{C}_\Lambda} \to T\mathcal{F}$
是满射；并且
\item 给定 $\widehat{\mathcal{F}}$ 中的图
$$
\vcenter{
\xymatrix{
            &  y \ar[d] \\
\xi \ar[r]  &  x
}
}
\quad\text{位于下图上方}\quad
\vcenter{
\xymatrix{
         &   B  \ar[d]^{f} \\
R \ar[r] &   A
}
}
$$
其中底图属于 $\widehat{\mathcal{C}}_\Lambda$，
且 $B \to A$ 是阿廷环的小扩张，则存在环映射
$R \to B$，使图
$$
\xymatrix{
         &   B  \ar[d]^{f} \\
R \ar[ur] \ar[r] &   A
}
$$
交换。
\end{enumerate}
\end{lemma}

\begin{proof}
若 $\xi$ 是万有的，则 (1) 由引理
\ref{formal-defos-lemma-smooth-morphism-essentially-surjective} 成立，
(2) 由说明 \ref{formal-defos-remark-versal-object} 成立。
反过来，假设 (1) 和 (2) 成立。由说明
\ref{formal-defos-remark-versal-object}，必须证明：给定 (2) 中那样的
$\widehat{\mathcal{F}}$ 中的图，存在 $\xi \to y$，使图
$$
\xymatrix{
            &  y \ar[d] \\
\xi \ar[ur] \ar[r]  &  x
}
$$
交换。设 $b : R \to B$ 是 (2) 所保证的映射。
记 $y' = b_*\xi$，并选取给定态射 $\xi \to x$
在 $R \to B \to A$ 上方的分解
$\xi \to y' \to x$。由 (S1)，得到交换图
$$
\vcenter{
\xymatrix{
z  \ar[r] \ar[d]          &  y \ar[d] \\
y' \ar[r]  &  x
}
}
\quad\text{位于下图上方}\quad
\vcenter{
\xymatrix{
B \times_A B \ar[d] \ar[r] &   B  \ar[d]^{f} \\
B \ar[r]^{f} &   A .
}
}
$$
令 $I = \Ker(f)$。令
$\overline{g} : B \times_A B \to k[I]$
为环映射 $(u, v) \mapsto \overline{u} \oplus (v - u)$；
参见引理 \ref{formal-defos-lemma-lifting-along-small-extension}。
由 (1)，存在位于某个环映射
$i : R \to k[\epsilon]$ 上方的态射
$\xi \to \overline{g}_*z$。
选取阿廷商 $b_1 : R \to B_1$，使
$b : R \to B$ 和 $i : R \to k[\epsilon]$ 都经
$R \to B_1$ 分解；也就是说，得到
$h : B_1 \to B$ 和 $i' : B_1 \to k[\epsilon]$。
选取推前 $y_1 = b_{1, *}\xi$，选取 $\xi \to y'$
在 $R \to B_1 \to B$ 上方的分解
$\xi \to y_1 \to y'$，并选取
$\xi \to \overline{g}_*z$ 在
$R \to B_1 \to k[\epsilon]$ 上方的分解
$\xi \to y_1 \to \overline{g}_*z$。
再次应用 (S1)，得到
$$
\vcenter{
\xymatrix{
z_1  \ar[r] \ar[d] & z \ar[r] \ar[d] & y \ar[d] \\
y_1 \ar[r] & y' \ar[r] &  x
}
}
\quad\text{位于下图上方}\quad
\vcenter{
\xymatrix{
B_1 \times_A B \ar[d] \ar[r] & B \times_A B \ar[r] \ar[d] & B \ar[d]^{f} \\
B_1 \ar[r]  & B \ar[r] &  A .
}
}
$$
注意，引理 \ref{formal-defos-lemma-lifting-along-small-extension}
中（用 $h$ 定义的）映射
$g : B_1 \times_A B \to k[I]$，
是 $B_1 \times_A B \to B \times_A B$
与上述映射 $\overline{g}$ 的复合。
由构造，存在态射
$y_1 \to g_*z_1 \cong \overline{g}_*z$！
因此，引理 \ref{formal-defos-lemma-lifting-along-small-extension}
可应用于上述图中的外侧矩形，给出态射
$y_1 \to y$；再前复合 $\xi \to y_1$，
便得到所需态射 $\xi \to y$。
\end{proof}

\noindent
若 $\mathcal{F}$ 具有性质 (S1)，则“存在提升的最大商”存在。
精确表述如下。

\begin{lemma}
\label{formal-defos-lemma-largest-closed-where-lift}
设 $\mathcal{F}$ 是 $\mathcal{C}_\Lambda$ 上满足 (S1)
的群胚余纤维范畴。设 $B \to A$ 是
$\mathcal{C}_\Lambda$ 中的满射，其核 $I$ 被
$\mathfrak m_B$ 零化。设 $x \in \mathcal{F}(A)$。
理想集合
$$
\mathcal{J} = \{ J \subset I \mid
\text{存在位于 }B/J \to A\text{ 上方的 }y \to x\}
$$
有最小元。% P10-E0455
\end{lemma}

\begin{proof}
注意，$\mathcal{J}$ 非空，因为 $I \in \mathcal{J}$。
此外，若 $J \in \mathcal{J}$ 且 $J \subset J' \subset I$，
则 $J' \in \mathcal{J}$，因为可以把对象 $y$ 推前为
$B/J'$ 上方的对象 $y'$。设 $J$ 和 $K$ 是上述集合的元素。
我们断言 $J \cap K \in \mathcal{J}$，这将证明引理。
因为 $I$ 是 $k$-向量空间，可以找到理想
$J \subset J' \subset I$，使 $J \cap K = J' \cap K$ 且
$J' + K = I$。由上可用 $J'$ 代替 $J$，并假设
$J + K = I$。此时
$$
A/(J \cap K) = A/J \times_{A/I} A/K.
$$
% P10-E0454
因此，由 (S1) 以及假设存在于 $A/J$ 和 $A/K$ 上方的元素，
可知存在映到 $x$ 的元素
$z \in \mathcal{F}(A/(J \cap K))$。
\end{proof}

\noindent
稍后将改进下述结果。

\begin{lemma}
\label{formal-defos-lemma-versal-object-existence}
设 $\mathcal{F}$ 是 $\mathcal{C}_\Lambda$ 上的群胚余纤维范畴。
假设以下条件成立：
\begin{enumerate}
\item $\mathcal{F}$ 是预形变范畴。
\item $\mathcal{F}$ 满足 (S1)。
\item $\mathcal{F}$ 满足 (S2)。
\item $\dim_k T\mathcal{F}$ 有限。
\end{enumerate}
则 $\mathcal{F}$ 有万有形式对象。
\end{lemma}

\begin{proof}
假设 (1)、(2)、(3) 和 (4) 成立。
选取对象 $R \in \Ob(\widehat{\mathcal{C}}_\Lambda)$，
使 $\underline{R}|_{\mathcal{C}_\Lambda}$ 光滑；
见引理 \ref{formal-defos-lemma-exists-smooth}。
令 $r = \dim_k T\mathcal{F}$，并置
$S = R[[X_1, \ldots, X_r]]$。

\medskip\noindent
我们将对 $n \geq 2$ 归纳构造二元组
$(J_n, f_{n - 1} : \xi_n \to \xi_{n - 1})$，
其中 $J_n \subset S$ 构成递减理想序列，% P10-E0456
而 $f_{n - 1} : \xi_n \to \xi_{n - 1}$ 是
$\mathcal{F}$ 中位于投影
$S/J_n \to S/J_{n - 1}$ 上方的态射。

\medskip\noindent
第 1 步。令 $J_1 = \mathfrak m_S$。
令 $\xi_1$ 为 $k = S/J_1 = S/\mathfrak m_S$ 上方的
$\mathcal{F}$ 的对象；它在唯一同构意义下唯一。% P10-E0457

\medskip\noindent
第 2 步。令
$J_2 = \mathfrak m_S^2 + \mathfrak{m}_R S$。
则 $S/J_2 = k[V]$，其中
$V = kX_1 \oplus \ldots \oplus kX_r$。
由 $\overline{\mathcal{F}}$ 的 (S2)，得到双射
$$
\overline{\mathcal{F}}(S/J_2)
\longrightarrow
T\mathcal{F} \otimes_k V,
$$
见引理 \ref{formal-defos-lemma-S1-S2-associated-functor} 和
\ref{formal-defos-lemma-tangent-space-vector-space}。
选取 $T\mathcal{F}$ 的一组基
$\theta_1, \ldots, \theta_r$，并置
$\xi_2 = \sum \theta_i \otimes X_i \in \Ob(\mathcal{F}(S/J_2))$。% P10-E0458
这样选取的要点在于
$$
d\xi_2 :
\Mor_{\mathcal{C}_\Lambda}(S/J_2, k[\epsilon])
\longrightarrow
T\mathcal{F}
$$
是满射。令 $f_1 : \xi_2 \to \xi_1$ 为唯一态射。

\medskip\noindent
归纳步骤。假设对某个 $n \geq 2$，已经构造了
$(J_n, f_{n - 1} : \xi_n \to \xi_{n - 1})$。
满足以下条件的理想 $J \subset S$ 的集合有最小元
$J_{n + 1}$：
(a) $\mathfrak m_S J_n \subset J \subset J_n$；并且
(b) 存在位于 $S/J \to S/J_n$ 上方的态射
$\xi_{n + 1} \to \xi_n$；见引理
\ref{formal-defos-lemma-largest-closed-where-lift}。
令 $f_n : \xi_{n + 1} \to \xi_n$ 为
$\mathcal{F}$ 中位于 $S/J_{n + 1} \to S/J_n$ 上方的任意态射。

\medskip\noindent
令 $J = \bigcap J_n$，$\overline{S} = S/J$，并令
$\overline{J}_n = J_n/J$。由引理
\ref{formal-defos-lemma-m-adic-topology}，理想序列
$(\overline{J}_n)$ 在 $\overline{S}$ 上诱导
$\mathfrak m_{\overline{S}}$-进拓扑。
由于 $(\xi_n, f_n)$ 是
$\widehat{\mathcal{F}}_\mathcal{I}(\overline{S})$ 的对象，
其中 $\mathcal{I}$ 是 $\overline{S}$ 的滤过
$(\overline{J}_n)$，可见 $(\xi_n, f_n)$ 诱导
$\widehat{\mathcal{F}}(\overline{S})$ 的对象 $\xi$；
见引理 \ref{formal-defos-lemma-formal-objects-different-filtration}。% P10-E0459

\medskip\noindent
我们证明 $\xi$ 是万有的。为证明万有性，只需核验引理
\ref{formal-defos-lemma-versal-criterion} 的条件 (1) 与 (2)。
条件 (1) 来自上述第 2 步中对 $\xi_2$ 的选取。
假设给定 $\widehat{\mathcal{F}}$ 中的图
$$
\vcenter{
\xymatrix{
            &  y \ar[d] \\
\xi \ar[r]  &  x
}
}
\quad\text{位于下图上方}\quad
\vcenter{
\xymatrix{
         &   B  \ar[d]^{f} \\
\overline{S} \ar[r] &   A
}
}
$$
其中底图属于 $\widehat{\mathcal{C}}_\Lambda$，
且 $f: B \to A$ 是阿廷环的小扩张。
必须证明存在使右图补全的映射 $\overline{S} \to B$。
选取 $n$，使 $\overline{S} \to A$ 经
$\overline{S} \to S/J_n$ 分解。
这是可能的，因为如上所见，序列 $(\overline{J}_n)$
诱导 $\mathfrak m_{\overline{S}}$-进拓扑。
$\xi$ 沿 $\overline{S} \to S/J_n$ 的推前为 $\xi_n$。
可以把 $\xi \to x$ 分解为 $\xi \to \xi_n \to x$，
从而得到 $\mathcal{F}$ 中的图
$$
\vcenter{
\xymatrix{
            &  y \ar[d] \\
\xi_n \ar[r]  &  x
}
}
\quad\text{位于下图上方}\quad
\vcenter{
\xymatrix{
         &   B  \ar[d]^{f} \\
S/J_n \ar[r] &   A .
}
}
$$
为核验引理 \ref{formal-defos-lemma-versal-criterion} 的条件 (2)，
只需补全图
$$
\xymatrix{
S/J_{n + 1} \ar[d] \ar@{-->}[r] & B \ar[d]^{f} \\
S/J_n   \ar[r] & A
}
$$
或者等价地，补全图
$$
\xymatrix{
  &  S/J_n \times_A B \ar[d]^{p_1} \\
S/J_{n + 1} \ar@{-->}[ur] \ar[r] &  S/J_n.
}
$$
若 $p_1$ 有截面，则已经完成。否则，由引理
\ref{formal-defos-lemma-fiber-product-CLambda} (2)，
$p_1$ 是小扩张；所以由引理
\ref{formal-defos-lemma-essential-surjection} (4)，
$p_1$ 是本质满射。回忆
$S = R[[X_1, \ldots, X_r]]$，且选取 $R$ 时使
$\underline{R}|_{\mathcal{C}_\Lambda}$ 光滑。
因此，存在映射 $h : R \to B$，提升映射
$R \to S \to S/J_n \to A$。
由幂级数环的泛性质，存在 $R$-代数映射
$h : S = R[[X_1, \ldots, X_r]] \to B$，% P10-E0460
提升给定映射 $S \to S/J_n \to A$。
这诱导映射 $g: S \to S/J_n \times_A B$，
使图中的实线方块
$$
\xymatrix{
S \ar[d] \ar[r]_-g  &  S/J_n \times_A B \ar[d]^{p_1} \\
S/J_{n + 1} \ar@{-->}[ur] \ar[r] &  S/J_n
}
$$
交换。因为 $p_1$ 是本质满射，所以 $g$ 是满射。
我们断言 $S$ 的理想 $K = \Ker(g)$ 满足上述归纳步骤构造
$J_{n + 1}$ 时的条件 (a) 与 (b)。
确实，$K \subset J_n$ 显然，而因为 $p_1$ 是小扩张，
$\mathfrak m_SJ_n \subset K$；这证明了 (a)。
把 (S1) 应用于
$$
\xymatrix{
            &  y \ar[d] \\
\xi_n \ar[r]  &  x,
}
$$
可得 $\xi_n$ 到 $S/K \cong S/J_n \times_A B$ 的提升，
所以 (b) 成立。由于 $J_{n + 1}$ 是具有性质 (a) 与 (b)
的最小理想，这蕴含 $J_{n + 1} \subset K$。
因此，所需映射
$S/J_{n+1} \to S/K \cong S/J_n \times_A B$ 存在。
\end{proof}

\begin{remark}
\label{formal-defos-remark-compare-schlessinger-H3-pre}
设 $F : \mathcal{C}_\Lambda \to \textit{Sets}$ 是满足 (S1)
和 (S2) 的预形变函子。条件 $\dim_k TF < \infty$
正是 Schlessinger 论文中的条件 (H3)。
回忆，(S1) 与 (S2) 分别对应于条件 (H1) 与 (H2)；
见说明 \ref{formal-defos-remark-compare-S1-S2-schlessinger}。
因此，引理 \ref{formal-defos-lemma-versal-object-existence} 告诉我们：
$$
(H1) + (H2) + (H3)
\Rightarrow
\text{存在万有形式对象}
$$
对预形变函子成立。我们将在说明
\ref{formal-defos-remark-compare-schlessinger-H3} 中说明它与包络的联系。
\end{remark}





\section{极小的万有形式对象}
\label{formal-defos-section-minimal-versal}

\noindent
我们稍作准备，以理解万有形式对象的（非）唯一性。
结果表明，若预形变范畴有万有形式对象，
则它有极小的万有形式对象，并且任意两个这样的对象都同构。
此外，在用幂级数扩张替换基环的意义下，
所有万有形式对象都“差不多”相同。

\medskip\noindent
设 $\mathcal{F}$ 是 $\mathcal{C}_\Lambda$ 上的群胚余纤维范畴。
对 $\mathcal{F}$ 中位于
$A \in \Ob(\mathcal{C}_\Lambda)$ 上方的每个对象 $x$，
考虑范畴 $\mathcal{S}_x$，其对象为
$$
\Ob(\mathcal{S}_x) =
\{x' \to x \mid x' \to x\text{ 所在的子环为 }A' \subset A\}
$$
而态射是 $x$ 上方的态射。
对 $\mathcal{F}$ 中位于
$\mathcal{C}_\Lambda$ 的 $f : B \to A$ 上方的每个
$y \to x$，有如下定义的函子
$f_* : \mathcal{S}_y \to \mathcal{S}_x$：
给定 $B' \subset B$ 上方的 $y' \to y$，令
$A' = f(B')$，并令位于 $B' \to f(B')$ 上方的
$y' \to x'$ 是 $y'$ 的推前。% P10-E0461
由群胚余纤维范畴的公理，得到位于
$f(B') \to A$ 上方的唯一态射 $x' \to x$，使图
$$
\xymatrix{
y' \ar[d] \ar[r] & x' \ar[d] \\
y \ar[r] & x
}
$$
交换。于是 $x' \to x$ 是 $\mathcal{S}_x$ 的对象。
若 $\mathcal{S}_x$ 中的任意态射
$(x'_1 \to x) \to (x' \to x)$ 都是同构，
即 $x'$ 与 $x'_1$ 定义在 $A$ 的同一子环上，
则称 $\mathcal{S}_x$ 的对象 $x' \to x$ 是{\it 极小的}。
因为 $A$ 作为 $\Lambda$-模具有有限长度，
所以极小对象总是存在。

\begin{lemma}
\label{formal-defos-lemma-smallest-where-descends}
设 $\mathcal{F}$ 是 $\mathcal{C}_\Lambda$ 上满足 (S1)
的群胚余纤维范畴。
\begin{enumerate}
\item 对 $\mathcal{F}$ 中的 $y \to x$，
$\mathcal{S}_y$ 的极小对象映到 $\mathcal{S}_x$ 的极小对象。
\item 对 $\mathcal{F}$ 中位于
$\mathcal{C}_\Lambda$ 的满射 $f : B \to A$ 上方的
$y \to x$，$\mathcal{S}_x$ 的每个极小对象都是
$\mathcal{S}_y$ 的某个极小对象的像。
\end{enumerate}
\end{lemma}

\begin{proof}
(1) 的证明。设 $y \to x$ 位于 $f : B \to A$ 上方。
设位于 $B' \subset B$ 上方的 $y' \to y$ 是
$\mathcal{S}_y$ 的极小对象。令
$$
\vcenter{
\xymatrix{
y' \ar[d] \ar[r] & x' \ar[d] \\
y \ar[r] & x
}
}
\quad\text{位于下图上方}\quad
\vcenter{
\xymatrix{
B' \ar[d] \ar[r] & f(B') \ar[d] \\
B \ar[r] & A
}
}
$$
如上述 $f_*$ 的构造所给。假设
$(x'' \to x) \to (x' \to x)$ 是 $\mathcal{S}_x$ 的态射，
且 $x'' \to x'$ 位于 $A'' \subset f(B')$ 上方。
由 (S1)，存在位于
$B' \times_{f(B')} A'' \to B'$ 上方的 $y'' \to y'$。
由于 $y' \to y$ 极小，可知
$B' \times_{f(B')} A'' \to B'$ 是同构，
这蕴含 $A'' = f(B')$；即 $x' \to x$ 极小。

\medskip\noindent
(2) 的证明。假设 $f : B \to A$ 满射，且
$y \to x$ 位于 $f$ 上方。
设 $x' \to x$ 是 $\mathcal{S}_x$ 中位于
$A' \subset A$ 上方的极小对象。
由 (S1)，存在位于
$B' = f^{-1}(A') = B \times_A A' \to B$ 上方的
$y' \to y$，它在 $\mathcal{S}_x$ 中的像是 $x' \to x$。
所以 $f_*(y' \to y) = x' \to x$。
选取 $\mathcal{S}_y$ 中的态射
$(y'' \to y) \to (y' \to y)$，其中 $y'' \to y$ 是极小对象
（由引理之前关于长度的说明，这是可能的）。
于是 $f_*(y'' \to y)$ 是 $\mathcal{S}_x$ 中映到
$x' \to x$ 的对象（由 $f_*$ 的函子性）；
由于 $x' \to x$ 极小，它与 $x' \to x$ 同构。
\end{proof}

\begin{lemma}
\label{formal-defos-lemma-smallest-where-descends-versal}
设 $\mathcal{F}$ 是 $\mathcal{C}_\Lambda$ 上满足 (S1)
的群胚余纤维范畴。设 $\xi$ 是位于 $R$ 上方的
$\mathcal{F}$ 的万有形式对象。
存在位于 $R' \subset R$ 上方的态射 $\xi' \to \xi$，
具有以下极小性性质：
\begin{enumerate}
\item 对每个 $f : R \to A$，其中
$A \in \Ob(\mathcal{C}_\Lambda)$，推前图
$$
\vcenter{
\xymatrix{
\xi' \ar[d] \ar[r] & x' \ar[d] \\
\xi \ar[r] & x
}
}
\quad\text{位于下图上方}\quad
\vcenter{
\xymatrix{
R' \ar[d] \ar[r] & f(R') \ar[d] \\
R \ar[r] & A
}
}
$$
产生 $\mathcal{S}_x$ 的极小对象 $x' \to x$；并且
\item 对任意形式对象态射 $\xi'' \to \xi'$，
相应环映射 $R'' \to R'$ 都是满射。
\end{enumerate}
\end{lemma}

\begin{proof}
写成 $\xi = (R, \xi_n, f_n)$。令 $R'_1 = k$ 且
$\xi'_1 = \xi_1$。假设对 $m \leq n$，已经构造了
$\mathcal{S}_{\xi_m}$ 中位于
$R'_m \subset R/\mathfrak m_R^m$ 上方的极小对象
$\xi'_m \to \xi_m$；并且对 $m \leq n - 1$，
已经构造了与 $f_m$ 相容的态射
$f'_m : \xi'_{m + 1} \to \xi'_m$。
由引理 \ref{formal-defos-lemma-smallest-where-descends} (2)，
存在位于
$R'_{n + 1} \subset R/\mathfrak m_R^{n + 1}$ 上方的极小对象
$\xi'_{n + 1} \to \xi_{n + 1}$，其像是位于
$R'_n \subset R/\mathfrak m_R^n$ 上方的
$\xi'_n \to \xi_n$。这按构造给出交换图
$$
\xymatrix{
\xi'_{n + 1} \ar[r]_{f'_n} \ar[d] & \xi'_n \ar[d] \\
\xi_{n + 1} \ar[r]^{f_n} & \xi_n
}
$$
此外，环映射 $R'_{n + 1} \to R'_n$ 是满射。
令 $R' = \lim_n R'_n$。则 $R' \to R$ 是单射。

\medskip\noindent
不过，先验上并不清楚 $R'$ 是否诺特。
为证明这一点，利用 $\xi$ 的万有性。
确切地，万有性蕴含 $\widehat{\mathcal{F}}$ 中存在态射
$\xi \to \xi'_n$；见引理
\ref{formal-defos-lemma-versal-object-quasi-initial}。
相应映射 $R \to R'_n$ 必为满射
（因为 $\xi'_n \to \xi_n$ 在 $\mathcal{S}_{\xi_n}$ 中极小）。
因此余切空间的维数有界，而引理
\ref{formal-defos-lemma-limit-artinian} 蕴含 $R'$ 诺特，
即它是 $\widehat{\mathcal{C}}_\Lambda$ 的对象。
由引理 \ref{formal-defos-lemma-formal-objects-different-filtration}
（再加上引理 \ref{formal-defos-lemma-limit-artinian} 关于滤过的结果），
元素序列 $\xi'_n$ 定义 $R'$ 上方的形式对象 $\xi'$，
并且有映射 $\xi' \to \xi$。

\medskip\noindent
按构造，对每个 $n$，(1) 对
$R \to R/\mathfrak m_R^n$ 成立。
因为 (1) 中的每个 $R \to A$ 都经
$R \to R/\mathfrak m_R^n \to A$ 分解，
所以由引理 \ref{formal-defos-lemma-smallest-where-descends} (1)，
$\xi'_n \to \xi_n$ 在
$R'_n \to R/\mathfrak m_R^n$ 上方的极小性蕴含：
(1) 对位于 $f(R) \subset A$ 上方的 $x' \to x$ 成立。% P10-E0462

\medskip\noindent
若 (2) 中的 $R'' \to R'$ 不是满射，
则对某个 $n$，$R'' \to R' \to R'_n$ 不是满射，
于是 $\xi'_n \to \xi_n$ 不会是极小的，矛盾。
这个矛盾证明了 (2)。
\end{proof}

\begin{lemma}
\label{formal-defos-lemma-descends-versal}
设 $\mathcal{F}$ 是 $\mathcal{C}_\Lambda$ 上满足 (S1)
的群胚余纤维范畴。设 $\xi$ 是位于 $R$ 上方的
$\mathcal{F}$ 的万有形式对象。
设 $\xi' \to \xi$ 是引理
\ref{formal-defos-lemma-smallest-where-descends-versal} 所构造的、
位于 $R' \subset R$ 上方的形式对象态射。则
$$
R \cong R'[[x_1, \ldots, x_r]]
$$
是 $R'$ 上的幂级数环。
此外，$\xi'$ 也是万有形式对象。
\end{lemma}

\begin{proof}
由引理 \ref{formal-defos-lemma-versal-object-quasi-initial}，
存在态射 $\xi \to \xi'$。
由引理 \ref{formal-defos-lemma-smallest-where-descends-versal}，
相应映射 $f : R \to R'$ 诱导满射
$f|_{R'} : R' \to R'$。
由《代数》引理
\ref{algebra-lemma-surjective-endo-noetherian-ring-is-iso}，
这是同构。因此 $I = \Ker(f)$ 是 $R$ 的理想，
并且 $R = R' \oplus I$。
选取 $x_1, \ldots, x_n \in I$，使其构成
$I/\mathfrak m_RI$ 的一组基。考虑把 $X_i$ 映到 $x_i$ 的映射
$S = R'[[X_1, \ldots, X_r]] \to R$。% P10-E0463
对每个 $n \geq 1$，得到阿廷 $R'$-代数的满射
$B = S/\mathfrak m_S^n \to R/\mathfrak m_R^n = A$。
用 $y \in \Ob(\mathcal{F}(B)$、相应地
$x \in \Ob(\mathcal{F}(A))$ 表示 $\xi'$ 分别沿
$R' \to S \to B$、相应地沿 $R' \to S \to A$ 的推前。% P10-E0464
注意，$x$ 也是 $\xi$ 沿 $R \to A$ 的推前，
因为 $\xi$ 是 $\xi'$ 沿 $R' \to R$ 的推前。
因此有实线图
$$
\vcenter{
\xymatrix{
& y \ar[d] \\
\xi \ar[r] \ar@{..>}[ru] & x
}
}
\quad\text{位于下图上方}\quad
\vcenter{
\xymatrix{
& S/\mathfrak m_S^n \ar[d] \\
R \ar[r] \ar@{..>}[ru] & R/\mathfrak m_R^n
}
}
$$
因为 $\xi$ 万有，使用说明 \ref{formal-defos-remark-versal-object}，
得到补全这些图的虚线箭头。
特别地，映射
$S/\mathfrak m_S^n \to R/\mathfrak m_R^n$
有截面 $h_n : R/\mathfrak m_R^n \to S/\mathfrak m_S^n$。
由引理 \ref{formal-defos-lemma-power-series}，$S \to R$ 是同构。

\medskip\noindent
由于 $\xi$ 是 $\xi'$ 沿 $R' \to R$ 的推前，
说明 \ref{formal-defos-remark-formal-objects-yoneda-map} 给出交换图
$$
\xymatrix{
\underline{R}|_{\mathcal{C}_\Lambda} \ar[rr] \ar[rd]_{\underline{\xi}} & &
\underline{R'}|_{\mathcal{C}_\Lambda} \ar[ld]^{\underline{\xi'}} \\
& \mathcal{F}
}
$$
因为 $R' \to R$ 有左逆（即 $R \to R/I = R'$），
可知
$\underline{R}|_{\mathcal{C}_\Lambda} \to
\underline{R'}|_{\mathcal{C}_\Lambda}$ 本质满。
因此，由引理 \ref{formal-defos-lemma-smooth-properties}，
$\underline{\xi'}$ 光滑；即 $\xi'$ 是万有形式对象。
\end{proof}

\noindent
受前述引理启发，作如下定义。

\begin{definition}
\label{formal-defos-definition-minimal-versal}
设 $\mathcal{F}$ 是预形变范畴。
若对任意形式对象态射 $\xi' \to \xi$，
底层环映射都是满射，则称 $\mathcal{F}$ 的万有形式对象
$\xi$ 是{\it 极小的}\footnote{这可能不是标准术语。
许多作者把这一概念与切空间的性质联系起来。
我们将在第 \ref{formal-defos-section-miniversal-objects-existence} 节说明这种联系。}。
有时把极小的万有形式对象称为{\it 极小万有的}。
\end{definition}

\noindent
本节的工作说明这个定义是合理的。
首先，万有形式对象的存在蕴含 $\mathcal{F}$ 具有 (S1)。
其次，前述引理表明极小的万有形式对象存在。
最后，任意两个极小的万有形式对象都同构。
下面汇总这些结果（并给出详细证明）。

\begin{lemma}
\label{formal-defos-lemma-minimal-versal}
设 $\mathcal{F}$ 是具有万有形式对象的预形变范畴。则
\begin{enumerate}
\item $\mathcal{F}$ 有极小的万有形式对象；
\item 极小的万有形式对象在同构意义下唯一；并且
\item 任意万有形式对象都是某个极小的万有形式对象
沿幂级数环扩张的推前。
\end{enumerate}
\end{lemma}

\begin{proof}
假设 $\mathcal{F}$ 有位于 $R$ 上方的万有形式对象 $\xi$。
则它满足 (S1)；见引理 \ref{formal-defos-lemma-versal-object-S1}。
设位于 $R' \subset R$ 上方的 $\xi' \to \xi$
是引理 \ref{formal-defos-lemma-smallest-where-descends-versal}
所构造的任一态射。
由引理 \ref{formal-defos-lemma-descends-versal}，$\xi'$ 是万有的；
因此按构造，它是极小的万有形式对象。这证明了 (1)。
此外，$R \cong R'[[x_1, \ldots, x_n]]$，这证明了 (3)。

\medskip\noindent
假设 $\xi_i/R_i$ 是两个极小的万有形式对象。
由引理 \ref{formal-defos-lemma-versal-object-quasi-initial}，
存在态射 $\xi_1 \to \xi_2$ 和 $\xi_2 \to \xi_1$。
由极小性，相应环映射
$f : R_1 \to R_2$ 和 $g : R_2 \to R_1$ 都是满射。
因此，由《代数》引理
\ref{algebra-lemma-surjective-endo-noetherian-ring-is-iso}，
复合 $g \circ f : R_1 \to R_1$ 和
$f \circ g : R_2 \to R_2$ 都是同构。
从而 $f$ 与 $g$ 是同构，于是映射
$\xi_1 \to \xi_2$ 和 $\xi_2 \to \xi_1$ 也是同构
（因为由引理 \ref{formal-defos-lemma-completion-cofibred}，
$\widehat{\mathcal{F}}$ 余纤维于群胚）。
这证明了 (2)，并完成引理的证明。
\end{proof}








\section{极小万有形式对象与切空间}
\label{formal-defos-section-miniversal-objects-existence}

\noindent
定义 \ref{formal-defos-definition-minimal-versal} 所引入的一般极小性概念，
有时可以从切空间上的表现推导出来。
设 $\xi$ 是预形变范畴 $\mathcal{F}$ 的形式对象，并设
$\underline{\xi} : \underline{R}|_{\mathcal{C}_\Lambda} \to \mathcal{F}$
为相应态射。于是可以考虑条件% P10-E0465
\begin{equation}
\label{formal-defos-equation-bijective}
d\underline{\xi} : \text{Der}_\Lambda(R, k) \to T\mathcal{F}
\text{ 是双射}
\end{equation}
以及条件
\begin{equation}
\label{formal-defos-equation-bijective-orbits}
d\underline{\xi} : \text{Der}_\Lambda(R, k) \to T\mathcal{F}
\text{ 在 }\text{Der}_\Lambda(k, k)\text{-轨道上是双射。}
\end{equation}
这里使用了例 \ref{formal-defos-example-tangent-space-prorepresentable-functor}
中的识别
$T\underline{R}|_{\mathcal{C}_\Lambda} = \text{Der}_\Lambda(R, k)$，
以及导子在切空间上的作用（\ref{formal-defos-equation-action}）。
若 $k' \subset k$ 可分，则
$\text{Der}_\Lambda(k, k) = 0$，两个条件等价。
事实表明，在条件 (S2) 成立时，万有形式对象极小，
当且仅当 $\underline{\xi}$ 满足
（\ref{formal-defos-equation-bijective-orbits}）。
此外，若 $\underline{\xi}$ 满足（\ref{formal-defos-equation-bijective}），
则 $\mathcal{F}$ 满足 (S2)。

\begin{lemma}
\label{formal-defos-lemma-miniversal-object-unique}
设 $\mathcal{F}$ 是预形变范畴。
设 $\xi$ 是 $\mathcal{F}$ 的万有形式对象，
且满足（\ref{formal-defos-equation-bijective-orbits}）。
则 $\xi$ 是极小的万有形式对象。
特别地，这样的 $\xi$ 在同构意义下唯一。
\end{lemma}

\begin{proof}
若 $\xi$ 不极小，则存在位于 $R' \to R$ 上方的态射
$\xi' \to \xi$，使得
$R = R'[[x_1, \ldots, x_n]]$ 且 $n > 0$；
见引理 \ref{formal-defos-lemma-minimal-versal}。
因此 $d\underline{\xi}$ 分解为
$$
\text{Der}_\Lambda(R, k) \to
\text{Der}_\Lambda(R', k) \to T\mathcal{F}
$$
而（\ref{formal-defos-equation-bijective-orbits}）不可能成立，
因为
$D : f \mapsto \partial/\partial x_1(f) \bmod \mathfrak m_R$
属于第一个箭头的核，却不在
$\text{Der}_\Lambda(k, k) \to \text{Der}_\Lambda(R, k)$ 的像中。
\end{proof}

\begin{lemma}
\label{formal-defos-lemma-miniversal-object-existence-1}
设 $\mathcal{F}$ 是预形变范畴。
设 $\xi$ 是 $\mathcal{F}$ 的万有形式对象，
且满足（\ref{formal-defos-equation-bijective}）。则
\begin{enumerate}
\item $\mathcal{F}$ 满足 (S1)。
\item $\mathcal{F}$ 满足 (S2)。
\item $\dim_k T\mathcal{F}$ 有限。
\end{enumerate}
\end{lemma}

\begin{proof}
由引理 \ref{formal-defos-lemma-versal-object-S1}，条件 (S1) 成立。
由于 (S1) 成立，(S2) 的第一部分成立。设
$$
\vcenter{
\xymatrix{
y \ar[r]_c \ar[d]_a & x_\epsilon \ar[d]^e \\
x \ar[r]^d          & x_0
}
}
\quad\text{和}\quad
\vcenter{
\xymatrix{
y' \ar[r]_{c'} \ar[d]_{a'} & x_\epsilon \ar[d]^e \\
x \ar[r]^d                 & x_0
}
}
\quad\text{位于下图上方}\quad
\vcenter{
\xymatrix{
A \times_k k[\epsilon] \ar[r] \ar[d] & k[\epsilon] \ar[d] \\
A  \ar[r] & k
}
}
$$
为 (S2) 第二部分中的图。如上，可以找到态射
$b : \xi \to y$ 和 $b' : \xi \to y'$，使图
$$
\xymatrix{
\xi \ar[r]^{b'} \ar[d]_b          & y' \ar[d]^{a'} \\
y \ar[r]^{a} & x
}
$$
交换。用 $p : \mathcal{F} \to \mathcal{C}_\Lambda$
表示结构态射。设 $\widehat{p}(\xi) = R$，
即 $\xi$ 位于
$R \in \Ob(\widehat{\mathcal{C}}_\Lambda)$ 上方。
可见，$\xi$ 经 $p(c) \circ p(b)$ 的推前是 $x_\epsilon$，
而 $\xi$ 经 $p(c') \circ p(b')$ 的推前也是 $x_\epsilon$。
由于 $\xi$ 满足（\ref{formal-defos-equation-bijective}），可知
$p(c) \circ p(b) = p(c') \circ p(b')$
作为 $R \to k[\epsilon]$ 的映射相等。
因此 $p(b) = p(b')$ 作为
$R \to A \times_k k[\epsilon]$ 的映射相等。
于是 $y$ 与 $y'$ 都同构于 $\xi$ 沿此映射的推前，
从而如所需，得到与 $b$、$b'$ 相容的、
位于 $A \times_k k[\epsilon]$ 上方的唯一态射 $y \to y'$。

\medskip\noindent
最后，由例 \ref{formal-defos-example-tangent-space-prorepresentable-functor}，
可知
$\dim_k T\mathcal{F} = \dim_k T\underline{R}|_{\mathcal{C}_\Lambda}$
有限。
\end{proof}

\begin{example}
\label{formal-defos-example-do-not-get-S2}
存在这样的预形变范畴：它们有满足
（\ref{formal-defos-equation-bijective-orbits}）的万有形式对象，
但不满足 (S2)。一个简单例子是取
$F = \underline{k[\epsilon]}/G$，其中
$G \subset \text{Aut}_{\mathcal{C}_\Lambda}(k[\epsilon])$
是有限非平凡子群。确切地，映射
$\underline{k[\epsilon]} \to F$ 光滑，
但 $F$ 的切空间没有自然的 $k$-向量空间结构
（因为它是 $k$-向量空间关于有限群的商）。
\end{example}

\begin{lemma}
\label{formal-defos-lemma-construct-bijective-orbits}
设 $\mathcal{F}$ 是满足 (S2) 且具有万有形式对象的预形变范畴。
则它的极小万有形式对象满足
（\ref{formal-defos-equation-bijective-orbits}）。
\end{lemma}

\begin{proof}
设 $\xi$ 是 $\mathcal{F}$ 的极小万有形式对象；
见引理 \ref{formal-defos-lemma-minimal-versal}。
设 $\xi$ 位于
$R \in \Ob(\widehat{\mathcal{C}}_\Lambda)$ 上方。
为解释（\ref{formal-defos-equation-bijective-orbits}），指出如下事实：
$T\mathcal{F}$ 具有自然的 $k$-向量空间结构
（见引理 \ref{formal-defos-lemma-tangent-space-vector-space}）；
$d\underline{\xi} : \text{Der}_\Lambda(R, k) \to T\mathcal{F}$
是线性的（见引理 \ref{formal-defos-lemma-k-linear-differential}）；
而 $\text{Der}_\Lambda(k, k)$ 的作用由加法给出
（见引理 \ref{formal-defos-lemma-action-linear}）。
考虑图
$$
\xymatrix{
& \Hom_k(\mathfrak m_R/\mathfrak m_R^2, k) \\
K \ar[r] & \text{Der}_\Lambda(R, k) \ar[r]^{d\underline{\xi}} \ar[u] &
T\mathcal{F} \\
& \text{Der}_\Lambda(k, k) \ar[u] \ar[ru]
}
$$
向量空间 $K$ 是 $d\underline{\xi}$ 的核。
注意，中间一列在中项正合，因为它对偶于序列
（\ref{formal-defos-equation-sequence}）。
若（\ref{formal-defos-equation-bijective-orbits}）不成立，
则可以找到非零元素 $D \in K$，
它在 $\Hom_k(\mathfrak m_R/\mathfrak m_R^2, k)$ 中的像不为零。
这意味着存在 $t \in \mathfrak m_R$，使 $D(t) = 1$。% P10-E0466
令 $R' = \{a \in R \mid D(a) = 0\}$。
因为 $D$ 是导子，这是 $R$ 的子环。
由于 $D(t) = 1$，可知 $R' \to k$ 是满射
（比较引理 \ref{formal-defos-lemma-essential-surjection} 的证明）。
注意，$\mathfrak m_{R'} = \Ker(D : \mathfrak m_R \to k)$
是 $R$ 的理想，并且
$\mathfrak m_R^2 \subset \mathfrak m_{R'}$。因此
$$
\mathfrak m_R/\mathfrak m_R^2 =
\mathfrak m_{R'}/\mathfrak m_R^2 + k\overline{t}
$$
这蕴含把 $\epsilon$ 映到 $\overline{t}$ 的映射
$$
R'/\mathfrak m_R^2 \times_k k[\epsilon] \to R/\mathfrak m_R^2
$$
是同构。特别地，存在映射
$R/\mathfrak m_R^2 \to R'/\mathfrak m_R^2$。

\medskip\noindent
设 $\xi \to y$ 是位于
$R \to R/\mathfrak m_R^2$ 上方的态射。
设 $y \to x$ 是位于
$R/\mathfrak m_R^2 \to R'/\mathfrak m_R^2$ 上方的态射。
设 $y \to x_\epsilon$ 是位于
$R/\mathfrak m_R^2 \to k[\epsilon]$ 上方的态射。
设 $x_0$ 为 $\mathcal{F}$ 中位于 $k$ 上方的、
在唯一同构意义下唯一的对象。
由群胚余纤维范畴的公理，得到交换图
$$
\vcenter{
\xymatrix{
y \ar[r] \ar[d] & x_\epsilon \ar[d] \\
x \ar[r]   & x_0
}
}
\quad\text{位于下图上方}\quad
\vcenter{
\xymatrix{
R'/\mathfrak m_R^2 \times_k k[\epsilon] \ar[r] \ar[d] & k[\epsilon] \ar[d] \\
R'/\mathfrak m_R^2 \ar[r] & k.
}
}
$$
因为 $D \in K$，可知 $x_\epsilon$ 同构于
$0 \in \mathcal{F}(k[\epsilon])$；即 $x_\epsilon$ 是
$x_0$ 经 $k \to k[\epsilon], a \mapsto a$ 的推前。
因此，由引理 \ref{formal-defos-lemma-lifting-section}，
存在态射 $x \to y$。由于
$\text{length}_\Lambda(R'/\mathfrak m_R^2) <
\text{length}_\Lambda(R/\mathfrak m_R^2)$，
相应环映射
$R'/\mathfrak m_R^2 \to R/\mathfrak m_R^2$ 不是满射。
这与 $\xi/R$ 的极小性矛盾；见引理
\ref{formal-defos-lemma-smallest-where-descends-versal} 的 (1)。
此矛盾表明这样的 $D$ 不存在，结论得证。% P10-E0467
\end{proof}

\begin{theorem}
\label{formal-defos-theorem-miniversal-object-existence}
设 $\mathcal{F}$ 是预形变范畴。考虑以下条件：
\begin{enumerate}
\item $\mathcal{F}$ 有满足（\ref{formal-defos-equation-bijective}）
的极小万有形式对象；
\item $\mathcal{F}$ 有满足
（\ref{formal-defos-equation-bijective-orbits}）的极小万有形式对象；
\item 以下条件成立：
\begin{enumerate}
\item $\mathcal{F}$ 满足 (S1)。
\item $\mathcal{F}$ 满足 (S2)。
\item $\dim_k T\mathcal{F}$ 有限。
\end{enumerate}
\end{enumerate}
总有
$$
(1) \Rightarrow (3) \Rightarrow (2).
$$
若 $k' \subset k$ 可分，则三个条件等价。
\end{theorem}

\begin{proof}
引理 \ref{formal-defos-lemma-miniversal-object-existence-1}
表明 (1) $\Rightarrow$ (3)。
引理 \ref{formal-defos-lemma-versal-object-existence} 和
\ref{formal-defos-lemma-construct-bijective-orbits}
表明 (3) $\Rightarrow$ (2)。
若 $k' \subset k$ 可分，则
$\text{Der}_\Lambda(k, k) = 0$，从而
（\ref{formal-defos-equation-bijective}）$=$
（\ref{formal-defos-equation-bijective-orbits}）；即 (1) 与 (2) 相同。

\medskip\noindent
在经典情形下，(3) $\Rightarrow$ (1) 的另一种证明是：
在引理 \ref{formal-defos-lemma-versal-object-existence} 的证明中补充几句话，
即可看出此时可直接构造满足
（\ref{formal-defos-equation-bijective}）的万有对象。
这在经典情形下避免了使用引理
\ref{formal-defos-lemma-versal-object-existence}。细节从略。% P10-E0468
\end{proof}

\begin{remark}
\label{formal-defos-remark-compare-schlessinger-H3}
设 $F : \mathcal{C}_\Lambda \to \textit{Sets}$ 是满足
(S1)、(S2) 及 $\dim_k TF < \infty$ 的预形变函子。
回忆，这些条件分别对应于 Schlessinger 论文中的
(H1)、(H2) 与 (H3)；见说明
\ref{formal-defos-remark-compare-schlessinger-H3-pre}。
现在，在经典情形下（或当 $k' \subset k$ 可分时），
依照 Schlessinger 引入包络的概念：所谓{\it 包络}，
是满足如下条件的万有形式对象
$\xi \in \widehat{F}(R)$：
$d\underline{\xi} : T\underline{R}|_{\mathcal{C}_\Lambda} \to TF$
为同构；即（\ref{formal-defos-equation-bijective}）成立。
因此，定理 \ref{formal-defos-theorem-miniversal-object-existence} 表明，
在经典情形下
$$
(H1) + (H2) + (H3)
\Rightarrow
\text{存在包络}
$$
成立。换言之，我们的定理恢复了 Schlessinger
关于包络存在性的定理。
\end{remark}

\begin{remark}
\label{formal-defos-remark-compose-minimal-into-iso-classes}
设 $\mathcal{F}$ 是预形变范畴。
回忆，$\mathcal{F} \to \overline{\mathcal{F}}$ 光滑；
见说明 \ref{formal-defos-remark-smooth-to-iso-classes}。
因此，若
$\xi \in \widehat{\mathcal{F}}(R)$ 是万有形式对象，
则复合
$$
\underline{R}|_{\mathcal{C}_\Lambda} \longrightarrow
\mathcal{F} \longrightarrow \overline{\mathcal{F}}
$$
光滑（引理 \ref{formal-defos-lemma-smooth-properties}），
从而 $\xi$ 在 $\overline{\mathcal{F}}$ 中的像
$\overline{\xi}$ 是万有形式对象。
若（\ref{formal-defos-equation-bijective}）成立，
则 $\overline{\xi}$ 诱导同构
$T\underline{R}|_{\mathcal{C}_\Lambda} \to T\overline{\mathcal{F}}$，
因为 $\mathcal{F} \to \overline{\mathcal{F}}$ 识别切空间。
所以在这种情形下，$\overline{\xi}$ 是
$\overline{\mathcal{F}}$ 的包络；见说明
\ref{formal-defos-remark-compare-schlessinger-H3}。
由定理 \ref{formal-defos-theorem-miniversal-object-existence}，
若 $k' \subset k$ 可分，且 $\mathcal{F}$ 是满足
(S1)、(S2) 与 $\dim_k T\mathcal{F} < \infty$ 的预形变范畴，
则总能找到这样的 $\xi$。
\end{remark}

\begin{example}
\label{formal-defos-example-smooth-continued}
在引理 \ref{formal-defos-lemma-exists-smooth} 中，我们构造了对象
$R \in \widehat{\mathcal{C}}_\Lambda$，使
$\underline{R}|_{\mathcal{C}_\Lambda}$ 光滑，并且
$$
H_1(L_{k/\Lambda}) = \mathfrak m_R/\mathfrak m_R^2
\quad\text{和}\quad
\Omega_{R/\Lambda} \otimes_R k = \Omega_{k/\Lambda}
$$
下面用上述定理重新解释这一点。
确切地，把 $\mathcal{F} = \mathcal{C}_\Lambda$
视为其自身上方的群胚余纤维范畴（使用恒等函子）。
则 $\mathcal{F}$ 是预形变范畴，满足 (S1) 与 (S2)，
并且 $T\mathcal{F} = 0$。
因此，$\mathcal{F}$ 满足定理
\ref{formal-defos-theorem-miniversal-object-existence} 的条件 (3)。
该定理蕴含 (2) 成立；即可以找到某个
$S \in \widehat{\mathcal{C}}_\Lambda$ 上方、满足
（\ref{formal-defos-equation-bijective-orbits}）的极小万有形式对象
$\xi \in \widehat{\mathcal{F}}(S)$。
引理 \ref{formal-defos-lemma-smooth} 表明，
$\Lambda \to S$ 在 $\mathfrak m_S$-进拓扑中形式光滑
（因为
$\underline{\xi} : \underline{R}|_{\mathcal{C}_\Lambda} \to
\mathcal{F} = \mathcal{C}_\Lambda$ 光滑）。% P10-E0469
现在，条件（\ref{formal-defos-equation-bijective-orbits}）表明
$\text{Der}_\Lambda(S, k) \to 0$ 在
$\text{Der}_\Lambda(k, k)$-轨道上是双射。
这意味着单射
$\text{Der}_\Lambda(k, k) \to \text{Der}_\Lambda(S, k)$
也是满射。换言之，
$\Omega_{S/\Lambda} \otimes_S k = \Omega_{k/\Lambda}$。
由于 $\Lambda \to S$ 在 $\mathfrak m_S$-进拓扑中形式光滑，
可以应用《更多代数》引理
\ref{more-algebra-lemma-formally-smooth-JZ}，
得出正合序列（\ref{formal-defos-equation-sequence-extended}）
变成一对识别
$$
H_1(L_{k/\Lambda}) = \mathfrak m_S/\mathfrak m_S^2
\quad\text{和}\quad
\Omega_{S/\Lambda} \otimes_S k = \Omega_{k/\Lambda}
$$
逆向阅读该论证，可知引理 \ref{formal-defos-lemma-exists-smooth}
中构造的 $R$ 承载一个极小万有对象。
由极小万有对象的唯一性
（引理 \ref{formal-defos-lemma-minimal-versal}），
还可得 $R \cong S$；即两个构造给出相同答案。
\end{example}







\section{Rim--Schlessinger 条件与形变范畴}
\label{formal-defos-section-RS-condition}

\noindent
$\mathcal{C}_\Lambda$ 上的群胚余纤维范畴有一个非常自然的性质；
它在实践中容易核验，并且蕴含此前引入的 Schlessinger 性质
(S1) 与 (S2)。

\begin{definition}
\label{formal-defos-definition-RS}
设 $\mathcal{F}$ 是 $\mathcal C_\Lambda$ 上的群胚余纤维范畴。
若对 $\mathcal{F}$ 中的每个图
$$
\vcenter{
\xymatrix{
           & x_2 \ar[d] \\
x_1 \ar[r] & x
}
}
\quad\text{位于下图上方}\quad
\vcenter{
\xymatrix{
           & A_2 \ar[d] \\
A_1 \ar[r] & A
}
}
$$
其中底图属于 $\mathcal{C}_\Lambda$，且 $A_2 \to A$ 满射，
都存在 $\mathcal{F}$ 中的纤维积 $x_1 \times_x x_2$，
使图
$$
\vcenter{
\xymatrix{
x_1 \times_x x_2 \ar[r] \ar[d] & x_2 \ar[d] \\
x_1 \ar[r]      & x
}
}
\quad\text{位于下图上方}\quad
\vcenter{
\xymatrix{
A_1 \times_A A_2 \ar[r] \ar[d] & A_2 \ar[d] \\
A_1 \ar[r]      & A.
}
}
$$
则称 $\mathcal{F}$ 满足{\it 条件 (RS)}。
\end{definition}

\begin{lemma}
\label{formal-defos-lemma-RS-fiber-square}
设 $\mathcal{F}$ 是 $\mathcal{C}_\Lambda$ 上满足 (RS)
的群胚余纤维范畴。给定 $\mathcal{F}$ 中的交换图
$$
\vcenter{
\xymatrix{
y \ar[r] \ar[d] & x_2 \ar[d]   \\
x_1 \ar[r]      & x
}
}
\quad\text{位于下图上方}\quad
\vcenter{
\xymatrix{
A_1 \times_A A_2 \ar[r] \ar[d] & A_2 \ar[d] \\
A_1 \ar[r]      & A.
}
}
$$
若 $A_2 \to A$ 满射，则该图是纤维积方块。
\end{lemma}

\begin{proof}
因为 $\mathcal{F}$ 满足 (RS)，存在纤维积图
$$
\vcenter{
\xymatrix{
x_1 \times_x x_2 \ar[r] \ar[d] & x_2 \ar[d] \\
x_1 \ar[r]      & x
}
}
\quad\text{位于下图上方}\quad
\vcenter{
\xymatrix{
A_1 \times_A A_2 \ar[r] \ar[d] & A_2 \ar[d] \\
A_1 \ar[r]      & A.
}
}
$$
诱导映射 $y \to x_1 \times_x x_2$ 位于
$\text{id} : A_1 \times_A A_1 \to A_1 \times_A A_1$ 上方，% P10-E0470
因而是同构。
\end{proof}

\begin{lemma}
\label{formal-defos-lemma-RS-small-extension}
设 $\mathcal{F}$ 是 $\mathcal C_\Lambda$ 上的群胚余纤维范畴。
若只在 $A_2 \to A$ 是小扩张时假设定义
\ref{formal-defos-definition-RS} 中的条件成立，则 $\mathcal{F}$ 满足 (RS)。
\end{lemma}

\begin{proof}
应用引理 \ref{formal-defos-lemma-factor-small-extension}。
证明类似于引理 \ref{formal-defos-lemma-smoothness-small-extensions} 的证明。
\end{proof}

\begin{lemma}
\label{formal-defos-lemma-RS-2-categorical}
设 $\mathcal{F}$ 是 $\mathcal{C}_\Lambda$ 上的群胚余纤维范畴。
以下条件等价：
\begin{enumerate}
\item $\mathcal{F}$ 满足 (RS)；
\item 函子
$\mathcal{F}(A_1 \times_A A_2) \to
\mathcal{F}(A_1) \times_{\mathcal{F}(A)} \mathcal{F}(A_2)$
（见（\ref{formal-defos-equation-compare}））在 $A_2 \to A$ 满射时
都是范畴等价；% P10-E0471
\item 在 $A_2 \to A$ 是小扩张时，(2) 中的同一陈述成立。
\end{enumerate}
\end{lemma}

\begin{proof}
假设 (1)。由引理 \ref{formal-defos-lemma-RS-fiber-square}，
可知 $\mathcal{F}(A_1 \times_A A_2)$ 的每个对象都形如
$x_1 \times_x x_2$。此外，
$$
\Mor_{A_1 \times_A A_2}(x_1 \times_x x_2, y_1 \times_y y_2) =
\Mor_{A_1}(x_1, y_1) \times_{\Mor_A(x, y)}
\Mor_{A_2}(x_2, y_2).
$$
因此，由《范畴》说明
\ref{categories-remark-other-description-2-fibre-product}，
可知 $\mathcal{F}(A_1 \times_A A_2)$ 是
$\mathcal{F}(A_1)$ 与 $\mathcal{F}(A_2)$ 在
$\mathcal{F}(A)$ 上的 $2$-纤维积。
换言之，(2) 成立。

\medskip\noindent
蕴含 (2) $\Rightarrow$ (3) 是立即的。

\medskip\noindent
假设 (3)。给定 $q_1 : A_1 \to A$ 与 $q_2 : A_2 \to A$，
其中 $q_2$ 是小扩张。
将使用《范畴》说明
\ref{categories-remark-other-description-2-fibre-product}
中对 $2$-纤维积
$\mathcal{F}(A_1) \times_{\mathcal{F}(A)} \mathcal{F}(A_2)$ 的描述。
因此，设 $y \in \mathcal{F}(A_1 \times_A A_2)$ 对应于
$(x_1, x_2, x, a_1 : x_1 \to x, a_2 : x_2 \to x)$。
设 $z$ 是 $\mathcal{F}$ 中位于 $C$ 上方的对象。则
\begin{align*}
\Mor_\mathcal{F}(z, y) & =
\{(f, \alpha) \mid f : C \to A_1 \times_A A_2,
\alpha : f_*z \to y\} \\
& = \{(f_1, f_2, \alpha_1, \alpha_2) \mid
f_i : C \to A_i, \ \alpha_i : f_{i, *}z \to x_i, \\
& \quad\quad\quad\quad
q_1 \circ f_1 = q_2 \circ f_2, \ q_{1, *} \alpha_1 = q_{2, *}\alpha_2\} \\ % P10-E0472
& =
\Mor_\mathcal{F}(z, x_1) \times_{\Mor_\mathcal{F}(z, x)}
\Mor_\mathcal{F}(z, x_2)
\end{align*}
所以 $y$ 是 $x_1$ 与 $x_2$ 在 $x$ 上的纤维积。
由此可见，当 $A_2 \to A$ 是小扩张时，
$\mathcal{F}$ 满足 (RS)。于是由引理
\ref{formal-defos-lemma-RS-small-extension}，(RS) 成立。
\end{proof}

\begin{remark}
\label{formal-defos-remark-compare-schlessinger-H4}
当 $\mathcal{F}$ 余纤维于集合时，条件 (RS) 正是
Schlessinger 论文 \cite[Theorem 2.11]{Sch} 中的条件 (H4)。
确切地，对函子
$F: \mathcal{C}_\Lambda \to \textit{Sets}$，
条件 (RS) 表述为：若 $A_1 \to A$ 与 $A_2 \to A$
是 $\mathcal{C}_\Lambda$ 中的映射，且 $A_2 \to A$ 满射，
则诱导映射
$F(A_1 \times_A A_2) \to F(A_1) \times_{F(A)} F(A_2)$
是双射。
\end{remark}

\begin{lemma}
\label{formal-defos-lemma-RS-implies-S1-S2}
设 $\mathcal{F}$ 是 $\mathcal{C}_\Lambda$ 上的群胚余纤维范畴。
$\mathcal{F}$ 的条件 (RS) 蕴含 $\mathcal{F}$ 的 (S1) 与 (S2)。
\end{lemma}

\begin{proof}
使用引理 \ref{formal-defos-lemma-RS-2-categorical} 的重新表述，
以及定义 \ref{formal-defos-definition-S1-S2} 后对 (S1) 的解释，
立即可知 (RS) 蕴含 (S1)。这证明了 (S2) 的第一部分。
(S2) 的第二部分成立，是因为引理
\ref{formal-defos-lemma-RS-fiber-square} 告诉我们：若 $y,y'$ 如定义
\ref{formal-defos-definition-S1-S2} 中 (S2) 的第二部分所述，则
$y = x_1 \times_{d, x_0, e} x_2 = y'$。
（事实上，态射 $y \to y'$ 同时与 $a,a'$ 及 $c,c'$ 相容！）
\end{proof}

\noindent
下述引理类似于引理 \ref{formal-defos-lemma-S1-S2-associated-functor}。
回忆，若 $\mathcal{F}$ 是 $\mathcal{C}_\Lambda$
上的群胚余纤维范畴，且 $x$ 是 $\mathcal{F}$ 中位于 $A$ 上方的对象，
则记
$\text{Aut}_A(x) = \Mor_A(x, x) = \Mor_{\mathcal{F}(A)}(x, x)$。
若 $x' \to x$ 是 $\mathcal{F}$ 中位于
$A' \to A$ 上方的态射，则有良定义的群映射
$\text{Aut}_{A'}(x') \to \text{Aut}_A(x)$。

\begin{lemma}
\label{formal-defos-lemma-RS-associated-functor}
设 $\mathcal{F}$ 是 $\mathcal{C}_\Lambda$ 上满足 (RS)
的群胚余纤维范畴。以下条件等价：
\begin{enumerate}
\item $\overline{\mathcal{F}}$ 满足 (RS)。
\item 设 $f_1: A_1 \to A$ 和 $f_2: A_2 \to A$
是 $\mathcal{C}_\Lambda$ 中的环映射，且 $f_2$ 满射。
同构类集合的诱导映射
$$
\overline{\mathcal{F}(A_1) \times_{\mathcal{F}(A)} \mathcal{F}(A_2)}
\to \overline{\mathcal{F}}(A_1) \times_{\overline{\mathcal{F}}(A)}
\overline{\mathcal{F}}(A_2)
$$
是单射。
\item 对 $\mathcal{F}$ 中位于满环映射
$A' \to A$ 上方的每个态射 $x' \to x$，映射
$\text{Aut}_{A'}(x') \to \text{Aut}_A(x)$ 是满射。
\item 对 $\mathcal{F}$ 中位于小扩张
$A' \to A$ 上方的每个态射 $x' \to x$，映射
$\text{Aut}_{A'}(x') \to \text{Aut}_A(x)$ 是满射。
\end{enumerate}
\end{lemma}

\begin{proof}
我们证明 (1) 等价于 (2)，而 (2) 等价于 (3)。
(3) 与 (4) 的等价性由引理
\ref{formal-defos-lemma-factor-small-extension} 得出。

\medskip\noindent
设 $f_1: A_1 \to A$ 和 $f_2: A_2 \to A$
是 $\mathcal{C}_\Lambda$ 中的环映射，且 $f_2$ 满射。
由说明 \ref{formal-defos-remark-compare-schlessinger-H4}，
$\overline{\mathcal{F}}$ 满足 (RS)，当且仅当对任意这样的
$f_1,f_2$，映射
$$
\overline{\mathcal{F}}(A_1 \times_A A_2) \to \overline{\mathcal
F}(A_1) \times_{\overline{\mathcal{F}}(A)} \overline{\mathcal{F}}(A_2)
$$
都是双射。此映射至少是满射，因为这正是 (S1) 的条件，% P10-E0473
而由引理 \ref{formal-defos-lemma-RS-implies-S1-S2} 和
\ref{formal-defos-lemma-S1-S2-associated-functor}，
$\overline{\mathcal{F}}$ 满足 (S1)。
此外，该映射分解为
$$
\overline{\mathcal{F}}(A_1 \times_A A_2)
\longrightarrow
\overline{\mathcal{F}(A_1) \times_{\mathcal{F}(A)} \mathcal{F}(A_2)}
\longrightarrow
\overline{\mathcal{F}}(A_1) \times_{\overline{\mathcal{F}}(A)}
\overline{\mathcal{F}}(A_2),
$$
其中第一个映射是双射，因为
$$
\mathcal{F}(A_1 \times_A A_2)
\longrightarrow
\mathcal{F}(A_1) \times_{\mathcal{F}(A)} \mathcal{F}(A_2)
$$
由 $\mathcal{F}$ 的 (RS) 是等价。因此 (1) 等价于 (2)。

\medskip\noindent
假设 (2) 成立。设 $x' \to x$ 是 $\mathcal{F}$ 中
位于满环映射 $f: A' \to A$ 上方的态射。
设 $a \in \text{Aut}_A(x)$。对象
$$
(x', x', a : x \to x), \ (x', x', \text{id} : x \to x)
$$
属于
$\mathcal{F}(A') \times_{\mathcal{F}(A)} \mathcal{F}(A')$，
并在
$\overline{\mathcal{F}}(A') \times_{\overline{\mathcal{F}}(A)}
\overline{\mathcal{F}}(A')$ 中有相同的像。
由 (2)，存在映射 $b_1,b_2 : x' \to x'$，% P10-E0474
使图
$$
\xymatrix{
x \ar[r]_a \ar[d]_{f_*b_1} & x \ar[d]^{f_*b_2} \\
x \ar[r]^{\text{id}} & x
}
$$
交换。因此，
$b_2^{-1} \circ b_1 \in \text{Aut}_{A'}(x')$
的像是 $a \in \text{Aut}_A(x)$。所以 (3) 成立。

\medskip\noindent
假设 (3) 成立。设
$$
(x_1, x_2, a : (f_1)_*x_1 \to (f_2)_*x_2),
\ (x'_1, x'_2, a' : (f_1)_*x'_1 \to (f_2)_*x'_2)
$$
是
$\mathcal{F}(A_1) \times_{\mathcal{F}(A)} \mathcal{F}(A_2)$
中在
$\overline{\mathcal{F}}(A_1) \times_{\overline{\mathcal{F}}(A)}
\overline{\mathcal{F}}(A_2)$ 中具有相同像的对象。
则存在 $\mathcal{F}(A_1)$ 中的态射
$b_1: x_1 \to x'_1$，以及 $\mathcal F(A_2)$ 中的态射
$b_2: x_2 \to x'_2$。
由 (3)，可以用 $x_2$ 在 $A_2$ 上的自同构修改 $b_2$，
使图
$$
\xymatrix{
(f_1)_*x_1 \ar[r]_a \ar[d]_{(f_1)_*b_1} & (f_2)_*x_2 \ar[d]^{(f_2)_*b_2} \\
(f_1)_*x'_1 \ar[r]^{a'} & (f_2)_*x'_2.
}
$$
交换。这证明
$(x_1, x_2, a) \cong (x'_1, x'_2, a')$ 在
$\overline{\mathcal{F}(A_1) \times_{\mathcal{F}(A)} \mathcal{F}(A_2)}$
中成立。% P10-E0475
所以 (2) 成立。
\end{proof}

\noindent
最后，我们定义形变范畴的概念。

\begin{definition}
\label{formal-defos-definition-deformation-category}
满足 (RS) 的前形变范畴 $\mathcal{F}$ 称为一个{\it 形变范畴}。
形变范畴之间的态射，是它们作为 $\mathcal{C}_\Lambda$ 上范畴的态射。
\end{definition}

\begin{remark}
\label{formal-defos-remark-deformation-functor}
若函子 $F: \mathcal{C}_\Lambda \to \textit{Sets}$ 所对应的
余纤维集合是形变范畴，即 $F(k)$ 是单元素集合且 $F$ 满足 (RS)，
则称这样的函子为一个{\it 形变函子}。若 $\mathcal{F}$ 是形变范畴，
则 $\overline{\mathcal{F}}$ 是前形变函子，但未必是形变函子，
正如引理 \ref{formal-defos-lemma-RS-associated-functor} 所示。
\end{remark}

\begin{example}
\label{formal-defos-example-prorepresentable-deformation-functor}
pro-可表函子 $F$ 是形变函子。事实上，设
$R \in \Ob(\widehat{\mathcal{C}}_\Lambda)$ 且
$F(A) = \Mor_{\widehat{\mathcal{C}}_\Lambda}(R, A)$。
态射 $R \to k$ 唯一，所以 $F(k)$ 是单元素集合。由于
$$
\Hom_\Lambda(R, A_1 \times_A A_2) =
\Hom_\Lambda(R, A_1) \times_{\Hom_\Lambda(R, A)}
\Hom_\Lambda(R, A_2),
$$
对 $\widehat{\mathcal{C}}_\Lambda$ 中的映射也有同样的结论，
因而 $F$ 满足 (RS)。
\end{example}

\noindent
下面是我们关于以下过程的典型说明之一：从一个群胚余纤维范畴
转到其在 $k$ 上某一点处的前形变范畴。它说明该过程保持 (RS)。

\begin{lemma}
\label{formal-defos-lemma-localize-RS}
设 $\mathcal{F}$ 是 $\mathcal{C}_\Lambda$ 上的群胚余纤维范畴，
且 $x_0 \in \Ob(\mathcal{F}(k))$。设 $\mathcal{F}_{x_0}$ 是说明
\ref{formal-defos-remark-localize-cofibered-groupoid} 所构造的
$\mathcal{C}_\Lambda$ 上的群胚余纤维范畴。
若 $\mathcal{F}$ 满足 (RS)，则 $\mathcal{F}_{x_0}$ 也满足 (RS)。
特别地，$\mathcal{F}_{x_0}$ 是形变范畴。
\end{lemma}

\begin{proof}
$\mathcal{F}_{x_0}$ 中任何如定义 \ref{formal-defos-definition-RS} 所述的图表
都会给出 $\mathcal{F}$ 中的一个图表；在 $\mathcal{F}$ 中
对后一个图表应用 (RS) 所得的对象，
也可视为对 $\mathcal{F}_{x_0}$ 中原图表应用 (RS) 所得的对象。
\end{proof}

\noindent
下述引理类似于代数栈的 $2$-纤维积仍为代数栈这一事实。

\begin{lemma}
\label{formal-defos-lemma-RS-fiber-product-morphisms}
设
$$
\xymatrix{
\mathcal{H} \times_\mathcal{F} \mathcal{G} \ar[r] \ar[d] &
\mathcal{G} \ar[d]^g \\
\mathcal{H} \ar[r]^f & \mathcal{F}
}
$$
是 $\mathcal{C}_\Lambda$ 上群胚余纤维范畴的 $2$-纤维积。% P10-E0476
若 $\mathcal{F}, \mathcal{G}, \mathcal{H}$ 都满足 (RS)，
则 $\mathcal{H} \times_\mathcal{F} \mathcal{G}$ 满足 (RS)。
\end{lemma}

\begin{proof}
若 $A$ 是 $\mathcal{C}_\Lambda$ 的对象，则
$\mathcal{H} \times_\mathcal{F} \mathcal{G}$ 在 $A$ 上的纤维范畴中的对象
是三元组 $(u, v, a)$，其中 $u \in \mathcal{H}(A)$、
$v \in \mathcal{G}(A)$，且 $a : f(u) \to g(v)$ 是
$\mathcal{F}(A)$ 中的态射。考虑
$\mathcal{H} \times_\mathcal{F} \mathcal{G}$ 中的图表
$$
\vcenter{
\xymatrix{
           & (u_2, v_2, a_2) \ar[d] \\
(u_1, v_1, a_1) \ar[r] & (u, v, a)
}
}
\quad\text{位于}\quad
\vcenter{
\xymatrix{
           & A_2 \ar[d] \\
A_1 \ar[r] & A
}
}
$$
之上，其中后一个图表在 $\mathcal{C}_\Lambda$ 中且
$A_2 \to A$ 满射。由于 $\mathcal{H}$ 和 $\mathcal{G}$ 满足 (RS)，
存在位于 $A_1 \times_A A_2$ 上的纤维积
$u_1 \times_u u_2$ 与 $v_1 \times_v v_2$。由于
$\mathcal{F}$ 满足 (RS)，引理 \ref{formal-defos-lemma-RS-fiber-square} 表明
$$
\vcenter{
\xymatrix{
f(u_1 \times_u u_2) \ar[r] \ar[d] & f(u_2) \ar[d] \\
f(u_1) \ar[r] & f(u)
}
}
\quad\text{与}\quad
\vcenter{
\xymatrix{
g(v_1 \times_v v_2) \ar[r] \ar[d] & g(v_2) \ar[d] \\
g(v_1) \ar[r] & g(v)
}
}
$$
在 $\mathcal{F}$ 中都是纤维积方块。因此可将
$a_1 \times_a a_2$ 视为从 $f(u_1 \times_u u_2)$ 到
$g(v_1 \times_v v_2)$ 的、位于 $A_1 \times_A A_2$ 上的态射。
由此可知
$$
\xymatrix{
 (u_1 \times_u u_2, v_1 \times_v v_2, a_{1} \times_a a_2) \ar[d] \ar[r] &
(u_2, v_2, a_2) \ar[d] \\
(u_1, v_1, a_1) \ar[r] & (u, v, a)
}
$$
是 $\mathcal{H} \times_\mathcal{F} \mathcal{G}$ 中的纤维积方块，
这正是所需结论。
\end{proof}

\section{对象的提升}
\label{formal-defos-section-lifts}

\noindent
本节的内容是：在形变范畴的情形，切空间在沿小扩张的提升集合上
有一个单传递作用。

\begin{definition}
\label{formal-defos-definition-lifts}
设 $\mathcal{F}$ 是 $\mathcal{C}_\Lambda$ 上的群胚余纤维范畴。
设 $f: A' \to A$ 是 $\mathcal{C}_\Lambda$ 中的映射，并设
$x \in \mathcal{F}(A)$。$x$ 沿 $f$ 的提升范畴
$\textit{Lift}(x, f)$ 具有如下对象与态射。
\begin{enumerate}
\item 对象：$x$ 沿 $f$ 的一个{\it 提升}，是位于 $f$ 上的态射
$x' \to x$。
\item 态射：从提升 $a_1 : x'_1 \to x$ 到提升
$a_2 : x'_2 \to x$ 的一个{\it 提升态射}，是
$\mathcal{F}(A')$ 中满足 $a_2 = a_1 \circ b$ 的态射
$b : x'_1 \to x'_2$。% P10-E0477
\end{enumerate}
$x$ 沿 $f$ 的提升集合 $\text{Lift}(x, f)$，是
$\textit{Lift}(x, f)$ 的同构类集合。
\end{definition}

\begin{remark}
\label{formal-defos-remark-omit-arrow}
当映射 $f: A' \to A$ 由语境可知时，可以分别用
$\textit{Lift}(x, A')$ 和 $\text{Lift}(x, A')$ 代替
$\textit{Lift}(x, f)$ 和 $\text{Lift}(x, f)$。
\end{remark}

\begin{remark}
\label{formal-defos-remark-tangent-space-lifting}
设 $\mathcal{F}$ 是 $\mathcal{C}_\Lambda$ 上的群胚余纤维范畴，
且 $x_0 \in \Ob(\mathcal{F}(k))$。设 $V$ 是有限维向量空间。
则 $\text{Lift}(x_0, k[V])$ 是 $\mathcal{F}_{x_0}(k[V])$
的同构类集合，其中 $\mathcal{F}_{x_0}$ 是 $\mathcal{F}$ 中
位于 $x_0$ 上的对象所成的前形变范畴；见说明
\ref{formal-defos-remark-localize-cofibered-groupoid}。因此，若 $\mathcal{F}$
满足 (S2)，则 $\mathcal{F}_{x_0}$ 也满足 (S2)（见引理
\ref{formal-defos-lemma-S1-S2-localize}），并且由引理
\ref{formal-defos-lemma-tangent-space-vector-space} 可知
$$
\text{Lift}(x_0, k[V]) = T\mathcal{F}_{x_0} \otimes_k V
$$
是 $k$-向量空间之间的等式。
\end{remark}

\begin{remark}
\label{formal-defos-remark-lift-bijections}
设 $\mathcal{F}$ 是 $\mathcal
C_\Lambda$ 上满足 (RS) 的群胚余纤维范畴。设
$$
\xymatrix{
A_1 \times_A A_2 \ar[r] \ar[d] & A_2 \ar[d] \\
A_1 \ar[r] & A
}
$$
是 $\mathcal{C}_\Lambda$ 中的纤维积方块，并且
$A_1 \to A$ 或 $A_2 \to A$ 至少有一个满射。设
$x \in \Ob(\mathcal{F}(A))$。给定 $x$ 到 $A_1$ 和 $A_2$
的提升 $x_1 \to x$ 与 $x_2 \to x$，由 (RS) 得到
$x$ 到 $A_1 \times_A A_2$ 的提升
$x_1 \times_x x_2 \to x$。反过来，由引理
\ref{formal-defos-lemma-RS-fiber-square}，$x$ 到 $A_1 \times_A A_2$
的每个提升都具有这种形式。从而得到双射% P10-E0478
$$
\text{Lift}(x, A_1) \times \text{Lift}(x, A_2)
\longrightarrow
\text{Lift}(x, A_1 \times_A A_2).
$$
类似地，若 $x_1 \to x$ 是 $x$ 到 $A_1$ 的一个固定提升，
则有双射
$$
\text{Lift}(x_1, A_1 \times_A A_2)
\longrightarrow
\text{Lift}(x, A_2).
$$
现设
$$
\xymatrix{
A_1' \times_A A_2 \ar[r] \ar[d] & A_1 \times_A A_2 \ar[r] \ar[d] & A_2
\ar[d] \\
A_1' \ar[r] & A_1 \ar[r] & A
}
$$
是 $\mathcal{C}_\Lambda$ 中两个纤维积方块的复合，且
$A'_1 \to A_1$ 与 $A_1 \to A$ 都满射。设 $x_1 \to x$
是位于 $A_1 \to A$ 上的态射。由上可得双射
\begin{align*}
\text{Lift}(x_1, A_1' \times_A A_2)
& = \text{Lift}(x_1, A_1') \times \text{Lift}(x_1, A_1 \times_A A_2) \\
& = \text{Lift}(x_1, A_1') \times \text{Lift}(x, A_2).
\end{align*}
\end{remark}

\begin{lemma}
\label{formal-defos-lemma-free-transitive-action}
设 $\mathcal{F}$ 是形变范畴。设 $A' \to A$ 是
$\mathcal{C}_\Lambda$ 中的满环映射，其核 $I$ 被
$\mathfrak m_{A'}$ 零化。设 $x \in \Ob(\mathcal{F}(A))$。
若 $\text{Lift}(x, A')$ 非空，则
$T\mathcal{F} \otimes_k I$ 在 $\text{Lift}(x, A')$
上有一个自由且传递的作用。
\end{lemma}

\begin{proof}
考虑环映射 $g : A' \times_A A' \to k[I]$，其定义规则为
$g(a_1, a_2) = \overline{a_1} \oplus a_2 - a_1$（与引理
\ref{formal-defos-lemma-lifting-along-small-extension} 比较）。存在同构
$$
A' \times_A A' \xrightarrow{\sim} A' \times_k k[I]
$$
由
$(a_1, a_2) \mapsto (a_1, g(a_1, a_2))$
给出。
此同构与到第一因子 $A'$ 的投影相容，因而也与
$A' \times_A A'$ 和 $A' \times_k k[I]$ 到 $A$ 的投影相容。
所以有双射
\begin{equation}
\label{formal-defos-equation-one}
\text{Lift}(x, A' \times_A A')
\longrightarrow
\text{Lift}(x, A' \times_k k[I])
\end{equation}
由说明 \ref{formal-defos-remark-lift-bijections}，有双射
\begin{equation}
\label{formal-defos-equation-two}
\text{Lift}(x, A') \times \text{Lift}(x, A')
\longrightarrow
\text{Lift}(x, A' \times_A A')
\end{equation}
有交换图
$$
\xymatrix{
A' \times_k k[I] \ar[r] \ar[d] & A \times_k k[I] \ar[r] \ar[d] & k[I] \ar[d] \\
A' \ar[r] & A \ar[r] & k.
}
$$
因此，若选定 $x$ 沿 $A \to k$ 的一个推前 $x \to x_0$，
则由说明 \ref{formal-defos-remark-lift-bijections} 的最后部分得到双射
\begin{equation}
\label{formal-defos-equation-three}
\text{Lift}(x, A' \times_k k[I])
\longrightarrow
\text{Lift}(x, A') \times \text{Lift}(x_0, k[I])
\end{equation}
依次复合（\ref{formal-defos-equation-two}）、（\ref{formal-defos-equation-one}）和
（\ref{formal-defos-equation-three}），得到双射
$$
\Phi :
\text{Lift}(x, A') \times \text{Lift}(x, A')
\longrightarrow
\text{Lift}(x, A') \times \text{Lift}(x_0, k[I]).
$$
此双射与到第一因子的投影相容。由说明
\ref{formal-defos-remark-tangent-space-lifting} 可知
$\text{Lift}(x_0, k[I]) = T\mathcal{F} \otimes_k I$。
若 $\text{pr}_2$ 是
$\text{Lift}(x, A') \times \text{Lift}(x, A')$ 的第二投影，
则得到映射
$$
a = \text{pr}_2 \circ \Phi^{-1} :
\text{Lift}(x, A') \times (T\mathcal{F} \otimes_k I)
\longrightarrow
\text{Lift}(x, A').
$$
展开上述所有构造可知，$a(x' \to x, \theta)$ 是唯一的提升
$x'' \to x$，使得
$g_*(x', x'') = \theta$ 在
$\text{Lift}(x_0, k[I]) = T\mathcal{F} \otimes_k I$ 中成立。
为证明这是 $T\mathcal{F} \otimes_k I$ 在
$\text{Lift}(x, A')$ 上的作用，必须证明如下结论：
若 $x', x'', x'''$ 是 $x$ 的提升，并且
$g_*(x', x'') = \theta$、$g_*(x'', x''') = \theta'$，则
$g_*(x', x''') = \theta + \theta'$。这由交换图
$$
\xymatrix{
A' \times_A A' \times_A A'
\ar[rrrrr]_-{(a_1, a_2, a_3) \mapsto (g(a_1, a_2), g(a_2, a_3))}
\ar[rrrrrd]_{(a_1, a_2, a_3) \mapsto g(a_1, a_3)} & & & & &
k[I] \times_k k[I] = k[I \times I] \ar[d]^{+} \\
& & & & & k[I]
}
$$
得出。由于 $\Phi$ 是双射，此作用自由且传递。
\end{proof}

\begin{remark}
\label{formal-defos-remark-free-transitive-action-functorial}
引理 \ref{formal-defos-lemma-free-transitive-action} 的作用具有函子性。
设 $\varphi : \mathcal{F} \to \mathcal{G}$ 是形变范畴的态射。
设 $A' \to A$ 是满环映射，其核 $I$ 被
$\mathfrak m_{A'}$ 零化。设 $x \in \Ob(\mathcal{F}(A))$。
在此情形，$\varphi$ 诱导下列交换图中的竖直箭头：
$$
\xymatrix{
\text{Lift}(x, A') \times (T\mathcal{F} \otimes_k I)
\ar[d]_{(\varphi, d\varphi \otimes \text{id}_I)} \ar[r] &
\text{Lift}(x, A') \ar[d]^\varphi \\
\text{Lift}(\varphi(x), A') \times (T\mathcal{G} \otimes_k I) \ar[r] &
\text{Lift}(\varphi(x), A')
}
$$
该图交换，是因为映射（\ref{formal-defos-equation-two}）、
（\ref{formal-defos-equation-one}）和（\ref{formal-defos-equation-three}）中的每一个
（见引理 \ref{formal-defos-lemma-free-transitive-action} 的证明）
都给出一个类似的交换图。
\end{remark}

\section{Schlessinger 关于 pro-可表函子的定理}
\label{formal-defos-section-schlessingers-theorem}

\noindent
下面推出 Schlessinger 刻画 $\mathcal{C}_\Lambda$ 上
pro-可表函子的定理。

\begin{lemma}
\label{formal-defos-lemma-minimal-smooth-morphism-functors}
设 $F, G: \mathcal{C}_\Lambda \to \textit{Sets}$ 是形变函子。
设 $\varphi : F \to G$ 是光滑态射，并且诱导切空间的同构
$d\varphi : TF \to TG$。则 $\varphi$ 是同构。
\end{lemma}

\begin{proof}
对 $\text{length}_A(A)$ 作归纳，证明对所有
$A \in \Ob(\mathcal{C}_\Lambda)$，映射 $F(A) \to G(A)$
都是双射。当 $A = k$ 时，该陈述由 $F$ 和 $G$ 是形变函子的假设得出。
假设该陈述对长度小于 $n$ 的环成立，并设 $A'$ 是长度为 $n$ 的环。
选取小扩张 $f : A' \to A$。有交换图
$$
\xymatrix{
F(A') \ar[r] \ar[d]_{F(f)} & G(A') \ar[d]^{G(f)} \\
F(A) \ar[r]^{\sim} & G(A)
}
$$
其中映射 $F(A) \to G(A)$ 是双射。由 $F
\to G$ 的光滑性，$F(A') \to G(A')$ 是满射（引理
\ref{formal-defos-lemma-smooth-morphism-essentially-surjective}）。因此，只需对满足
$F(f)^{-1}(x)$ 非空的 $x \in F(A)$，在纤维
$F(f)^{-1}(x) \to G(f)^{-1}(\varphi(x))$ 上检验双射性。
这些纤维恰为 $\text{Lift}(x, A')$ 和
$\text{Lift}(\varphi(x), A')$，并且由假设有同构
$d\varphi \otimes \text{id} :
TF \otimes_k \Ker(f)  \to TG \otimes_k \Ker(f)$。
所以，由引理 \ref{formal-defos-lemma-free-transitive-action} 和说明
\ref{formal-defos-remark-free-transitive-action-functorial}，对满足
$F(f)^{-1}(x)$ 非空的 $x \in F(A)$，映射
$F(f)^{-1}(x) \to G(f)^{-1}(\varphi(x))$ 与
$TF \otimes_k \Ker(f)$ 的自由传递作用相容。因此它是双射。
\end{proof}

\noindent
注意，当 $k' \subset k$ 可分时，下述定理中的条件 (c) 是空条件。

\begin{theorem}
\label{formal-defos-theorem-Schlessinger-prorepresentability}
设 $F: \mathcal{C}_\Lambda \to \textit{Sets}$ 是函子。
则 $F$ pro-可表，当且仅当：(a) $F$ 是形变函子；
(b) $\dim_k TF$ 有限；以及
(c) $\gamma : \text{Der}_\Lambda(k, k) \to TF$ 是单射。
\end{theorem}

\begin{proof}
假设 $F$ 由 $R \in \widehat{\mathcal{C}}_\Lambda$ pro-表示。
由例 \ref{formal-defos-example-prorepresentable-deformation-functor} 可知
$F$ 是形变函子。由例
\ref{formal-defos-example-tangent-space-prorepresentable-functor} 可知
$\dim_k TF$ 有限。最后，由例
\ref{formal-defos-example-tangent-space-map-prorepresentable-functor}，
$\text{Der}_\Lambda(k, k) \to TF$ 等同于
$\text{Der}_\Lambda(k, k) \to \text{Der}_\Lambda(R, k)$；
后者是单射，因为 $R \to k$ 满射。

\medskip\noindent
反过来，假设 (a)、(b)、(c) 成立。由引理
\ref{formal-defos-lemma-RS-implies-S1-S2} 可知 (S1) 与 (S2) 成立。
所以由定理 \ref{formal-defos-theorem-miniversal-object-existence}，
存在 $F$ 的极小万有形式对象 $\xi$，使得
（\ref{formal-defos-equation-bijective-orbits}）成立。设 $\xi$ 位于 $R$ 上。
映射
$$
d\underline{\xi} : \text{Der}_\Lambda(R, k) \to T\mathcal{F}
$$
在 $\text{Der}_\Lambda(k, k)$-轨道上为双射。% P10-E0479
由 (c) 和引理 \ref{formal-defos-lemma-action-linear}，
$\text{Der}_\Lambda(k, k)$ 在左端的作用是自由的，
故该映射是双射。因此，由引理
\ref{formal-defos-lemma-minimal-smooth-morphism-functors}，
$\underline{\xi}$ 是同构。
\end{proof}

\section{无穷小自同构}
\label{formal-defos-section-infinitesimal-automorphisms}

\noindent
设 $\mathcal{F}$ 是 $\mathcal{C}_\Lambda$ 上的群胚余纤维范畴。
给定 $\mathcal{F}$ 中位于 $A' \to A$ 上的态射 $x' \to x$，
会诱导同态
$$
\text{Aut}_{A'}(x') \to \text{Aut}_A(x).
$$
引理 \ref{formal-defos-lemma-RS-associated-functor} 表明，此同态的余核决定
$\mathcal{F}$ 的条件 (RS) 是否能传递给
$\overline{\mathcal{F}}$。本节研究此同态的核。我们将看到，
它也衡量 $\mathcal{F}$ 与 $\overline{\mathcal{F}}$ 之间的差距。

\begin{definition}
\label{formal-defos-definition-relative-infinitesimal-auts}
设 $\mathcal{F}$ 是 $\mathcal
C_\Lambda$ 上的群胚余纤维范畴。设 $x' \to x$ 是
$\mathcal{F}$ 中位于 $A' \to A$ 上的态射。核
$$
\text{Inf}(x'/x) = \Ker(\text{Aut}_{A'}(x') \to \text{Aut}_A(x))
$$
称为 $x'$ 在 $x$ 上的{\it 无穷小自同构群}。
\end{definition}

\begin{definition}
\label{formal-defos-definition-infinitesimal-auts}
设 $\mathcal{F}$ 是 $\mathcal
C_\Lambda$ 上的群胚余纤维范畴。设
$x_0 \in \Ob(\mathcal{F}(k))$。假定已经选定 $x_0$ 沿映射
$k \to k[\epsilon], a \mapsto a$ 的一个推前 $x_0 \to x_0'$。
则存在唯一映射 $x'_0 \to x_0$，使得
$x_0 \to x_0' \to x_0$ 是 $x_0$ 上的恒等映射。于是
$$
\text{Inf}_{x_0}(\mathcal F) = \text{Inf}(x'_0/x_0)
$$
称为 $x_0$ 的{\it 无穷小自同构群}。% P10-E0480
\end{definition}

\begin{remark}
\label{formal-defos-remark-choice-pushforward-immaterial-infinitesimal-aut}
在典范同构的意义下，$\text{Inf}_{x_0}(\mathcal{F})$
与推前 $x_0 \to x_0'$ 的选择无关，因为任意两个推前典范同构。
此外，若 $y_0 \in \mathcal{F}(k)$ 且
$x_0 \cong y_0$ 在 $\mathcal{F}(k)$ 中成立，则
$\text{Inf}_{x_0}(\mathcal{F}) \cong \text{Inf}_{y_0}(\mathcal{F})$；
该同构（仅）依赖于同构 $x_0 \to y_0$ 的选择。
特别地，$\text{Aut}_k(x_0)$ 作用在
$\text{Inf}_{x_0}(\mathcal{F})$ 上。
\end{remark}

\begin{remark}
\label{formal-defos-remark-trivial-aut-point}
假设 $\mathcal{F}$ 是前形变范畴。则
\begin{enumerate}
\item 对 $x_0 \in \Ob(\mathcal{F}(k))$，自同构群
$\text{Aut}_k(x_0)$ 平凡，因而
$\text{Inf}_{x_0}(\mathcal{F}) = \text{Aut}_{k[\epsilon]}(x'_0)$；并且
\item 对 $x_0, y_0 \in \Ob(\mathcal{F}(k))$，存在唯一同构
$x_0 \to y_0$，因而有典范识别
$\text{Inf}_{x_0}(\mathcal{F}) = \text{Inf}_{y_0}(\mathcal{F})$。
\end{enumerate}
由于 $\mathcal{F}(k)$ 非空，选取
$x_0 \in \Ob(\mathcal{F}(k))$ 并置
$$
\text{Inf}(\mathcal{F}) = \text{Inf}_{x_0}(\mathcal{F}),
$$
便得到良定义的 $\mathcal{F}$ 的{\it 无穷小自同构群}。
用此记号，有
$\text{Inf}(\mathcal{F}_{x_0}) = \text{Inf}_{x_0}(\mathcal{F})$。
请与说明 \ref{formal-defos-remark-tangent-space-cofibered-groupoid} 中的等式
$T\mathcal{F}_{x_0} = T_{x_0}\mathcal{F}$ 比较。
\end{remark}

\noindent
我们将看到，当 $\mathcal{F}$ 满足 (RS) 时，
$\text{Inf}_{x_0}(\mathcal{F})$ 具有自然的 $k$-向量空间结构。
同时还将看到，若 $\mathcal{F}$ 满足 (RS)，则位于小扩张上方的态射
$x' \to x$ 的无穷小自同构 $\text{Inf}(x'/x)$ 由
$\text{Inf}_{x_0}(\mathcal{F})$ 控制，其中 $x_0$ 是 $x$
到 $\mathcal{F}(k)$ 的一个推前。为此，对任意对象
$x \in \Ob(\mathcal{F})$ 引入如下自同构函子。

\begin{definition}
\label{formal-defos-definition-automorphism-functor}
设 $p : \mathcal{F} \to \mathcal{C}$ 是任意基范畴
$\mathcal{C}$ 上的群胚余纤维范畴，并假定已经选定推前。
设 $x \in \Ob(\mathcal{F})$，并令 $U = p(x)$。
以 $U/\mathcal{C}$ 表示 $U$ 下的对象范畴。
$x$ 的{\it 自同构函子}是函子
$\mathit{Aut}(x) : U/\mathcal{C} \to \textit{Sets}$，
它把对象 $f : U \to V$ 送到 $\text{Aut}_V(f_*x)$，
并把态射
$$
\xymatrix{
V' \ar[rr] &                    & V\\
          & U \ar[ul]^{f'}  \ar[ur]_f &
}
$$
送到同态
$\text{Aut}_{V'}(f'_*x) \to \text{Aut}_V(f_*x)$；该同态来自
唯一的态射 $f'_*x \to f_*x$，它位于 $V' \to V$ 上，
并与 $x \to f'_*x$ 及 $x \to f_*x$ 相容。
\end{definition}

\noindent
我们将研究 $\mathcal{C}_\Lambda$ 上群胚余纤维范畴
$\mathcal{F}$ 中对象的自同构函子。若
$A \in \Ob(\mathcal{C}_\Lambda)$，则范畴
$A/\mathcal{C}_\Lambda$ 正是范畴 $\mathcal{C}_A$，
即第 \ref{formal-defos-section-CLambda} 节所定义的范畴，其中取
$\Lambda = A$ 且 $k = A/\mathfrak m_A$。因此，对象
$x \in \Ob(\mathcal{F}(A))$ 的自同构函子是函子
$\mathit{Aut}(x) : \mathcal{C}_A \to \textit{Sets}$。

\medskip\noindent
下述引理也可从引理 \ref{formal-defos-lemma-RS-fiber-product-morphisms} 推出：
只需考虑群胚余纤维范畴的“惯性”；例如见《栈》第
\ref{stacks-section-the-inertia-stack} 节和《范畴》第
\ref{categories-section-inertia} 节。不过，直接证明更容易。

\begin{lemma}
\label{formal-defos-lemma-Aut-functor-RS}
设 $\mathcal{F}$ 是 $\mathcal{C}_\Lambda$ 上满足 (RS)
的群胚余纤维范畴，并设 $x \in \Ob(\mathcal{F}(A))$。
则 $\mathit{Aut}(x): \mathcal{C}_A \to \textit{Sets}$ 满足 (RS)。
\end{lemma}

\begin{proof}
由引理 \ref{formal-defos-lemma-RS-2-categorical} 中函子
$\mathcal{F}(A_1 \times_A A_2) \to
\mathcal{F}(A_1) \times_{\mathcal{F}(A)} \mathcal{F}(A_2)$
的全忠实性可知，$\mathit{Aut}(x)$ 满足 (RS)。% P10-E0481
\end{proof}

\begin{lemma}
\label{formal-defos-lemma-Aut-functor-tangent-space}
设 $\mathcal{F}$ 是 $\mathcal{C}_\Lambda$ 上满足 (RS)
的群胚余纤维范畴。设 $x \in \Ob(\mathcal{F}(A))$，
并设 $x_0$ 是 $x$ 到 $\mathcal{F}(k)$ 的一个推前。
\begin{enumerate}
\item $T_{\text{id}_{x_0}} \mathit{Aut}(x)$ 具有自然的
$k$-向量空间结构，使得加法与
$T_{\text{id}_{x_0}} \mathit{Aut}(x)$ 中的复合一致。
特别地，$T_{\text{id}_{x_0}} \mathit{Aut}(x)$ 中的复合可交换。
\item 有 $k$-向量空间的典范同构
$T_{\text{id}_{x_0}} \mathit{Aut}(x) \to
T_{\text{id}_{x_0}} \mathit{Aut}(x_0)$。
\end{enumerate}
\end{lemma}

\begin{proof}
对函子 $\mathit{Aut}(x) : \mathcal{C}_A \to \textit{Sets}$
及元素 $\text{id}_{x_0} \in \mathit{Aut}(x)(k)$ 应用说明
\ref{formal-defos-remark-localize-cofibered-groupoid}，得到前形变函子
$F = \mathit{Aut}(x)_{\text{id}_{x_0}}$。由引理
\ref{formal-defos-lemma-Aut-functor-RS} 和 \ref{formal-defos-lemma-localize-RS}，
$F$ 是形变函子。根据定义，
$T_{\text{id}_{x_0}} \mathit{Aut}(x) = TF = F(k[\epsilon])$，
它具有引理 \ref{formal-defos-lemma-tangent-space-functor} 所给出的自然
$k$-向量空间结构。

\medskip\noindent
加法定义为复合
$$
F(k[\epsilon]) \times F(k[\epsilon]) \longrightarrow
F(k[\epsilon] \times_k k[\epsilon]) \longrightarrow
F(k[\epsilon]),
$$
其中第一个映射是由 (RS) 保证的双射的逆，而第二个映射由
$k$-代数映射
$k[\epsilon] \times_k k[\epsilon] \to k[\epsilon]$
诱导；该映射把 $(\epsilon, 0)$ 与 $(0, \epsilon)$ 都映到
$\epsilon$。若 $A \to B$ 是 $\mathcal{C}_\Lambda$ 中的环映射，
则 $F(A) \to F(B)$ 是同态，其中
$F(A) = \mathit{Aut}(x)_{\text{id}_{x_0}}(A)$ 与
$F(B) = \mathit{Aut}(x)_{\text{id}_{x_0}}(B)$
都在复合运算下成为群。因此，
$+ : F(k[\epsilon]) \times F(k[\epsilon])\to F(k[\epsilon])$
是同态，其中把 $F(k[\epsilon])$ 看作复合运算下的群。
以 $\text{id} \in F(k[\epsilon])$ 表示单位元。对任意
$v \in F(k[\epsilon])$，有
$+(v, \text{id}) = +(\text{id}, v) = v$，因为
$(\text{id}, v)$ 是 $v$ 沿环映射
$k[\epsilon] \to k[\epsilon] \times_k k[\epsilon]$ 的推前，
该映射满足 $\epsilon \mapsto (\epsilon, 0)$。
一般地，给定乘法为 $\circ$ 的群 $G$，若同态
$+ : G \times G \to G$ 满足
$+(g, 1) = +(1, g) = g$，其中 $1$ 是 $G$ 的单位元，
则 $+ = \circ$。% P10-E0482
这说明 $F(k[\epsilon])$ 的 $k$-向量空间结构中的加法与复合一致。

\medskip\noindent
最后，(2) 只需展开定义。具体地，
$T_{\text{id}_{x_0}} \mathit{Aut}(x)$ 是如下自同构
$\alpha$ 所成的集合：它们是 $x$ 沿
$A \to k \to k[\epsilon]$ 的推前的自同构，并且模
$\epsilon$ 为恒等。另一方面，
$T_{\text{id}_{x_0}} \mathit{Aut}(x_0)$ 是 $x_0$ 沿
$k \to k[\epsilon]$ 的推前的、模 $\epsilon$ 为恒等的自同构集合。
由于 $x_0$ 是 $x$ 沿 $A \to k$ 的推前，结论显然。
\end{proof}

\begin{remark}
\label{formal-defos-remark-infaut-lifting-equalities}
下面指出无穷小自同构群、提升以及自同构函子切空间之间的一些基本关系。
设 $\mathcal{F}$ 是 $\mathcal{C}_\Lambda$ 上的群胚余纤维范畴。
设 $x' \to x$ 是位于环映射 $A' \to A$ 上的态射。
由定义有等式
$$
\text{Inf}(x'/x) = \text{Lift}(\text{id}_x, A'),
$$
其中右端是把 $\text{id}_x$ 视为 $\mathit{Aut}(x')$ 的对象时
所得的提升。若 $x_0 \in \Ob(\mathcal{F}(k))$，且 $x'_0$
是到 $\mathcal{F}(k[\epsilon])$ 的推前，则对
$x'_0 \to x_0$ 应用上述等式，得到
$$
\text{Inf}_{x_0}(\mathcal{F}) =
\text{Lift}(\text{id}_{x_0}, k[\epsilon]) =
T_{\text{id}_{x_0}} \mathit{Aut}(x_0),
$$
最后一个等式直接由定义得出。
\end{remark}

\begin{lemma}
\label{formal-defos-lemma-infaut-vector-space}
设 $\mathcal{F}$ 是 $\mathcal{C}_\Lambda$ 上满足 (RS)
的群胚余纤维范畴，并设 $x_0 \in \Ob(\mathcal{F}(k))$。
则作为集合，$\text{Inf}_{x_0}(\mathcal{F})$ 等于
$T_{\text{id}_{x_0}} \mathit{Aut}(x_0)$，因而具有自然的
$k$-向量空间结构，使得加法与自同构的复合一致。
\end{lemma}

\begin{proof}
集合的等式如说明 \ref{formal-defos-remark-infaut-lifting-equalities} 末尾所述，
而关于向量空间结构的陈述由引理
\ref{formal-defos-lemma-Aut-functor-tangent-space} 得出。
\end{proof}

\begin{lemma}
\label{formal-defos-lemma-k-linear-infaut}
设 $\varphi : \mathcal{F} \to \mathcal{G}$ 是
$\mathcal{C}_\Lambda$ 上满足 (RS) 的群胚余纤维范畴之间的态射。
设 $x_0 \in \Ob(\mathcal{F}(k))$。则 $\varphi$ 诱导
$k$-线性映射
$\text{Inf}_{x_0}(\mathcal{F}) \to
\text{Inf}_{\varphi(x_0)}(\mathcal{G})$。
\end{lemma}

\begin{proof}
显然，$\varphi$ 诱导把恒等映到恒等的态射
$\mathit{Aut}(x_0) \to \mathit{Aut}(\varphi(x_0))$。
因此该结论由切空间的相应结论得出；见引理
\ref{formal-defos-lemma-k-linear-differential}。
\end{proof}

\begin{lemma}
\label{formal-defos-lemma-lifted-automorphisms-torsor}
设 $\mathcal{F}$ 是 $\mathcal{C}_\Lambda$ 上满足 (RS)
的群胚余纤维范畴。设 $x' \to x$ 是位于满环映射
$A' \to A$ 上的态射，该映射的核 $I$ 被
$\mathfrak m_{A'}$ 零化。设 $x_0$ 是 $x$ 到
$\mathcal{F}(k)$ 的一个推前。则
$T_{\text{id}_{x_0}} \mathit{Aut}(x') \otimes_k I
= \text{Inf}_{x_0}(\mathcal{F}) \otimes_k I$
在 $\text{Inf}(x'/x)$ 上有一个自由且传递的作用。
\end{lemma}

\begin{proof}
这就是引理 \ref{formal-defos-lemma-free-transitive-action} 在自同构层情形的类似结论。
确切地，对函子
$\mathit{Aut}(x') : \mathcal{C}_{A'} \to \textit{Sets}$
和元素 $\text{id}_{x_0} \in \mathit{Aut}(x)(k)$ 应用说明
\ref{formal-defos-remark-localize-cofibered-groupoid}，得到前形变函子
$F = \mathit{Aut}(x')_{\text{id}_{x_0}}$。% P10-E0483
由引理 \ref{formal-defos-lemma-Aut-functor-RS} 和 \ref{formal-defos-lemma-localize-RS}，
$F$ 是形变函子。因此，引理
\ref{formal-defos-lemma-free-transitive-action} 给出
$TF \otimes_k I$ 在 $\text{Lift}(\text{id}_x, A')$ 上的
自由且传递作用；由于 $\text{Lift}(\text{id}_x, A')$
是群，它总是非空。注意，由引理
\ref{formal-defos-lemma-Aut-functor-tangent-space}，有向量空间的等式
$$
TF = T_{\text{id}_{x_0}} \mathit{Aut}(x') \otimes_k I =
\text{Inf}_{x_0}(\mathcal{F}) \otimes_k I.
$$ % P10-E0484
最后，说明 \ref{formal-defos-remark-infaut-lifting-equalities} 中的等式
$\text{Inf}(x'/x) = \text{Lift}(\text{id}_x, A')$
完成证明。
\end{proof}

\begin{lemma}
\label{formal-defos-lemma-infaut-trivial}
设 $\mathcal{F}$ 是 $\mathcal{C}_\Lambda$ 上满足 (RS)
的群胚余纤维范畴。设 $x' \to x$ 是 $\mathcal{F}$ 中位于
一个满环映射上的态射。设 $x_0$ 是 $x$ 到
$\mathcal{F}(k)$ 的一个推前。若
$\text{Inf}_{x_0}(\mathcal{F}) = 0$，则
$\text{Inf}(x'/x) = 0$。
\end{lemma}

\begin{proof}
由引理 \ref{formal-defos-lemma-factor-small-extension} 和
\ref{formal-defos-lemma-lifted-automorphisms-torsor} 得出。
\end{proof}

\begin{lemma}
\label{formal-defos-lemma-infdef-trivial}
设 $\mathcal{F}$ 是 $\mathcal{C}_\Lambda$ 上满足 (RS)
的群胚余纤维范畴，并设
$x_0 \in \Ob(\mathcal{F}(k))$。则
$\text{Inf}_{x_0}(\mathcal{F}) = 0$，当且仅当群胚余纤维范畴的自然态射
$\mathcal{F}_{x_0} \to \overline{\mathcal{F}_{x_0}}$ 是等价。
\end{lemma}

\begin{proof}
态射 $\mathcal{F}_{x_0} \to \overline{\mathcal{F}_{x_0}}$ 是等价，
当且仅当 $\mathcal{F}_{x_0}$ 纤维化于集合胚；参见《范畴》第
\ref{categories-section-fibred-in-setoids} 节（根据定义，集合胚是
每个对象的自同构都只有恒等映射的群胚）。我们证明
$\text{Inf}_{x_0}(\mathcal{F}) = 0$ 当且仅当
$\mathcal{F}_{x_0}$ 满足此条件。显然，若
$\mathcal{F}_{x_0}$ 纤维化于集合胚，则
$\text{Inf}_{x_0}(\mathcal{F}) = 0$。反过来，假设
$\text{Inf}_{x_0}(\mathcal{F}) = 0$。设 $A$ 是
$\mathcal{C}_\Lambda$ 的对象。由引理
\ref{formal-defos-lemma-infaut-trivial}，对 $\mathcal{F}_{x_0}(A)$ 中的
任意对象 $x \to x_0$，都有 $\text{Inf}(x/x_0) = 0$。
根据定义，$\text{Inf}(x/x_0)$ 等于
$\mathcal{F}_{x_0}(A)$ 中 $x \to x_0$ 的自同构群，
所以 $\mathcal{F}_{x_0}(A)$ 是集合胚。
\end{proof}

\section{应用}
\label{formal-defos-section-applications}

\noindent
这里汇集一些稍后将用到的形变范畴结果。

\begin{lemma}
\label{formal-defos-lemma-deformation-categories-fiber-product-morphisms}
设 $f : \mathcal{H} \to \mathcal{F}$ 和
$g : \mathcal{G} \to \mathcal{F}$ 是形变范畴的 $1$-态射。则
\begin{enumerate}
\item $\mathcal{W} = \mathcal{H} \times_\mathcal{F} \mathcal{G}$
是形变范畴；并且
\item 有向量空间的 $6$ 项正合序列
$$
0 \to \text{Inf}(\mathcal{W})
\to \text{Inf}(\mathcal{H}) \oplus \text{Inf}(\mathcal{G})
\to \text{Inf}(\mathcal{F}) \to
T\mathcal{W} \to T\mathcal{H} \oplus T\mathcal{G} \to T\mathcal{F}.
$$
\end{enumerate}
\end{lemma}

\begin{proof}
(1) 由引理 \ref{formal-defos-lemma-RS-fiber-product-morphisms} 以及下述事实得出：
$\mathcal{W}(k)$ 是两个具有唯一同构类的集合胚在另一个具有唯一同构类的
集合胚上的纤维积。

\medskip\noindent
证明 (2)。设 $w_0 \in \Ob(\mathcal{W}(k))$，并设
$x_0, y_0, z_0$ 是 $w_0$ 在
$\mathcal{F}, \mathcal{H}, \mathcal{G}$ 中的像。
则 $\text{Inf}(\mathcal{W}) = \text{Inf}_{w_0}(\mathcal{W})$，
对 $\mathcal{H}$、$\mathcal{G}$ 和 $\mathcal{F}$ 也类似；见说明
\ref{formal-defos-remark-trivial-aut-point}。应用引理
\ref{formal-defos-lemma-k-linear-differential} 和
\ref{formal-defos-lemma-k-linear-infaut}，得到除“边界映射”
$\delta : \text{Inf}_{x_0}(\mathcal{F}) \to T\mathcal{W}$
之外的所有线性映射。稍后再插入适当的符号。

\medskip\noindent
构造 $\delta$。选取沿 $k \to k[\epsilon]$
的一个推前 $w_0 \to w'_0$。以 $x'_0,y'_0,z'_0$ 表示
$w'_0$ 在 $\mathcal{F},\mathcal{H},\mathcal{G}$ 中的像。
特别地，得到同构 $b' : f(y'_0) \to x'_0$ 和
$c' : x'_0 \to g(z'_0)$。以 $b : f(y_0) \to x_0$ 和
$c : x_0 \to g(z_0)$ 表示它们沿 $k[\epsilon] \to k$ 的推前。
注意，按照范畴纤维积的显式形式，这意味着
$w'_0 = (k[\epsilon], y'_0, z'_0, c' \circ b')$ 且
$w_0 = (k, y_0, z_0, c \circ b)$；见说明
\ref{formal-defos-remarks-cofibered-groupoids}（\ref{formal-defos-item-fibre-product}）。
给定 $\alpha : x'_0 \to x'_0$，置
$\delta(\alpha) = (k[\epsilon], y'_0, z'_0, c' \circ \alpha \circ b')$。
这确实是 $\mathcal{W}$ 中位于 $k[\epsilon]$ 上的对象，并且带有
位于 $k[\epsilon] \to k$ 上的态射
$(k[\epsilon], y'_0, z'_0, c' \circ \alpha \circ b') \to w_0$，
因为 $\alpha$ 推前到 $k$ 上成为恒等。更一般地，对任意
$k$-向量空间 $V$，可以用完全相同的公式定义映射
$$
\text{Lift}(\text{id}_{x_0}, k[V])
\longrightarrow
\text{Lift}(w_0, k[V]).
$$
此构造关于向量空间 $V$ 具有函子性（略去细节）。因此，
应用引理 \ref{formal-defos-lemma-morphism-linear-functors} 可知
$\delta$ 是 $k$-线性的。

\medskip\noindent
构造这些映射后，容易证明序列正合。第一个映射的单射性来自
$f \times g : \mathcal{W} \to \mathcal{H} \times \mathcal{G}$
是忠实的这一事实。若
$(\beta, \gamma) \in
\text{Inf}_{y_0}(\mathcal{H}) \oplus \text{Inf}_{z_0}(\mathcal{G})$
映到 $\text{Inf}_{x_0}(\mathcal{F})$ 中的同一元素，则
$(\beta, \gamma)$ 定义 $w'_0 =
(k[\epsilon], y'_0, z'_0, c' \circ b')$ 的自同构，
从而得到第二处的正合性。若上述 $\alpha$ 给出 $w_0$ 的平凡形变
$(k[\epsilon], y'_0, z'_0, c' \circ \alpha \circ b')$，
则同构
$w'_0 = (k[\epsilon], y'_0, z'_0, c' \circ b') \to
(k[\epsilon], y'_0, z'_0, c' \circ \alpha \circ b')$
产生一个作为 $\alpha$ 原像的对 $(\beta, \gamma)$。
若 $w = (k[\epsilon], y, z, \phi)$ 是 $w_0$ 的形变，且
$y'_0 \cong y$ 与 $z \cong z'_0$，则映射
$$
f(y'_0) \to f(y) \xrightarrow{\phi} g(z) \to g(z'_0)
$$
是一个在 $\delta$ 下映到 $w$ 的 $\alpha$。最后，若 $y$ 与 $z$
分别是 $y_0$ 与 $z_0$ 的形变，并且存在 $f(y_0)=x_0=g(z_0)$
的形变之间的同构 $\phi : f(y) \to g(z)$，则得到
$T\mathcal{W}$ 中 $(x, y)$ 的原像
$w = (k[\epsilon], y, z, \phi)$。% P10-E0485
证明完毕。
\end{proof}

\begin{lemma}
\label{formal-defos-lemma-map-fibre-products-smooth}
设 $\mathcal{H}_1 \to \mathcal{G}$、
$\mathcal{H}_2 \to \mathcal{G}$ 和 $\mathcal{G} \to \mathcal{F}$
是 $\mathcal{C}_\Lambda$ 上群胚余纤维范畴之间的映射。假设
\begin{enumerate}
\item $\mathcal{F}$ 和 $\mathcal{G}$ 是形变范畴；
\item $T\mathcal{G} \to T\mathcal{F}$ 是单射；并且
\item $\text{Inf}(\mathcal{G}) \to \text{Inf}(\mathcal{F})$ 是满射。
\end{enumerate}
则
$\mathcal{H}_1 \times_\mathcal{G} \mathcal{H}_2 \to
\mathcal{H}_1 \times_\mathcal{F} \mathcal{H}_2$ 是光滑的。
\end{lemma}

\begin{proof}
以 $p_i : \mathcal{H}_i \to \mathcal{G}$ 和
$q : \mathcal{G} \to \mathcal{F}$ 表示给定的映射。% P10-E0486
设 $A' \to A$ 是 $\mathcal{C}_\Lambda$ 中的小扩张。
$\mathcal{H}_1 \times_\mathcal{F} \mathcal{H}_2$ 中位于
$A'$ 上的对象，是三元组 $(x'_1, x'_2, a')$，其中
$x'_i$ 是 $\mathcal{H}_i$ 中位于 $A'$ 上的对象，而
$a' : q(p_1(x'_1)) \to q(p_2(x'_2))$ 是 $\mathcal{F}$
在 $A'$ 上的纤维范畴中的态射。沿 $A' \to A$ 推前，
得到 $(x_1,x_2,a)$。将它提升为
$\mathcal{H}_1 \times_\mathcal{G} \mathcal{H}_2$ 中位于 $A$
上的对象，意指找到位于 $A$ 上且满足 $q(b)=a$ 的态射
$b : p_1(x_1) \to p_2(x_2)$。因此必须证明，可以把 $b$
提升为态射 $b' : p_1(x'_1) \to p_2(x'_2)$，且其在 $q$ 下的像是
$a'$。

\medskip\noindent
注意，可以把
$$
p_1(x'_1) \to p_1(x_1) \xrightarrow{b} p_2(x_2)
\quad\text{和}\quad
p_2(x'_2) \to p_2(x_2)
$$
视为 $\textit{Lift}(p_2(x_2), A' \to A)$ 的两个对象。
函子 $q$ 把它们送到
$$
q(p_1(x'_1)) \to q(p_1(x_1)) \xrightarrow{b} q(p_2(x_2))
\quad\text{和}\quad % P10-E0487
q(p_2(x'_2)) \to q(p_2(x_2))
$$
这两个 $\textit{Lift}(q(p_2(x_2)), A' \to A)$ 的对象；
利用映射 $a' : q(p_1(x'_1)) \to q(p_2(x'_2))$，
它们彼此同构。另一方面，函子
$$
q : \textit{Lift}(p_2(x_2), A' \to A) \to
\textit{Lift}(q(p_2(x_2)), A' \to A)
$$
由引理 \ref{formal-defos-lemma-free-transitive-action} 和关于切空间的假设，
在同构类上诱导单射。因此存在态射
$b' : p_1(x_1') \to p_2(x'_2)$，其到 $A$ 的推前是 $b$。
不过，可能还需要调整 $b'$ 的选择，才能使 $q(b')=a'$。
为此，只需证明
$q : \text{Inf}(p_2(x'_2)/p_2(x_2)) \to
\text{Inf}(q(p_2(x'_2))/q(p_2(x_2)))$ 是满射。
这由关于无穷小自同构的假设和引理
\ref{formal-defos-lemma-lifted-automorphisms-torsor} 得出。
\end{proof}

\begin{lemma}
\label{formal-defos-lemma-easy-check-smooth}
设 $f : \mathcal{F} \to \mathcal{G}$ 是形变范畴之间的映射。
设 $x_0 \in \Ob(\mathcal{F}(k))$，其像为
$y_0 \in \Ob(\mathcal{G}(k))$。若
\begin{enumerate}
\item 映射 $T\mathcal{F} \to T\mathcal{G}$ 是满射；并且
\item 对 $\mathcal{C}_\Lambda$ 中的每个小扩张
$A' \to A$，以及像为 $y \in \mathcal{G}(A)$ 的每个
$x \in \mathcal{F}(A)$，若 $y$ 可以提升到 $A'$，
则 $x$ 也可以提升到 $A'$，
\end{enumerate}
则 $\mathcal{F} \to \mathcal{G}$ 光滑（反之亦然）。
\end{lemma}

\begin{proof}
设 $A' \to A$ 是小扩张，且 $x \in \mathcal{F}(A)$。
设 $y' \to f(x)$ 是 $\mathcal{G}$ 中位于 $A' \to A$ 上的态射。
考虑 $f$ 所诱导的函子
$\text{Lift}(A', x) \to \text{Lift}(A', f(x))$。% P10-E0488
必须证明存在 $\text{Lift}(A', x)$ 的对象 $x' \to x$，
其像为 $y' \to f(x)$；见引理
\ref{formal-defos-lemma-smoothness-small-extensions}。由条件 (2)，
$\text{Lift}(A', x)$ 不是空范畴。由条件 (2) 和引理
\ref{formal-defos-lemma-free-transitive-action}，可知同构类上的映射是满射，
正如所需。% P10-E0489 P10-E0490
\end{proof}

\begin{lemma}
\label{formal-defos-lemma-map-between-smooth}
设 $\mathcal{F} \to \mathcal{G} \to \mathcal{H}$ 是
$\mathcal{C}_\Lambda$ 上群胚余纤维范畴之间的映射。若
\begin{enumerate}
\item $\mathcal{F}$、$\mathcal{G}$ 是形变范畴；
\item 映射 $T\mathcal{F} \to T\mathcal{G}$ 是满射；并且
\item $\mathcal{F} \to \mathcal{H}$ 光滑，
\end{enumerate}
则 $\mathcal{F} \to \mathcal{G}$ 光滑。
\end{lemma}

\begin{proof}
设 $A' \to A$ 是 $\mathcal{C}_\Lambda$ 中的小扩张，
并设 $x \in \mathcal{F}(A)$，其像为
$y \in \mathcal{G}(A)$。假设存在提升
$y' \in \mathcal{G}(A')$。根据引理
\ref{formal-defos-lemma-easy-check-smooth}，只需检验 $x$ 也有提升。
取 $y'$ 的像 $z' \in \mathcal{H}(A')$。由于
$\mathcal{F} \to \mathcal{H}$ 光滑，存在
$x' \in \mathcal{F}(A')$，它同时映到
$x \in \mathcal{F}(A)$ 和 $z' \in \mathcal{H}(A')$；
见定义 \ref{formal-defos-definition-smooth-morphism}。证明完毕。
\end{proof}

\section{任意范畴上的函子中的群胚}
\label{formal-defos-section-groupoids-arbitrary}

\noindent
首先讨论任意范畴上的函子中的群胚的一般性质。下一节再转到范畴
$\mathcal{C}_\Lambda$。为清楚起见，有时把通常的群胚，
即所有态射都是同构的范畴，称为群胚范畴。

\begin{definition}
\label{formal-defos-definition-groupoid-in-functors}
设 $\mathcal{C}$ 是范畴。$\mathcal{C}$ 上的{\it 函子中的群胚范畴}
具有如下对象和态射。
\begin{enumerate}
\item 对象：$\mathcal{C}$ 上的一个{\it 函子中的群胚}是五元组
$(U, R, s, t, c)$，其中 $U, R : \mathcal{C} \to \textit{Sets}$
是函子，而 $s, t : R \to U$ 和
$c : R \times_{s, U, t} R \to R$ 是具有如下性质的态射：
对 $\mathcal{C}$ 的任意对象 $T$，五元组
$$
(U(T), R(T), s, t, c)
$$
是群胚范畴。
\item 态射：$\mathcal{C}$ 上函子中的群胚之间的一个
{\it 态射 $(U, R, s, t, c) \to (U', R', s', t', c')$}
由态射 $U \to U'$ 和 $R \to R'$ 组成，并具有如下性质：
对 $\mathcal{C}$ 的任意对象 $T$，诱导映射
$U(T) \to U'(T)$ 与 $R(T) \to R'(T)$ 定义群胚范畴之间的函子
$$
(U(T), R(T), s, t, c) \to (U'(T), R'(T), s', t', c').
$$
\end{enumerate}
\end{definition}

\begin{remark}
\label{formal-defos-remark-confusion-groupoids-in-functors}
$\mathcal{C}$ 上函子中的群胚等价于函子
$\mathcal{C} \to \textit{Groupoids}$ 的数据，而
$\mathcal{C}$ 上函子中的群胚态射等价于相应函子
$\mathcal{C} \to \textit{Groupoids}$ 之间的态射
（这里把 $\textit{Groupoids}$ 看作 1-范畴）。
不过，为了我们的目的，使用函子中的群胚这一术语更方便。
事实上，把函子中的群胚想成相应的函子
$\mathcal{C} \to \textit{Groupoids}$，或等价地想成与该函子关联的
群胚余纤维范畴，可能导致混淆（说明
\ref{formal-defos-remark-smooth-groupoid-in-functors-warning}）。
\end{remark}

\begin{remark}
\label{formal-defos-remark-identity-inverse}
设 $(U, R, s, t, c)$ 是范畴 $\mathcal{C}$ 上函子中的群胚。
存在唯一的态射 $e : U \to R$ 和 $i : R \to R$，使得对
$\mathcal{C}$ 的每个对象 $T$，$e: U(T) \to R(T)$ 把
$x \in U(T)$ 送到 $x$ 上的恒等态射，而
$i: R(T) \to R(T)$ 把 $a \in U(T)$ 送到 $a$ 在群胚范畴
$(U(T), R(T), s, t, c)$ 中的逆。% P10-E0491
有时分别把 $s$、$t$、$c$、$e$、$i$ 称为“源”、“靶”、
“复合”、“恒等”、“逆”。
\end{remark}

\begin{definition}
\label{formal-defos-definition-representable}
设 $\mathcal{C}$ 是范畴。若 $\mathcal{C}$ 上函子中的群胚
同构于形如 $(\underline{U}, \underline{R}, s, t, c)$ 的群胚，
其中 $U$ 和 $R$ 是 $\mathcal{C}$ 的对象，且推出
$R \amalg_{s, U, t} R$ 存在，则称它{\it 可表}。
\end{definition}

\begin{remark}
\label{formal-defos-remark-reason-existence-coproduct}
因此，$\mathcal{C}$ 上函子中的可表群胚由
$\mathcal{C}$ 的对象 $U$、$R$ 以及态射 $s,t:U\to R$、
$c:R\to R\amalg_{s,U,t}R$ 给出，并且
$(\underline{U}, \underline{R}, s, t, c)$ 满足定义
\ref{formal-defos-definition-groupoid-in-functors} 的条件。要求推出
$R \amalg_{s, U, t} R$ 存在，是为了使复合态射 $c$
能在 $\mathcal{C}$ 的态射层面定义。下文考虑
$\widehat{\mathcal{C}}_\Lambda$ 上函子中的可表群胚时，
此要求总会满足，因为由引理
\ref{formal-defos-lemma-CLambdahat-pushouts}，范畴
$\widehat{\mathcal{C}}_\Lambda$ 具有推出。
\end{remark}

\begin{remark}
\label{formal-defos-remark-simplify-terminology}
我们会用“{\it 设 $(\underline{U}, \underline{R}, s, t, c)$ 是
$\mathcal{C}$ 上函子中的群胚}”来表示所给的是函子中的可表群胚。
因此，这意味着 $U$、$R$ 是 $\mathcal{C}$ 的对象，存在态射
$s,t:U\to R$，推出 $R\amalg_{s,U,t}R$ 存在，
存在态射 $c:R\to R\amalg_{s,U,t}R$，并且
$(\underline{U}, \underline{R}, s, t, c)$ 是
$\mathcal{C}$ 上函子中的群胚。
\end{remark}

\noindent
下面引入函子中群胚的限制记号。稍后把
$\widehat{\mathcal
C}_\Lambda$ 限制到 $\mathcal{C}_\Lambda$ 时将用到它。

\begin{definition}
\label{formal-defos-definition-restricting-groupoids-in-functors}
设 $(U, R, s, t, c)$ 是范畴 $\mathcal{C}$ 上函子中的群胚，
并设 $\mathcal{C}'$ 是 $\mathcal{C}$ 的子范畴。
$(U, R, s, t, c)$ 到 $\mathcal{C}'$ 的{\it 限制
$(U, R, s, t, c)|_{\mathcal{C}'}$}，是 $\mathcal{C}'$
上函子中的群胚
$(U|_{\mathcal{C}'}, R|_{\mathcal
C'}, s|_{\mathcal{C}'}, t|_{\mathcal{C}'}, c|_{\mathcal{C}'})$。
\end{definition}

\begin{remark}
\label{formal-defos-remark-notation-restriction}
在定义 \ref{formal-defos-definition-restricting-groupoids-in-functors} 的情形，
常把 $s|_{\mathcal{C}'}, t|_{\mathcal{C}'}, c|_{\mathcal{C}'}$
简记为 $s,t,c$。
\end{remark}

\begin{definition}
\label{formal-defos-definition-quotient}
设 $(U, R, s, t, c)$ 是范畴 $\mathcal{C}$ 上函子中的群胚。
\begin{enumerate}
\item 对应 $T \mapsto (U(T), R(T), s, t, c)$ 确定函子
$\mathcal{C} \to \textit{Groupoids}$。{\it 商群胚余纤维范畴
$[U/R] \to \mathcal{C}$}，是与该函子关联的
$\mathcal{C}$ 上群胚余纤维范畴（如说明
\ref{formal-defos-remarks-cofibered-groupoids}（\ref{formal-defos-item-construction-associated-cofibered-groupoid}）
所述）。
\item {\it 商态射 $U \to [U/R]$} 是由如下规则定义的
$\mathcal{C}$ 上群胚余纤维范畴之间的态射：
\begin{enumerate}
\item $x \in U(T)$ 映到对象 $(T, x) \in \Ob([U/R](T))$；并且
\item $x \in U(T)$ 和 $f : T \to T'$ 给出位于
$f : T \to T'$ 上的态射
$(f, \text{id}_{U(f)(x)}): (T, x) \to (T, U(f)(x))$。
\end{enumerate}
\end{enumerate}
\end{definition}

\section{基范畴上的函子中的群胚}
\label{formal-defos-section-prorepresentable-groupoids-in-functors}

\noindent
本节讨论 $\mathcal{C}_\Lambda$ 上的函子中的群胚。
最终目标是说明，$\mathcal{C}_\Lambda$ 上函子中的 pro-可表群胚，
对性质良好的形变范畴起“呈现”作用，正如代数空间中的光滑群胚
对代数栈起呈现作用一样；参见《代数栈》第
\ref{algebraic-section-stack-to-presentation} 节。

\begin{definition}
\label{formal-defos-definition-prorepresentable-groupoid-in-functors}
若 $\mathcal{C}_\Lambda$ 上函子中的群胚同构于
$(\underline{R_0}, \underline{R_1}, s, t, c)|_{\mathcal{C}_\Lambda}$，
且 $(\underline{R_0}, \underline{R_1}, s, t, c)$ 是范畴
$\widehat{\mathcal{C}}_\Lambda$ 上函子中的某个可表群胚，
则称原群胚{\it pro-可表}。
\end{definition}

\noindent
设 $(U, R, s, t, c)$ 是 $\mathcal{C}_\Lambda$
上的函子中的群胚。取完备化，得到五元组
$(\widehat{U}, \widehat{R}, \widehat{s}, \widehat{t}, \widehat{c})$。
由说明 \ref{formal-defos-remark-completion-restriction-cofset-adjoint}，
$\text{CofSet}(\mathcal{C}_\Lambda)$ 上的完备化函子是右伴随，
所以与极限可交换。特别地，有典范同构
$$
\widehat{R \times_{s, U, t} R}
\longrightarrow
\widehat{R} \times_{\widehat{s}, \widehat{U}, \widehat{t}} \widehat{R},
$$
故可把 $\widehat{c}$ 看作函子
$\widehat{R} \times_{\widehat{s}, \widehat{U}, \widehat{t}} \widehat{R} \to
\widehat{R}$。于是
$(\widehat{U}, \widehat{R}, \widehat{s}, \widehat{t}, \widehat{c})$
是 $\widehat{\mathcal{C}}_\Lambda$ 上函子中的群胚，
其恒等态射和逆态射分别是 $(U,R,s,t,c)$ 相应态射的完备化。

\begin{definition}
\label{formal-defos-definition-completion-groupoid-in-functors}
设 $(U, R, s, t, c)$ 是 $\mathcal{C}_\Lambda$
上的函子中的群胚。$(U,R,s,t,c)$ 的{\it 完备化
$(U, R, s, t, c)^{\wedge}$}，是上面所述的
$\widehat{\mathcal{C}}_\Lambda$ 上函子中的群胚
$(\widehat{U}, \widehat{R}, \widehat{s}, \widehat{t}, \widehat{c})$。
\end{definition}

\begin{remark}
\label{formal-defos-remark-groupoid-in-functors-complete-restrict}
设 $(U, R, s, t, c)$ 是 $\mathcal{C}_\Lambda$
上的函子中的群胚。则有典范同构
$(U, R, s, t, c)^{\wedge}|_{\mathcal{C}_\Lambda}
\cong (U, R, s, t, c)$；见说明
\ref{formal-defos-remark-restrict-completion}。另一方面，设
$(U, R, s, t, c)$ 是 $\widehat{\mathcal{C}}_\Lambda$
上的函子中的群胚，并且
$U, R : \widehat{\mathcal{C}}_\Lambda \to \textit{Sets}$
都与极限可交换，例如 $U,R$ 可表。则有典范同构
$((U, R, s, t, c)|_{\mathcal{C}_\Lambda})^{\wedge}
\cong (U, R, s, t, c)$。这由说明
\ref{formal-defos-remark-restrict-complete-continuous-functor} 得出。
\end{remark}

\begin{lemma}
\label{formal-defos-lemma-groupoid-in-functors-prorep-equivalences}
设 $(U, R, s, t, c)$ 是 $\mathcal{C}_\Lambda$
上的函子中的群胚。
\begin{enumerate}
\item $(U,R,s,t,c)$ pro-可表，当且仅当其完备化作为
$\widehat{\mathcal{C}}_\Lambda$ 上函子中的群胚可表。
\item $(U,R,s,t,c)$ pro-可表，当且仅当 $U$ 和 $R$ pro-可表。
\end{enumerate}
\end{lemma}

\begin{proof}
(1) 由说明 \ref{formal-defos-remark-groupoid-in-functors-complete-restrict} 得出。
对 (2)，“仅当”方向由函子中 pro-可表群胚的定义立即得出。
反过来，假设 $U$、$R$ pro-可表，即对
$\widehat{\mathcal{C}}_\Lambda$ 的对象 $R_0$、$R_1$，有
$U \cong \underline{R_0}|_{\mathcal{C}_\Lambda}$ 且
$R \cong \underline{R_1}|_{\mathcal{C}_\Lambda}$。由说明
\ref{formal-defos-remark-restrict-complete-continuous-functor}，有
$\underline{R_0} \cong \widehat{\underline{R_0}|_{\mathcal{C}_\Lambda}}$
和
$\underline{R_1} \cong \widehat{\underline{R_1}|_{\mathcal{C}_\Lambda}}$，
所以完备化 $(U,R,s,t,c)^\wedge$ 是形如
$(\underline{R_0}, \underline{R_1}, \widehat{s}, \widehat{t}, \widehat{c})$
的函子中的群胚。由引理 \ref{formal-defos-lemma-CLambdahat-pushouts}，推出
$\underline{R_1} \times_{\widehat{s}, \underline{R_1}, \widehat{t}}
\underline{R_1}$ 存在。% P10-E0492
所以
$(\underline{R_0}, \underline{R_1}, \widehat{s}, \widehat{t}, \widehat{c})$
是 $\widehat{\mathcal{C}}_\Lambda$ 上函子中的可表群胚。
最后，由说明 \ref{formal-defos-remark-groupoid-in-functors-complete-restrict}，
限制
$(\underline{R_0}, \underline{R_1}, s, t, c)|_{\mathcal{C}_\Lambda}$
恢复 $(U,R,s,t,c)$，故根据定义 $(U,R,s,t,c)$ pro-可表。
\end{proof}

\section{基范畴上的函子中的光滑群胚}
\label{formal-defos-section-smooth-minimal-groupoids-in-functors}

\noindent
$\mathcal{C}_\Lambda$ 上函子中群胚的光滑性概念定义如下。

\begin{definition}
\label{formal-defos-definition-smooth-groupoid-in-functors}
设 $(U, R, s, t, c)$ 是 $\mathcal{C}_\Lambda$
上的函子中的群胚。若 $s,t:R\to U$ 光滑，
则称 $(U,R,s,t,c)$ {\it 光滑}。
\end{definition}

\begin{remark}
\label{formal-defos-remark-smooth-groupoid-in-functors-warning}
注意，此术语可能引起混淆：若 $(U,R,s,t,c)$ 是函子中的光滑群胚，
则商 $[U/R]$ 未必是定义
\ref{formal-defos-definition-cofibered-groupoid-projection-smooth} 所指的
光滑群胚余纤维范畴。不过，$(U,R,s,t,c)$ 的光滑性确实蕴含
商态射 $U\to[U/R]$ 光滑（事实上二者等价），正如引理
\ref{formal-defos-lemma-smooth-quotient-morphism} 将要说明的那样。
\end{remark}

\begin{remark}
\label{formal-defos-remark-smooth-power-series-prorepresentable-smooth-groupoid-in-functors}
设
$(\underline{R_0}, \underline{R_1}, s, t, c)|_{\mathcal{C}_\Lambda}$
是 $\mathcal{C}_\Lambda$ 上函子中的 pro-可表群胚。
则 $(\underline{R_0}, \underline{R_1}, s, t, c)|_{\mathcal{C}_\Lambda}$
光滑，当且仅当经 $s$ 和 $t$，$R_1$ 都是
$R_0$ 上的幂级数环。% P10-E0493
这由引理 \ref{formal-defos-lemma-smooth-morphism-power-series} 得出。
\end{remark}

\begin{lemma}
\label{formal-defos-lemma-smooth-quotient-morphism}
设 $(U,R,s,t,c)$ 是 $\mathcal{C}_\Lambda$ 上函子中的群胚。
以下条件等价：
\begin{enumerate}
\item 函子中的群胚 $(U,R,s,t,c)$ 光滑。
\item 态射 $s:R\to U$ 光滑。
\item 态射 $t:R\to U$ 光滑。
\item 商态射 $U\to[U/R]$ 光滑。
\end{enumerate}
\end{lemma}

\begin{proof}
(2) 等价于 (3)，因为 $(U,R,s,t,c)$ 的逆
$i:R\to R$ 是同构，且 $t=s\circ i$。根据定义，(1) 等价于
(2) 与 (3) 同时成立，因而等价于其中任意一个单独成立。

\medskip\noindent
最后证明 (2) 等价于 (4)。展开定义：
\begin{enumerate}
\item[(2)] $s:R\to U$ 光滑等价于下述条件：若
$f:B\to A$ 是 $\mathcal{C}_\Lambda$ 中的满环映射，
$a\in R(A)$、$y\in U(B)$ 且 $s(a)=U(f)(y)$，
则存在 $a'\in R(B)$，使 $R(f)(a')=a$ 且 $s(a')=y$。
\item[(4)] $U\to[U/R]$ 光滑等价于下述条件：若
$f:B\to A$ 是 $\mathcal{C}_\Lambda$ 中的满环映射，且
$(f,a):(B,y)\to(A,x)$ 是 $[U/R]$ 的态射，则存在
$x'\in U(B)$ 和 $b\in R(B)$，使 $s(b)=x'$、$t(b)=y$，
且 $c(a,R(f)(b))=e(x)$。这里 $e:U\to R$ 表示恒等，
而记号 $(f,a)$ 如说明
\ref{formal-defos-remarks-cofibered-groupoids}（\ref{formal-defos-item-construction-associated-cofibered-groupoid}）
所述；特别地，$a\in R(A)$，且
$s(a)=U(f)(y)$、$t(a)=x$。
\end{enumerate}
若 (4) 成立，并给定 (2) 中那样的 $f,a,y$，令 $x=t(a)$，
于是有态射 $(f,a):(B,y)\to(A,x)$。(4) 给出 $x'$ 和 $b$，
而 $a'=i(b)$ 满足 (2) 的要求。反过来，假设 (2) 成立，
并给定 (4) 中那样的 $(f,a):(B,y)\to(A,x)$。
则 (2) 给出 $a'\in R(B)$，而 $x'=t(a')$ 和 $b=i(a')$
满足 (4) 的要求。
\end{proof}

\section{作为函子中群胚之商的形变范畴}
\label{formal-defos-section-deformation-categories-as-quotients}

\noindent
我们讨论 $\mathcal{C}_\Lambda$ 上函子中的群胚应满足哪些条件，
才能保证其商是形变范畴，并计算这种商的切空间和无穷小自同构空间。

\begin{lemma}
\label{formal-defos-lemma-smooth-RS-groupoid-in-functors-quotient}
设 $(U,R,s,t,c)$ 是 $\mathcal{C}_\Lambda$ 上函子中的光滑群胚。
假设 $U$ 和 $R$ 满足 (RS)。则 $[U/R]$ 满足 (RS)。
\end{lemma}

\begin{proof}
设
$$
\xymatrix{
                           &     (A_2, x_2) \ar[d]^{(f_2, a_2)} \\
(A_1, x_1) \ar[r]^{(f_1, a_1)} &     (A, x)
}
$$
是 $[U/R]$ 中的图表，其中 $f_2:A_2\to A$ 为满射。
记号如说明 \ref{formal-defos-remarks-cofibered-groupoids}
（\ref{formal-defos-item-construction-associated-cofibered-groupoid}）所述。
因此 $f_1:A_1\to A, f_2:A_2\to A$ 是
$\mathcal{C}_\Lambda$ 中的映射，
$x\in U(A)$、$x_1\in U(A_1)$、$x_2\in U(A_2)$，
并且 $a_1,a_2\in R(A)$，其中
$s(a_1)=U(f_1)(x_1)$、$t(a_1)=x$，且
$s(a_2)=U(f_2)(x_2)$、$t(a_2)=x$。
我们如下构造该 $[U/R]$ 中图表的一个纤维积，
它位于 $A_1\times_A A_2$ 上方。

\medskip\noindent
令 $a=c(i(a_1),a_2)$，其中 $i:R\to R$ 是逆态射。
则 $a\in R(A)$、$x_2\in U(A_2)$，且
$s(a)=U(f_2)(x_2)$。因而有元素
$(a,x_2)\in R(A)\times_{s,U(A),U(f_2)}U(A_2)$。
由于 $s:R\to U$ 光滑，存在元素
$\widetilde{a}\in R(A_2)$，使
$R(f_2)(\widetilde{a})=a$ 且 $s(\widetilde{a})=x_2$。
特别地，
$U(f_2)(t(\widetilde{a}))=t(a)=U(f_1)(x_1)$。
所以 $x_1$ 和 $t(\widetilde{a})$ 确定元素
$$
(x_1, t(\widetilde{a})) \in U(A_1) \times_{U(A)} U(A_2).
$$
由 $U$ 满足 (RS) 的假设，有等同
$U(A_1)\times_{U(A)}U(A_2)=U(A_1\times_A A_2)$。
我们用
$x_1\times t(\widetilde{a})\in U(A_1\times_A A_2)$
表示与
$(x_1,t(\widetilde{a}))\in U(A_1)\times_{U(A)}U(A_2)$
对应的元素。% P10-E0494
令 $p_1,p_2$ 为 $A_1\times_A A_2$ 的投影。我们断言
$$
\xymatrix{
(A_1 \times_A A_2, x_1 \times t(\widetilde{a})) \ar[d]_{(p_1, e(x_1))}
\ar[rr]_-{(p_2, i(\widetilde{a}))} & & (A_2, x_2) \ar[d]^{(f_2, a_2)} \\
(A_1, x_1) \ar[rr]^{(f_1, a_1)} & & (A, x)
}
$$
是 $[U/R]$ 中的纤维方块。（注意，$e:U\to R$ 表示恒等。）

\medskip\noindent
该图表可交换，因为
$c(a_2,R(f_2)(i(\widetilde{a})))=c(a_2,i(a))=a_1$。
为了验证它是纤维方块，设
$$
\xymatrix{
(B, z) \ar[d]_{(g_1, b_1)} \ar[rr]_{(g_2, b_2)} & & (A_2, x_2)
\ar[d]^{(f_2, a_2)} \\
(A_1, x_1) \ar[rr]^{(f_1, a_1)} & & (A, x)
}
$$
是 $[U/R]$ 中的可交换图表。我们将证明存在唯一态射
$(g,b):(B,z)\to(A_1\times_A A_2,x_1\times t(\widetilde{a}))$，
与到 $(A_1,x_1)$ 和 $(A_2,x_2)$ 的态射相容。
必有 $g=(g_1,g_2):B\to A_1\times_A A_2$。
由于根据假设 $R$ 满足 (RS)，有等同
$R(A_1\times_A A_2)=R(A_1)\times_{R(A)}R(A_2)$。
因此可写成 $b=(b'_1,b'_2)$，其中
$b'_1\in R(A_1)$、$b'_2\in R(A_2)$，且二者在 $R(A)$ 中相等。
于是
$((g_1,g_2),(b'_1,b'_2)):(B,z)\to
(A_1\times_A A_2,x_1\times t(\widetilde{a}))$
与投影可交换，当且仅当
$b'_1=b_1$ 且 $b'_2=c(\widetilde{a},b_2)$；
这就证明了唯一性和存在性。
\end{proof}

\begin{lemma}
\label{formal-defos-lemma-deformation-groupoid-quotient}
设 $(U,R,s,t,c)$ 是 $\mathcal{C}_\Lambda$ 上函子中的光滑群胚。
假设 $U$ 和 $R$ 是形变函子。则：
\begin{enumerate}
\item 商 $[U/R]$ 是形变范畴。
\item $[U/R]$ 的切空间为
$$
T[U/R] = \Coker(ds-dt: TR \to TU).
$$
\item $[U/R]$ 的无穷小自同构空间为
$$
\text{Inf}([U/R]) =
\Ker(ds \oplus dt : TR \to TU \oplus TU).
$$
\end{enumerate}
\end{lemma}

\begin{proof}
由于 $U$ 和 $R$ 是形变函子，$[U/R]$ 是预形变范畴。
根据定义，形变函子满足 (RS)，而由引理
\ref{formal-defos-lemma-smooth-RS-groupoid-in-functors-quotient} 可见
[U/R] 满足 (RS)。% P10-E0495
所以 $[U/R]$ 是形变范畴。
(2) 和 (3) 直接由定义得出。
\end{proof}

\section{群胚余纤维范畴的呈现}
\label{formal-defos-section-presentation-categories-cofibred-in-groupoids}

\noindent
呈现定义如下。

\begin{definition}
\label{formal-defos-definition-presentation}
设 $\mathcal{F}$ 是范畴 $\mathcal{C}$ 上的群胚余纤维范畴。
设 $(U,R,s,t,c)$ 是 $\mathcal{C}$ 上函子中的群胚。
所谓{\it 由 $(U,R,s,t,c)$ 给出的 $\mathcal{F}$ 的呈现}，
是指 $\mathcal{C}$ 上群胚余纤维范畴的一个等价
$\varphi:[U/R]\to\mathcal{F}$。
\end{definition}

\noindent
下面两个一般性引理将用于得到呈现。

\begin{lemma}
\label{formal-defos-lemma-presentation-construction}
设 $\mathcal{F}$ 是范畴 $\mathcal{C}$ 上的群胚余纤维范畴。% P10-E0496
设 $U:\mathcal{C}\to\textit{Sets}$ 是函子。
设 $f:U\to\mathcal{F}$ 是 $\mathcal{C}$
上群胚余纤维范畴的态射。如下定义 $R,s,t,c$：
\begin{enumerate}
\item $R:\mathcal{C}\to\textit{Sets}$ 是函子
$U\times_{f,\mathcal{F},f}U$。
\item $t,s:R\to U$ 分别为第一投影和第二投影。
\item $c:R\times_{s,U,t}R\to R$ 是如下态射：
在典范同构
$R\times_{s,U,t}R\to
U\times_{f,\mathcal{F},f}U\times_{f,\mathcal{F},f}U$
下，它由
$U\times_{f,\mathcal{F},f}U\times_{f,\mathcal{F},f}U$
到第一因子和最后一个因子的投影给出。
\end{enumerate}
则 $(U,R,s,t,c)$ 是 $\mathcal{C}$ 上函子中的群胚。
\end{lemma}

\begin{proof}
略。
\end{proof}

\begin{lemma}
\label{formal-defos-lemma-presentation-morphism}
设 $\mathcal{F}$ 是范畴 $\mathcal{C}$ 上的群胚余纤维范畴。% P10-E0496
设 $U:\mathcal{C}\to\textit{Sets}$ 是函子。
设 $f:U\to\mathcal{F}$ 是 $\mathcal{C}$
上群胚余纤维范畴的态射。设 $(U,R,s,t,c)$ 是引理
\ref{formal-defos-lemma-presentation-construction} 中由
$f:U\to\mathcal{F}$ 构造的 $\mathcal{C}$ 上函子中的群胚。
则存在自然态射 $[f]:[U/R]\to\mathcal{F}$，满足：
\begin{enumerate}
\item $[f]:[U/R]\to\mathcal{F}$ 全忠实。
\item $[f]:[U/R]\to\mathcal{F}$ 是等价，当且仅当
$f:U\to\mathcal{F}$ 本质满。
\end{enumerate}
\end{lemma}

\begin{proof}
略。
\end{proof}

\section{形变范畴的呈现}
\label{formal-defos-section-presentation-deformation-categories}

\noindent
根据下一个引理，从预形变函子到预形变范畴 $\mathcal{F}$
的光滑态射，会给出一个由函子中光滑群胚构成的
$\mathcal{F}$ 的呈现。

\begin{lemma}
\label{formal-defos-lemma-smooth-groupoid-in-functors-construction}
设 $\mathcal{F}$ 是 $\mathcal{C}_\Lambda$ 上的群胚余纤维范畴。
设 $U:\mathcal{C}_\Lambda\to\textit{Sets}$ 是函子。
设 $f:U\to\mathcal{F}$ 是群胚余纤维范畴的光滑态射。则：
\begin{enumerate}
\item 若 $(U,R,s,t,c)$ 是引理
\ref{formal-defos-lemma-presentation-construction} 中由
$f:U\to\mathcal{F}$ 构造的 $\mathcal{C}_\Lambda$
上函子中的群胚，则 $(U,R,s,t,c)$ 光滑。
\item 若 $f:U(k)\to\mathcal{F}(k)$ 本质满，则引理
\ref{formal-defos-lemma-presentation-morphism} 中的态射
$[f]:[U/R]\to\mathcal{F}$ 是等价。
\end{enumerate}
\end{lemma}

\begin{proof}
根据引理 \ref{formal-defos-lemma-presentation-construction} 的构造，
有可交换图表
$$
\xymatrix{
R = U \times_{f, \mathcal{F}, f} U \ar[r]_-s \ar[d]_t & U
\ar[d]^f \\
U \ar[r]^f & \mathcal{F}
}
$$
其中 $t,s$ 分别为第一投影和第二投影。于是由引理
\ref{formal-defos-lemma-smooth-properties}，$t,s$ 光滑。因此 (1) 成立。

\medskip\noindent
若 (2) 的假设成立，则由引理
\ref{formal-defos-lemma-smooth-morphism-essentially-surjective}，
态射 $f:U\to\mathcal{F}$ 本质满。因此由引理
\ref{formal-defos-lemma-presentation-morphism}，态射
$[f]:[U/R]\to\mathcal{F}$ 是等价。
\end{proof}

\begin{lemma}
\label{formal-defos-lemma-deformation-functor-diagonal}
设 $\mathcal{F}$ 是形变范畴。
设 $U:\mathcal{C}_\Lambda\to\textit{Sets}$ 是形变函子。
设 $f:U\to\mathcal{F}$ 是群胚余纤维范畴的态射。
则 $U\times_{f,\mathcal{F},f}U$ 是形变函子，
其切空间嵌入如下 $k$-向量空间正合序列：
$$
0 \to \text{Inf}(\mathcal{F}) \to
T(U \times_{f, \mathcal{F}, f} U) \to TU \oplus TU
$$
\end{lemma}

\begin{proof}
这由引理 \ref{formal-defos-lemma-deformation-categories-fiber-product-morphisms}
以及 $\text{Inf}(U)=(0)$ 得出。
\end{proof}

\begin{lemma}
\label{formal-defos-lemma-prorepresentable-groupoid-in-functors-construction}
设 $\mathcal{F}$ 是形变范畴。
设 $U:\mathcal{C}_\Lambda\to\textit{Sets}$ 是 pro-可表函子。
设 $f:U\to\mathcal{F}$ 是群胚余纤维范畴的态射。
设 $(U,R,s,t,c)$ 是引理 \ref{formal-defos-lemma-presentation-construction}
中由 $f:U\to\mathcal{F}$ 构造的
$\mathcal{C}_\Lambda$ 上函子中的群胚。若
$\dim_k\text{Inf}(\mathcal{F})<\infty$，则
$(U,R,s,t,c)$ pro-可表。
\end{lemma}

\begin{proof}
注意，由例 \ref{formal-defos-example-prorepresentable-deformation-functor}，
$U$ 是形变函子。由引理
\ref{formal-defos-lemma-deformation-functor-diagonal} 可见，
$R=U\times_{f,\mathcal{F},f}U$ 是形变函子，其切空间
$TR=T(U\times_{f,\mathcal{F},f}U)$ 位于正合序列
$0\to\text{Inf}(\mathcal{F})\to TR\to TU\oplus TU$ 中。
由于我们已假设第一个空间维数有限，并且由例
\ref{formal-defos-example-tangent-space-prorepresentable-functor}，
$TU$ 的维数有限，所以 $\dim TR<\infty$。
映射 $\gamma:\text{Der}_\Lambda(k,k)\to TR$，见
（\ref{formal-defos-equation-map}），是单射；% P10-E0497
这是因为它与 $TR\to TU$ 的复合，由关于 pro-可表函子 $U$
的定理 \ref{formal-defos-theorem-Schlessinger-prorepresentability} 可知是单射。
于是由定理 \ref{formal-defos-theorem-Schlessinger-prorepresentability}，
$R$ pro-可表。再由引理
\ref{formal-defos-lemma-groupoid-in-functors-prorep-equivalences}，
$(U,R,s,t,c)$ pro-可表。
\end{proof}

\begin{theorem}
\label{formal-defos-theorem-presentation-deformation-groupoid}
设 $\mathcal{F}$ 是 $\mathcal{C}_\Lambda$ 上的群胚余纤维范畴。
则 $\mathcal{F}$ 存在一个由 $\mathcal{C}_\Lambda$
上函子中的光滑 pro-可表群胚给出的呈现，当且仅当以下条件成立：
\begin{enumerate}
\item $\mathcal{F}$ 是形变范畴。
\item $\dim_kT\mathcal{F}$ 有限。
\item $\dim_k\text{Inf}(\mathcal{F})$ 有限。
\end{enumerate}
\end{theorem}

\begin{proof}
回顾例 \ref{formal-defos-example-prorepresentable-deformation-functor}：
pro-可表函子是形变函子。因此，若 $\mathcal{F}$
等价于函子中的一个光滑 pro-可表群胚，% P10-E0498
则条件 (1)、(2)、(3) 分别由引理
\ref{formal-defos-lemma-deformation-groupoid-quotient} 的 (1)、(2)、(3) 得出。

\medskip\noindent
反过来，假设条件 (1)、(2)、(3) 成立。由条件 (1) 可知
(S1) 和 (S2) 成立；见引理 \ref{formal-defos-lemma-RS-implies-S1-S2}。
由引理 \ref{formal-defos-lemma-versal-object-existence}，存在万有形式对象
$\xi$。令 $U=\underline{R}|_{\mathcal{C}_\Lambda}$，% P10-E0499
则相应映射 $\underline{\xi}:U\to\mathcal{F}$ 光滑
（这正是万有形式对象的定义）。
设 $(U,R,s,t,c)$ 是引理
\ref{formal-defos-lemma-presentation-construction} 中由映射
$\underline{\xi}$ 构造的函子中的群胚。
由引理 \ref{formal-defos-lemma-smooth-groupoid-in-functors-construction} 可见，
$(U,R,s,t,c)$ 是函子中的光滑群胚，且
$[U/R]\to\mathcal{F}$ 是等价。由引理
\ref{formal-defos-lemma-prorepresentable-groupoid-in-functors-construction} 可见，
$(U,R,s,t,c)$ pro-可表。因此
$[U/R]\to\mathcal{F}$ 是所求的 $\mathcal{F}$ 的呈现。
\end{proof}

\section{关于极小性的说明}
\label{formal-defos-section-minimality}

\noindent
本章的主要定理是上面的定理
\ref{formal-defos-theorem-presentation-deformation-groupoid}。
它完全刻画了 $\mathcal{C}_\Lambda$ 上那些存在函子中
光滑 pro-可表群胚呈现的群胚余纤维范畴。
本节简要讨论如何利用第 \ref{formal-defos-section-minimal-versal} 节和
第 \ref{formal-defos-section-miniversal-objects-existence} 节所讨论的极小性，
得到一个“极小”光滑 pro-可表呈现。

\begin{definition}
\label{formal-defos-definition-minimal-groupoid-in-functors}
设 $(U,R,s,t,c)$ 是 $\mathcal{C}_\Lambda$
上函子中的光滑 pro-可表群胚。
\begin{enumerate}
\item 若群胚 $(U(k[\epsilon]),R(k[\epsilon]),s,t,c)$
全不连通，即任意两个不同对象之间都没有态射，
则称 $(U,R,s,t,c)$ {\it 规范化}。
\item 若 $U\to[U/R]$ 由 $[U/R]$
的一个极小万有形式对象给出，% P10-E0500
则称 $(U,R,s,t,c)$ {\it 极小}。
\end{enumerate}
\end{definition}

\noindent
这两个概念的差别与条件（\ref{formal-defos-equation-bijective}）和
（\ref{formal-defos-equation-bijective-orbits}）之间的差别有关，
并在 $k'\subset k$ 可分时消失。此外，函子中的规范化光滑
pro-可表群胚是极小的，正如下一个引理所示。准确陈述如下。

\begin{lemma}
\label{formal-defos-lemma-characterize-minimal-groupoid-in-functors}
设 $(U,R,s,t,c)$ 是 $\mathcal{C}_\Lambda$
上函子中的光滑 pro-可表群胚。
\begin{enumerate}
\item $(U,R,s,t,c)$ 规范化，当且仅当态射
$U\to[U/R]$ 在切空间上诱导同构；
\item $(U,R,s,t,c)$ 极小，当且仅当
$TU\to T[U/R]$ 的核包含在
$\text{Der}_\Lambda(k,k)\to TU$ 的像中。
\end{enumerate}
\end{lemma}

\begin{proof}
(1) 直接由定义得出。为证明 (2)，令
$\mathcal{F}=[U/R]$。由于 $\mathcal{F}$ 有一个呈现，
它是形变范畴；见定理
\ref{formal-defos-theorem-presentation-deformation-groupoid}。
特别地，它满足 (RS)、(S1)、(S2)；见引理
\ref{formal-defos-lemma-RS-implies-S1-S2}。
回顾引理 \ref{formal-defos-lemma-minimal-versal}：
极小万有形式对象在同构意义下唯一。
由定理 \ref{formal-defos-theorem-miniversal-object-existence}，
一个极小万有对象诱导满足
（\ref{formal-defos-equation-bijective-orbits}）的映射
$\underline{\xi}:\underline{R}|_{\mathcal{C}_\Lambda}
\to\mathcal{F}$。由于在 $\mathcal{F}$ 上有
$U\cong\underline{R}|_{\mathcal{C}_\Lambda}$，% P10-E0501
可见 $TU\to T\mathcal{F}=T[U/R]$
满足引理所述的性质。
\end{proof}

\noindent
$\mathcal C_\Lambda$ 上函子中的极小 pro-可表群胚之商，
不存在不是自同构的自等价。为证明这一点，先注意下列引理。

\begin{lemma}
\label{formal-defos-lemma-surjective-morphism-prorepresentable-functor}
设 $U:\mathcal{C}_\Lambda\to\textit{Sets}$ 是 pro-可表函子。
设 $\varphi:U\to U$ 是态射，并且
$d\varphi:TU\to TU$ 是同构。则 $\varphi$ 是同构。
\end{lemma}

\begin{proof}
若对某个 $R\in\Ob(\widehat{\mathcal{C}}_\Lambda)$，
$U\cong\underline{R}|_{\mathcal{C}_\Lambda}$，
则完备化 $\varphi$ 得到态射
$\underline{R}\to\underline{R}$。
若 $f:R\to R$ 是 $\widehat{\mathcal{C}}_\Lambda$
中相应的态射，则 $f$ 诱导同构
$\text{Der}_\Lambda(R,k)\to\text{Der}_\Lambda(R,k)$；
见例 \ref{formal-defos-example-tangent-space-map-prorepresentable-functor}。
特别地，由引理 \ref{formal-defos-lemma-derivations-surjective}，$f$ 是满射。
诺特环的满自同态是同构（见《代数》引理
\ref{algebra-lemma-surjective-endo-noetherian-ring-is-iso}），
所以 $f$ 是同构，从而 $\underline{R}\to\underline{R}$
是同构，进而 $\varphi:U\to U$ 是同构。
\end{proof}

\begin{lemma}
\label{formal-defos-lemma-minimal-prorepresentable-groupoid-autoequivalence}
设 $(U,R,s,t,c)$ 是 $\mathcal{C}_\Lambda$
上函子中的极小光滑 pro-可表群胚。
若 $\varphi:[U/R]\to[U/R]$ 是群胚余纤维范畴的等价，
则 $\varphi$ 是同构。
\end{lemma}

\begin{proof}
态射 $\varphi:[U/R]\to[U/R]$ 与 $\mathcal{C}_\Lambda$
上函子中群胚的态射
$\varphi:(U,R,s,t,c)\to(U,R,s,t,c)$ 是同一回事，% P10-E0502
这里后者如定义 \ref{formal-defos-definition-groupoid-in-functors} 所定义。
把相应态射记为 $\phi:U\to U$ 和 $\psi:R\to R$。
由于图表
$$
\xymatrix{
& \text{Der}_\Lambda(k, k) \ar[dr]_\gamma \ar[dl]^\gamma \\
TU \ar[rr]_{d\phi} \ar[d] & & TU \ar[d]  \\
T[U/R] \ar[rr]^{d\varphi} & & T[U/R]
}
$$
可交换，$d\varphi$ 双射，并且有引理
\ref{formal-defos-lemma-characterize-minimal-groupoid-in-functors}
中的极小性刻画，所以 $d\phi$ 是单射
（因维数理由，它也是双射）。因此由引理
\ref{formal-defos-lemma-surjective-morphism-prorepresentable-functor}，
$\phi:U\to U$ 是同构。
利用引理 \ref{formal-defos-lemma-deformation-functor-diagonal} 的正合序列
$$
0 \to \text{Inf}([U/R]) \to TR \to TU \oplus TU
$$
作类似论证，可证明 $\psi:R\to R$ 是同构。
但这也可由 $R=U\times_{[U/R]}U$ 以及
$\varphi$ 和 $\phi$ 都是同构直接得出。% P10-E0503
\end{proof}

\begin{lemma}
\label{formal-defos-lemma-minimal-prorepresentable-groupoid-equivalence}
设 $(U,R,s,t,c)$ 和 $(U',R',s',t',c')$ 是
$\mathcal{C}_\Lambda$ 上函子中的极小光滑 pro-可表群胚。
若 $\varphi:[U/R]\to[U'/R']$ 是群胚余纤维范畴的等价，
则 $\varphi$ 是同构。
\end{lemma}

\begin{proof}
设 $\psi:[U'/R']\to[U/R]$ 是 $\varphi$ 的拟逆。
由引理 \ref{formal-defos-lemma-minimal-prorepresentable-groupoid-autoequivalence}，
$\psi\circ\varphi$ 和 $\varphi\circ\psi$ 都是同构，
因此 $\varphi$ 和 $\psi$ 都是同构。
\end{proof}

\noindent
下一个引理概括了本章前面得到的一些结果。

\begin{lemma}
\label{formal-defos-lemma-minimal-groupoid-in-functors-construction}
设 $\mathcal{F}$ 是形变范畴，并且
$\dim_kT\mathcal{F}<\infty$ 且
$\dim_k\text{Inf}(\mathcal{F})<\infty$。
则 $\mathcal{F}$ 存在一个极小万有形式对象 $\xi$。
设 $\xi$ 位于 $R\in\Ob(\widehat{\mathcal{C}}_\Lambda)$ 上方。
令 $U=\underline{R}|_{\mathcal{C}_\Lambda}$。
令 $f=\underline{\xi}:U\to\mathcal{F}$ 为相应态射。
设 $(U,R,s,t,c)$ 是引理 \ref{formal-defos-lemma-presentation-construction}
中由 $f:U\to\mathcal{F}$ 构造的
$\mathcal{C}_\Lambda$ 上函子中的群胚。% P10-E0504
则 $(U,R,s,t,c)$ 是 $\mathcal{C}_\Lambda$
上函子中的极小光滑 pro-可表群胚，并且存在等价
$[U/R]\to\mathcal{F}$。
\end{lemma}

\begin{proof}
由于 $\mathcal{F}$ 是形变范畴，它满足 (S1) 和 (S2)；
见引理 \ref{formal-defos-lemma-RS-implies-S1-S2}。
由引理 \ref{formal-defos-lemma-versal-object-existence}，存在万有形式对象。
由引理 \ref{formal-defos-lemma-minimal-versal}，存在引理陈述中那样的
极小万有形式对象 $\xi/R$。令
$U=\underline{R}|_{\mathcal{C}_\Lambda}$，% P10-E0505
则相应映射 $\underline{\xi}:U\to\mathcal{F}$ 光滑
（这正是万有形式对象的定义）。
设 $(U,R,s,t,c)$ 是引理
\ref{formal-defos-lemma-presentation-construction} 中由映射
$\underline{\xi}$ 构造的函子中的群胚。由引理
\ref{formal-defos-lemma-smooth-groupoid-in-functors-construction} 可见，
$(U,R,s,t,c)$ 是函子中的光滑群胚，并且
$[U/R]\to\mathcal{F}$ 是等价。由引理
\ref{formal-defos-lemma-prorepresentable-groupoid-in-functors-construction} 可见，
$(U,R,s,t,c)$ pro-可表。
最后，$(U,R,s,t,c)$ 极小，因为
$U\to[U/R]=\mathcal{F}$ 对应于极小万有形式对象 $\xi$。% P10-E0506
\end{proof}

\noindent
由函子中的极小 pro-可表群胚给出的呈现满足下列唯一性。

\begin{lemma}
\label{formal-defos-lemma-minimal-presentations-equivalent}
设 $\mathcal{F}$ 是 $\mathcal{C}_\Lambda$
上的群胚余纤维范畴。% P10-E0507
假设 $\mathcal{F}$ 存在分别由函子中的极小光滑 pro-可表群胚
$(U,R,s,t,c)$ 和 $(U',R',s',t',c')$ 给出的呈现。
则 $(U,R,s,t,c)$ 与 $(U',R',s',t',c')$ 同构。
\end{lemma}

\begin{proof}
这由引理 \ref{formal-defos-lemma-minimal-prorepresentable-groupoid-equivalence}
以及下述观察得出：态射 $[U/R]\to[U'/R']$
与函子中群胚的态射是同一回事；% P10-E0502
这是因为定义 \ref{formal-defos-definition-quotient}
明确构造了 $[U/R]$。
\end{proof}

\noindent
综上，我们证明了下列定理。

\begin{theorem}
\label{formal-defos-theorem-minimal-smooth-prorepresentable-presentations}
设 $\mathcal{F}$ 是 $\mathcal{C}_\Lambda$ 上的群胚余纤维范畴。
考虑下列条件：
\begin{enumerate}
\item $\mathcal{F}$ 存在一个由 $\mathcal{C}_\Lambda$
上函子中的规范化光滑 pro-可表群胚给出的呈现；
\item $\mathcal{F}$ 存在一个由 $\mathcal{C}_\Lambda$
上函子中的光滑 pro-可表群胚给出的呈现；
\item $\mathcal{F}$ 存在一个由 $\mathcal{C}_\Lambda$
上函子中的极小光滑 pro-可表群胚给出的呈现；并且
\item $\mathcal{F}$ 满足以下条件：
\begin{enumerate}
\item $\mathcal{F}$ 是形变范畴。
\item $\dim_kT\mathcal{F}$ 有限。
\item $\dim_k\text{Inf}(\mathcal{F})$ 有限。
\end{enumerate}
\end{enumerate}
则 (2)、(3)、(4) 等价，并且都由 (1) 蕴含。
若 $k'\subset k$ 可分，则 (1)、(2)、(3)、(4) 全部等价。
此外，为 $\mathcal{F}$ 提供呈现的函子中极小光滑
pro-可表群胚在同构意义下唯一。
\end{theorem}

\begin{proof}
由引理 \ref{formal-defos-lemma-characterize-minimal-groupoid-in-functors}，
(1) 蕴含 (3)，并且在 $k'\subset k$ 可分时 (1) 等价于 (3)。
显然 (3) 蕴含 (2)。由定理
\ref{formal-defos-theorem-presentation-deformation-groupoid}，
(2) 蕴含 (4)。由引理
\ref{formal-defos-lemma-minimal-groupoid-in-functors-construction}，
(4) 蕴含 (3)。这就证明了所有蕴含关系。
最后的唯一性陈述由引理
\ref{formal-defos-lemma-minimal-presentations-equivalent} 得出。
\end{proof}

\section{万有环的唯一性}
\label{formal-defos-section-uniqueness}

\noindent
给定 $\widehat{\mathcal{C}}_\Lambda$ 中的 $R,S$，
若存在 $r\geq0$ 以及 $\widehat{\mathcal{C}}_\Lambda$ 中的映射
$h:R\to R[[t_1,\ldots,t_r]]$ 和
$k:R[[t_1,\ldots,t_r]]\to S$，使得对所有 $a\in R$，
\begin{enumerate}
\item $h(a)\bmod(t_1,\ldots,t_r)=a$；
\item $f(a)=k(a)$；
\item $g(a)=k(h(a))$，
\end{enumerate}
则称映射 $f,g:R\to S$ {\it 形式同伦}。
我们称 $(r,h,k)$ 是 $f$ 与 $g$ 之间的一个{\it 形式同伦}。

\begin{lemma}
\label{formal-defos-lemma-homotopy}
形式同伦在 $\widehat{\mathcal{C}}_\Lambda$
中的态射集合上构成等价关系。
\end{lemma}

\begin{proof}
假设任取 $r\geq1$，并有两个映射
$h_1,h_2:R\to R[[t_1,\ldots,t_r]]$，使得对所有 $a\in R$，
$h_1(a)\bmod(t_1,\ldots,t_r)
=h_2(a)\bmod(t_1,\ldots,t_r)=a$，
以及映射 $k:R[[t_1,\ldots,t_r]]\to S$。
我们断言 $k\circ h_1$ 与 $k\circ h_2$ 形式同伦。
该断言内在的对称性将说明形式同伦概念是对称的。% P10-E0508
事实上，映射
$$
\Psi :
R[[t_1, \ldots, t_r]]
\longrightarrow
R[[t_1, \ldots, t_r]],\quad
\sum a_I t^I \longmapsto \sum h_1(a_I)t^I
$$
是同构。对 $a\in R$，令
$h(a)=\Psi^{-1}(h_2(a))$，并令 $k'=k\circ\Psi$，
则可见 $(r,h,k')$ 是 $k\circ h_1$ 与
$k\circ h_2$ 之间的形式同伦，从而证明了该断言。% P10-E0509

\medskip\noindent
设有如上三个映射 $f_1,f_2,f_3:R\to S$，
以及 $f_1$ 与 $f_2$ 之间的形式同伦 $(r_1,h_1,k_1)$，
和 $f_3$ 与 $f_2$ 之间的形式同伦 $(r_2,h_2,k_2)$（！）。
重新标记坐标后，可假设
$h_2:R\to R[[t_{r_1+1},\ldots,t_{r_1+r_2}]]$，
且
$k_2:R[[t_{r_1+1},\ldots,t_{r_1+r_2}]]\to S$。
选择一个适当的同构
$$
R[[t_1, \ldots, t_{r_1 + r_2}]]
\longrightarrow
R[[t_{r_1 + 1}, \ldots, t_{r_1 + r_2}]]
\widehat{\otimes}_{h_2, R, h_1}
R[[t_1, \ldots, t_{r_1}]]
$$
后，可把这些映射置入可交换图表
$$
\xymatrix{
R \ar[r]_{h_1} \ar[d]^{h_2} &
R[[t_1, \ldots, t_{r_1}]] \ar[d]^{h_2'} \\
R[[t_{r_1 + 1}, \ldots, t_{r_1 + r_2}]] \ar[r]^{h_1'} &
R[[t_1, \ldots, t_{r_1 + r_2}]]
}
$$
中，其中当 $1\leq i\leq r_1$ 时 $h_2'(t_i)=t_i$，
当 $r_1+1\leq i\leq r_2$ 时 $h_1'(t_i)=t_i$。% P10-E0510
略去若干细节。
由于该图表是范畴 $\widehat{\mathcal{C}}_\Lambda$ 中的推出
（见引理 \ref{formal-defos-lemma-CLambdahat-pushouts} 的证明），
并且 $k_1\circ h_1=f_2=k_2\circ h_2$，
可知存在映射
$$
k : R[[t_1, \ldots, t_{r_1 + r_2}]] \to S
$$
使 $k_1=k\circ h_2'$ 且 $k_2=k\circ h_1'$。
记 $h=h_1'\circ h_2=h_2'\circ h_1$。于是
\begin{enumerate}
\item $k(h_1'(a))=k_2(a)=f_3(a)$；且
\item $k(h_2'(a))=k_1(a)=f_1(a)$。
\end{enumerate}
由证明第一段中的断言，这说明 $f_1$ 与 $f_3$ 形式同伦。
\end{proof}

\begin{lemma}
\label{formal-defos-lemma-composition-homotopic}
在范畴 $\widehat{\mathcal{C}}_\Lambda$ 中，若
$f_1,f_2:R\to S$ 形式同伦，且 $g:S\to S'$ 是态射，
则 $g\circ f_1$ 与 $g\circ f_2$ 形式同伦。
\end{lemma}

\begin{proof}
事实上，若 $(r,h,k)$ 是 $f_1$ 与 $f_2$ 之间的形式同伦，
则 $(r,h,g\circ k)$ 是 $g\circ f_1$ 与
$g\circ f_2$ 之间的形式同伦。
\end{proof}

\begin{lemma}
\label{formal-defos-lemma-versal-unique-up-to-homotopy}
设 $\mathcal{F}$ 是 $\mathcal{C}_\Lambda$ 上的形变范畴，
并且 $\dim_kT\mathcal{F}<\infty$、
$\dim_k\text{Inf}(\mathcal{F})<\infty$。
设 $\xi$ 是位于 $R$ 上方的万有形式对象。
设 $\eta$ 是位于 $S$ 上方的形式对象。
则任意两个满足
$$
f, g : R \to S
$$
且 $f_*\xi\cong\eta\cong g_*\xi$ 的映射都形式同伦。
\end{lemma}

\begin{proof}
根据定理 \ref{formal-defos-theorem-presentation-deformation-groupoid}
及其证明，$\mathcal{F}$ 有一个由函子中的光滑 pro-可表群胚
$$
(\underline{R}, \underline{R_1}, s, t, c, e, i)|_{\mathcal{C}_\Lambda}
$$
给出的呈现；该群胚位于 $\mathcal{C}_\lambda$ 上，
并且 $\mathcal{F}$。% P10-E0511 P10-E0512
于是映射 $s:R\to R_1$ 和 $t:R\to R_1$
是形式光滑环映射，且 $e:R_1\to R$ 是一个截面。
特别地，可对某个 $r\geq0$ 选择同构
$R_1=R[[t_1,\ldots,t_r]]$，
使 $s$ 是嵌入 $R\subset R[[t_1,\ldots,t_r]]$，
而 $t$ 对应于映射
$h:R\to R[[t_1,\ldots,t_r]]$，且对所有 $a\in R$，
$h(a)\bmod(t_1,\ldots,t_r)=a$。
同构 $\alpha:f_*\xi\to g_*\xi$ 的存在，恰好意味着
存在映射 $k:R_1\to S$，使
$f=k\circ s$ 且 $g=k\circ t$。
这恰好说明 $(r,h,k)$ 是 $f$ 与 $g$ 之间的形式同伦。
\end{proof}

\begin{lemma}
\label{formal-defos-lemma-homotopic-minimal-prime}
在范畴 $\widehat{\mathcal{C}}_\Lambda$ 中，若
$f_1,f_2:R\to S$ 形式同伦，且
$\mathfrak p\subset R$ 是极小素理想，则作为理想有
$f_1(\mathfrak p)S=f_2(\mathfrak p)S$。
\end{lemma}

\begin{proof}
设 $(r,h,k)$ 是 $f_1$ 与 $f_2$ 之间的形式同伦。
我们断言
$\mathfrak pR[[t_1,\ldots,t_r]]
=h(\mathfrak p)R[[t_1,\ldots,t_r]]$。
再与 $k$ 复合，该断言即蕴含本引理。
为证明断言，注意映射
$\mathfrak p\mapsto\mathfrak pR[[t_1,\ldots,t_r]]$
是 $R$ 的极小素理想与
$R[[t_1,\ldots,t_r]]$ 的极小素理想之间的双射。
最后，因为 $h$ 平坦，
$h(\mathfrak p)R[[t_1,\ldots,t_r]]$ 是极小素理想，
所以根据刚才所述，它形如
$\mathfrak qR[[t_1,\ldots,t_r]]$，
其中 $\mathfrak q\subset R$ 是某个极小素理想。
但根据形式同伦的定义，
$h\bmod(t_1,\ldots,t_r)=\text{id}_R$，
故 $\mathfrak q=\mathfrak p$，这正是所求。
\end{proof}

\section{更换剩余域}
\label{formal-defos-section-change-of-field}

\noindent
本节简要讨论用有限扩张替换剩余域 $k$ 后会发生什么。
设 $\Lambda$ 是诺特环，并设 $\Lambda\to k$ 是有限环映射，
其中 $k$ 是域。在本章中，我们一直用
$\mathcal{C}_\Lambda$ 表示剩余域与 $k$ 等同的
阿廷局部 $\Lambda$-代数范畴；见定义
\ref{formal-defos-definition-CLambda}。不过，本节将讨论更换 $k$ 后
会发生什么，% P10-E0513
所以改用记号 $\mathcal{C}_{\Lambda,k}$ 来指明对 $k$ 的依赖。

\begin{situation}
\label{formal-defos-situation-change-of-fields}
设 $\Lambda$ 是诺特环，并设
$\Lambda\to k\to l$ 是有限环映射，其中 $k$、$l$ 都是域。% P10-E0514
因此 $l/k$ 是有限域扩张。% P10-E0515
$\mathcal{C}_{\Lambda,l}$ 的典型对象记为 $B$，
$\mathcal{C}_{\Lambda,k}$ 的典型对象记为 $A$。定义
\begin{equation}
\label{formal-defos-equation-comparison}
\mathcal{C}_{\Lambda, l} \longrightarrow \mathcal{C}_{\Lambda, k},
\quad
B \longmapsto B \times_l k
\end{equation}
给定群胚余纤维范畴
$p:\mathcal{F}\to\mathcal{C}_{\Lambda,k}$，
令 $\mathcal{F}_{l/k}(B)=\mathcal{F}(B\times_l k)$，
便得到相应的群胚余纤维范畴
$$
p_{l/k} : \mathcal{F}_{l/k} \longrightarrow \mathcal{C}_{\Lambda, l}.
$$
\end{situation}

\noindent
函子（\ref{formal-defos-equation-comparison}）确有意义：
由于 $B\times_l k\subset B$，有
\begin{align*}
[k : k']\ \text{length}_{B \times_l k}(B \times_l k) & =
\text{length}_\Lambda(B \times_l k) \\
& \leq \text{length}_\Lambda(B) \\
& = [l : k']\ \text{length}_B(B) < \infty
\end{align*}
（见引理 \ref{formal-defos-lemma-length}），所以 $B\times_l k$ 是阿廷环
（见《代数》引理
\ref{algebra-lemma-artinian-finite-length}）。
因此 $B\times_l k$ 是剩余域为 $k$ 的阿廷局部环。
注意，（\ref{formal-defos-equation-comparison}）与纤维积可交换：
$$
(B_1 \times_B B_2) \times_l k =
(B_1 \times_l k) \times_{(B \times_l k)} (B_2 \times_l k),
$$
并且把满环映射变为满环映射。
我们使用较为“昂贵”的记号 $\mathcal{F}_{l/k}$，
以免与说明 \ref{formal-defos-remark-localize-cofibered-groupoid}
中的构造混淆。下面给出一些初等观察。

\begin{lemma}
\label{formal-defos-lemma-elementary-properties-change-of-field}
采用情形 \ref{formal-defos-situation-change-of-fields} 中的记号和假设。
\begin{enumerate}
\item 有
$\overline{\mathcal{F}_{l/k}}=(\overline{\mathcal{F}})_{l/k}$。
\item 若 $\mathcal{F}$ 是预形变范畴，则
$\mathcal{F}_{l/k}$ 是预形变范畴。
\item 若 $\mathcal{F}$ 满足 (S1)，则
$\mathcal{F}_{l/k}$ 满足 (S1)。
\item 若 $\mathcal{F}$ 满足 (S2)，则
$\mathcal{F}_{l/k}$ 满足 (S2)。
\item 若 $\mathcal{F}$ 满足 (RS)，则
$\mathcal{F}_{l/k}$ 满足 (RS)。
\end{enumerate}
\end{lemma}

\begin{proof}
(1) 直接由定义得出。

\medskip\noindent
由于 $\mathcal{F}_{l/k}(l)=\mathcal{F}(k)$，
(2) 由定义得出；见定义 \ref{formal-defos-definition-predeformation-category}。

\medskip\noindent
因为函子（\ref{formal-defos-equation-comparison}）与纤维积可交换，
且把满射变为满射，(3) 得出；见定义
\ref{formal-defos-definition-S1-S2}。

\medskip\noindent
证明 (4)。考虑 $\mathcal{C}_{\Lambda,l}$ 中如定义
\ref{formal-defos-definition-S1-S2} 所述的图表
$$
\xymatrix{
          & l[\epsilon] \ar[d] \\
B  \ar[r] & l
}.
$$
应用函子（\ref{formal-defos-equation-comparison}），得到
$$
\xymatrix{
          & k[l\epsilon] \ar[d] \\
B \times_l k  \ar[r] & k
},
$$
其中 $l\epsilon$ 表示有限维 $k$-向量空间
$l\epsilon\subset l[\epsilon]$。根据引理
\ref{formal-defos-lemma-S2-extensions}，$\mathcal{F}$ 的 (S2)
条件也对该图表成立。因此 (S2) 对
$\mathcal{F}_{l/k}$ 成立。

\medskip\noindent
(5) 由引理 \ref{formal-defos-lemma-RS-2-categorical} 的 (2)
对 (RS) 的刻画，以及
（\ref{formal-defos-equation-comparison}）与纤维积可交换这一事实得出。
\end{proof}

\noindent
下列引理尤其适用于 $\mathcal{F}$ 满足 (S2)
且为预形变范畴的情形；见引理
\ref{formal-defos-lemma-S1-S2-associated-functor}。

\begin{lemma}
\label{formal-defos-lemma-tangent-space-change-of-field}
采用情形 \ref{formal-defos-situation-change-of-fields} 中的记号和假设。
假设 $\mathcal{F}$ 是预形变范畴，且
$\overline{\mathcal{F}}$ 满足 (S2)。则切空间之间存在典范
$l$-向量空间同构
$$
T\mathcal{F} \otimes_k l \longrightarrow T\mathcal{F}_{l/k}.
$$
\end{lemma}

\begin{proof}
由引理 \ref{formal-defos-lemma-elementary-properties-change-of-field}，
可用 $\overline{\mathcal{F}}$ 替换 $\mathcal{F}$。
此外，由引理 \ref{formal-defos-lemma-tangent-space-vector-space}，
$T\mathcal{F}$、$T\mathcal{F}_{l/k}$ 分别具有典范
$k$-向量空间结构、$l$-向量空间结构。于是
$$
T\mathcal{F}_{l/k} = \mathcal{F}_{l/k}(l[\epsilon])
= \mathcal{F}(k[l\epsilon]) = T\mathcal{F} \otimes_k l,
$$
最后一个等号由引理
\ref{formal-defos-lemma-tangent-space-vector-space} 得出。
更一般地，给定有限维 $l$-向量空间 $V$，有
$$
\mathcal{F}_{l/k}(l[V]) = \mathcal{F}(k[V_k]) = T\mathcal{F} \otimes_k V_k,
$$
其中 $V_k$ 表示把 $V$ 看作 $k$-向量空间。
因此，作为取值于集合范畴的函子，函子
$V\mapsto\mathcal{F}_{l/k}(l[V])$ 和
$V\mapsto T\mathcal{F}\otimes_k V_k$ 典范等同。
由引理 \ref{formal-defos-lemma-linear-functor}，可见至多有一种方法
使任一函子成为 $l$-线性函子。因此这些同构与
$l$-向量空间结构相容，结论得证。
\end{proof}

\begin{lemma}
\label{formal-defos-lemma-inf-aut-change-of-field}
采用情形 \ref{formal-defos-situation-change-of-fields} 中的记号和假设。
假设 $\mathcal{F}$ 是形变范畴。
则无穷小自同构空间之间存在典范 $l$-向量空间同构
$$
\text{Inf}(\mathcal{F}) \otimes_k l
\longrightarrow
\text{Inf}(\mathcal{F}_{l/k}).
$$
\end{lemma}

\begin{proof}
设 $x_0\in\Ob(\mathcal{F}(k))$，并用 $x_{l,0}$
表示 $\mathcal{F}_{l/k}$ 在 $l$ 上方的相应对象。% P10-E0516
回顾
$\text{Inf}(\mathcal{F})=\text{Inf}_{x_0}(\mathcal{F})$ 且
$\text{Inf}(\mathcal{F}_{l/k})
=\text{Inf}_{x_{l,0}}(\mathcal{F}_{l/k})$；
见说明 \ref{formal-defos-remark-trivial-aut-point}。
回顾 $\text{Inf}_{x_0}(\mathcal{F})$ 的向量空间结构来自
将其等同为函子 $\mathit{Aut}(x_0)$ 的切空间；
该函子定义在剩余域为 $k$ 的阿廷局部 $k$-代数范畴
$\mathcal{C}_{k,k}$ 上。类似地，
$\text{Inf}_{x_{l,0}}(\mathcal{F}_{l/k})$
是 $\mathit{Aut}(x_{l,0})$ 的切空间；
后者定义在剩余域为 $l$ 的阿廷局部 $l$-代数范畴
$\mathcal{C}_{l,l}$ 上。展开定义可见，
$\mathit{Aut}(x_{l,0})$ 是
$\mathit{Aut}(x_0)_{l/k}$（它定义在
$\mathcal{C}_{k,l}$ 上）到 $\mathcal{C}_{l,l}$ 的限制。
把 $\mathit{Aut}(x_0)_{l/k}$ 看作
$\mathcal{C}_{k,l}$ 或 $\mathcal{C}_{l,l}$ 上的函子，
其切空间没有差别。因此本引理由引理
\ref{formal-defos-lemma-tangent-space-change-of-field} 以及下述事实得出：
由引理 \ref{formal-defos-lemma-Aut-functor-RS}，
$\mathit{Aut}(x_0)$ 满足 (RS)
（因而由引理 \ref{formal-defos-lemma-RS-implies-S1-S2} 满足 (S2)）。
\end{proof}

\begin{lemma}
\label{formal-defos-lemma-change-of-fields-smooth}
采用情形 \ref{formal-defos-situation-change-of-fields} 中的记号和假设。
若 $\mathcal{F}\to\mathcal{G}$ 是
$\mathcal{C}_{\Lambda,k}$ 上群胚余纤维范畴的光滑态射，
则 $\mathcal{F}_{l/k}\to\mathcal{G}_{l/k}$ 是
$\mathcal{C}_{\Lambda,l}$ 上群胚余纤维范畴的光滑态射。
\end{lemma}

\begin{proof}
这直接由定义以及
（\ref{formal-defos-equation-comparison}）保持满射这一事实得出。
\end{proof}

\noindent
关于 $\mathcal{F}$ 与 $\mathcal{F}_{l/k}$ 的关系，
还有更多可以说明的内容（特别是关于万有形变之间的关系）；
我们会在需要时把这些内容添加到这里。

\begin{lemma}
\label{formal-defos-lemma-change-of-field-versal-ring}
采用情形 \ref{formal-defos-situation-change-of-fields} 中的记号和假设。
设 $\xi$ 是 $\mathcal{F}$ 的万有形式对象，位于
$R\in\Ob(\widehat{\mathcal{C}}_{\Lambda,k})$ 上方。则存在：
\begin{enumerate}
\item 一个 $S\in\Ob(\widehat{\mathcal{C}}_{\Lambda,l})$
以及局部 $\Lambda$-代数同态 $R\to S$；
该同态在 $\mathfrak m_S$-进拓扑下形式光滑，
并在剩余域上诱导给定的域扩张 $l/k$；以及
\item $\mathcal{F}_{l/k}$ 的一个位于 $S$ 上方的万有形式对象。
\end{enumerate}
\end{lemma}

\begin{proof}
构造 $S$。选择 $R$-代数的满射
$R[x_1,\ldots,x_n]\to l$。其核为极大理想 $\mathfrak m$。
令 $S$ 为诺特环 $R[x_1,\ldots,x_n]$ 的
$\mathfrak m$-进完备化。于是由《代数》引理
\ref{algebra-lemma-completion-Noetherian-Noetherian}，
$S$ 属于 $\widehat{\mathcal{C}}_{\Lambda,l}$。
由《代数进阶》引理
\ref{more-algebra-lemma-formally-smooth} 和
\ref{more-algebra-lemma-formally-smooth-completion}，
以及 $R\to R[x_1,\ldots,x_n]$ 形式光滑这一事实，
映射 $R\to S$ 在 $\mathfrak m_S$-进拓扑下形式光滑。
（比较引理 \ref{formal-defos-lemma-exists-smooth} 的证明。）% P10-E0517

\medskip\noindent
由于 $\xi$ 万有，变换
$\underline{\xi}:\underline{R}|_{\mathcal{C}_{\Lambda,k}}
\to\mathcal{F}$ 光滑。由引理
\ref{formal-defos-lemma-change-of-fields-smooth}，诱导映射
$$
(\underline{R}|_{\mathcal{C}_{\Lambda, k}})_{l/k}
\longrightarrow
\mathcal{F}_{l/k}
$$
光滑。因此只需构造光滑态射
$\underline{S}|_{\mathcal{C}_{\Lambda,l}}\to
(\underline{R}|_{\mathcal{C}_{\Lambda,k}})_{l/k}$。
给出这样的映射，意即对
$\mathcal{C}_{\Lambda,l}$ 的每个对象 $B$，
给出关于 $B$ 函子性的集合映射
$$
\Mor_{\widehat{\mathcal{C}}_{\Lambda, l}}(S, B)
\longrightarrow
\Mor_{\widehat{\mathcal{C}}_{\Lambda, k}}(R, B \times_l k).
$$
给定左侧元素 $\varphi:S\to B$，
把它映到复合 $R\to S\to B$；
该复合的像包含在子 $\Lambda$-代数 $B\times_l k$ 中。
该映射的光滑性意味着：给定满射 $B'\to B$ 和可交换图表
$$
\xymatrix{
S \ar[r] & B \ar@{=}[r] & B \\
R \ar[u] \ar[r] & B' \times_l k \ar[u] \ar[r] & B' \ar[u]
},
$$
必须找到环映射 $S\to B'$，使其嵌入外侧矩形。
由于我们选择的 $R\to S$ 在 $\mathfrak m_S$-进拓扑下
形式光滑，这一映射的存在性由《代数进阶》引理
\ref{more-algebra-lemma-lift-continuous} 保证。
\end{proof}
